\documentclass{article}
\usepackage[margin=0.8in]{geometry}
\usepackage{amsmath,amssymb,amsthm,mathtools}
\usepackage{mathrsfs,bm,bbm}
\usepackage{graphicx,booktabs,array,multirow,tabularx,longtable}
\usepackage{enumitem}
\usepackage{xcolor}
\usepackage{algorithm,algpseudocode}
\usepackage{subcaption,tikz}
\usepackage{siunitx,cancel,accents}
\usepackage{microtype}
\usepackage[numbers,sort&compress]{natbib}
\usepackage[colorlinks=true,linkcolor=blue,citecolor=blue,urlcolor=blue]{hyperref}
\usepackage{cleveref}
\setlength{\emergencystretch}{2em}
\newtheorem{assumption}{Assumption}
\newtheorem{definition}{Definition}
\newtheorem{theorem}{Theorem}
\newtheorem{lemma}{Lemma}
\newtheorem{proposition}{Proposition}
\newtheorem{openendedquestion}{Open-ended Question}
\newtheorem{conjecture}{Conjecture}
\newcommand{\coreref}[1]{\ref{#1}}

\begin{document}

\sloppy

\section*{eid\_\allowbreak{}extra\_\allowbreak{}parent\_\allowbreak{}degree}

\noindent\textit{Specialization:} v1\par

\paragraph{TL;DR.} Adding parents to a supplied rooted tree creates finite ambiguity measured by the single-coordinate complex degree \(D_k(n)\), not by whole-fiber cardinality, real root count, or local rank. For \(n\geq4\), the proved envelope is \(2^{\min\{k+1,\lfloor(n-2)/2\rfloor\}}\leq D_k(n)\leq\vartheta_{n-1,k+1}\leq\min\{2^{k+1},\tau_{n-1}\}\). The upper bound selects finitely many nonsingular conjugate zeros with distinct queried coordinates and applies a paper-owned same-multidegree perturbation bridge; only after its block-projective perturbation has a finite common-zero set is the cited multihomogeneous B\'ezout coefficient bound invoked. The exact maximality thresholds are \(\vartheta_{N,g}=2^g\) if and only if \(N\geq2g+1\), and \(D_k(n)=2^{k+1}\) if and only if \(n\geq2k+4\). Exact positive-excess values below the latter threshold remain open from six vertices onward. The zero-excess dispatcher retains the published \(1/2/\infty\) classification; positive excess unconditionally returns an algebraic-versus-transcendental tag and a size-binomial degree cap, while exact polynomials, exceptional sets, real lifting, and regular inference remain chart-conditional.

\section{Project justification}

\paragraph{Gap.} The published tree algorithms completely solve exactly \(\mathcal G_{n,0}\), not \(\mathcal G_{n,k}\) for positive \(k\). General mixed-graph work studies rational recovery, whole-parameter fibers, or unparameterized elimination, and does not give the maximum finite degree of one original directed coefficient as a function of both vertex size and excess-parent budget. General geometric-resolution algorithms also do not select the embedded SEM source prime or provide bounded descent over the implicitly presented covariance field.

\paragraph{Niche.} Distance from a supplied arborescence creates a parameterized algebraic regime in which the orientation system yields a projection-aware ceiling and explicit constructions yield sharp size-sensitive lower bounds. The zero-excess class inherits a complete published algorithm, while positive excess separates the unconditional algebraicity tag from chart-certified exact finite output.

\paragraph{Fill.} The paper proves the coupled envelope \(2^{\min\{k+1,\lfloor(n-2)/2\rfloor\}}\leq D_k(n)\leq\vartheta_{n-1,k+1}\leq\min\{2^{k+1},\tau_{n-1}\}\) for \(n\geq4\), together with the exact maximality threshold \(D_k(n)=2^{k+1}\) if and only if \(n\geq2k+4\). The upper proof does not assume a zero-dimensional whole fiber: it constructs a finite selected set of nonsingular conjugate zeros, preserves their distinct queried coordinates under specialization, and invokes the paper-owned finite-nonsingular perturbation lemma to pass to a finite block-projective system before using the cited finite-system multihomogeneous B\'ezout theorem. The paper also gives the exact five-vertex orientation recurrence and the resulting six-vertex caps, while leaving the remaining below-threshold frontier open. Operationally, it supplies the complete published tree output and an unconditional positive-excess rank tag, plus a conditional geometric-resolution theorem with explicit straight-line-program and intermediate-degree dependence. Quantifier elimination is used only for existence of an equivalent quantifier-free real-lifting condition; exact algebraic sampling and the paper's own reductions supply the specialization certificates.

\section{Related work}

Van der Zander, Wienöbst, Bläser, and Liśkiewicz introduced TreeID for directed-tree linear SEMs. Gupta and Bläser proved the exact \(1/2/\infty\) queried-coordinate classification on that zero-excess class, and Briefs and Bläser accelerated the same tree workflow in fasttreeid; none of these results covers positive excess. Foygel, Draisma, and Drton supply older higher-degree small SEM examples, which are not novelty claims here. Hollering, Misra, and Sturma study bounded-degree rational formulas rather than algebraic coordinate-fiber degree. The multihomogeneous B\'ezout citation used here bounds a finite common-zero set only; the passage from finitely many selected nonsingular zeros inside a possibly positive-dimensional residual system to that finite setting is proved by the paper-owned perturbation lemma. Giusti, Lecerf, and Salvy, and Lecerf, provide the geometric-resolution and lifting-fiber substrates with geometric-degree dependence, but do not select the embedded SEM source prime or construct its valid chart. Basu, Pollack, and Roy are cited only for existence of quantifier elimination over real closed fields, and the semialgebraic stratification citation is used only for a finite disjoint union into smooth manifolds; the SEM-specific regular cell and real-extension conclusions are derived in this paper.

\section{Setup}

Target (estimand / structural parameter): \(\lambda_e\).
Primary functional / estimator / diagnostic: \(D_k(n)=\max_{(G,T,e)\in\mathcal G_{n,k},\ d_e(G)<\infty}[K_G(\lambda_e):K_G]\). For \(n\geq4\), the proved two-sided envelope is \(2^{\min\{k+1,\lfloor(n-2)/2\rfloor\}}\leq D_k(n)\leq\vartheta_{n-1,k+1}\leq\min\{2^{k+1},\tau_{n-1}\}\). The exact maximality characterizations are \(\vartheta_{N,g}=2^g\) if and only if \(N\geq2g+1\), and, for every admissible \((n,k)\), \(D_k(n)=2^{k+1}\) if and only if \(n\geq2k+4\); in particular, \(D_k(n)<2^{k+1}\) below that size threshold. The exact five-vertex orientation recurrence is \((\vartheta_{5,g})_{g=0}^{6}=(1,2,4,6,8,12,24)\), giving \(D_2(6)\leq6\), \(D_3(6)\leq8\), \(D_4(6)\leq12\), and \(D_k(6)\leq24\) for admissible \(k\geq5\); the established equalities are \(D_0(6)=2\) and \(D_1(6)=4\). The exact table is \(D_0(2)=1\), \(D_k(3)=1\), \(D_k(4)=2\), \(D_0(5)=2\), \(D_1(5)=3\), and \(D_k(5)=4\) for \(k\geq2\). At zero excess the dispatcher returns the exact published \(1/2/\infty\) tree classification and finite \(\operatorname{FASTP}\) expressions; at positive excess its unconditional algebraic branch retains the explicit computable report \(d_e(G)\leq2^{\min\{|D|-n+2,\binom{n-2}{2}\}}\) and does not claim to compute \(\vartheta_{n-1,k+1}\). A rational-circuit representation of \(p_e(t)\) is delivered only under the valid-chart and successful-output conditions of \(\mathrm{thm:conditional\mbox{-}geometric\mbox{-}resolution}\)..
\begin{itemize}
  \item \texttt{V} --- \texttt{finite vertex set}
  \par\smallskip\noindent\emph{Definition.} \texttt{The observed-\allowbreak{}variable vertices.}
  \item \texttt{n} --- \texttt{positive integer}; \(\mathbb N\)
  \par\smallskip\noindent\emph{Definition.} \(n=|V|\).
  \item \texttt{D} --- \texttt{directed edge set}; \(V\times V\)
  \par\smallskip\noindent\emph{Definition.} \texttt{The directed support.}
  \item \texttt{B} --- \texttt{bidirected edge set}; \(\{\{i,j\}:i,j\in V,\ i\ne j\}\)
  \par\smallskip\noindent\emph{Definition.} \texttt{The pairs whose disturbances may be correlated.}
  \item \texttt{G} --- \texttt{mixed graph}
  \par\smallskip\noindent\emph{Definition.} \(G=(V,D,B)\).
  \item \texttt{T} --- \texttt{supplied directed edge subset}; \(2^D\)
  \par\smallskip\noindent\emph{Definition.} \texttt{The candidate directed backbone supplied with the graph.}
  \item \texttt{k} --- \texttt{nonnegative integer}; \(\mathbb N_0\)
  \par\smallskip\noindent\emph{Definition.} The permitted number of directed edges outside \(T\).
  \item \texttt{e} --- \texttt{queried directed edge}; \(D\)
  \par\smallskip\noindent\emph{Definition.} A fixed edge \(e=(p,q)\in D\).
  \item \texttt{Lambda} --- \texttt{structural coefficient matrix}; \(\mathbb C^{n\times n}\)
  \par\smallskip\noindent\emph{Definition.} The matrix \(\Lambda\) supported on \(D\), in the convention \(X=\Lambda^{\mathsf T}X+\varepsilon\).
  \item \texttt{lambda\_\allowbreak{}D} --- \texttt{family of directed source coordinates}
  \par\smallskip\noindent\emph{Definition.} \(\lambda_D=(\lambda_{ij}:i\to j\in D)\), the nonzero-support coordinates of \(\Lambda\).
  \item \texttt{Omega} --- \texttt{disturbance covariance matrix}; \(\mathbb C^{n\times n}_{\mathrm{sym}}\)
  \par\smallskip\noindent\emph{Definition.} \texttt{The symmetric disturbance matrix before its graph-\allowbreak{}support restriction is imposed.}
  \item \texttt{omega\_\allowbreak{}B} --- \texttt{family of bidirected source coordinates}
  \par\smallskip\noindent\emph{Definition.} \(\omega_B=(\omega_{ij}:i\leftrightarrow j\in B)\).
  \item \texttt{omega\_\allowbreak{}diag} --- \texttt{family of diagonal source coordinates}
  \par\smallskip\noindent\emph{Definition.} \(\omega_{\mathrm{diag}}=(\omega_{ii}:i\in V)\).
  \item \texttt{I\_\allowbreak{}n} --- \texttt{identity matrix}; \(\mathbb Q^{n\times n}\)
  \par\smallskip\noindent\emph{Definition.} The identity matrix of order \(n\).
  \item \texttt{Sigma} --- \texttt{observed covariance matrix}; \(\mathbb C^{n\times n}_{\mathrm{sym}}\)
  \par\smallskip\noindent\emph{Definition.} \(\Sigma=(I_n-\Lambda)^{-\mathsf T}\Omega(I_n-\Lambda)^{-1}\).
  \item \texttt{sigma\_\allowbreak{}ij} --- \texttt{family of observed covariance coordinates}
  \par\smallskip\noindent\emph{Definition.} \(\sigma_{ij}=\Sigma_{ij}\) for \(i,j\in V\).
  \item \texttt{R\_\allowbreak{}sigma} --- \texttt{covariance-\allowbreak{}coordinate ring}; polynomial ring over \(\mathbb C\)
  \par\smallskip\noindent\emph{Definition.} \(R_{\sigma}=\mathbb C[\sigma_{ij}:i\le j,\ i,j\in V]\).
  \item \texttt{R\_\allowbreak{}src} --- \texttt{acyclic SEM source ring}; polynomial ring over \(\mathbb C\)
  \par\smallskip\noindent\emph{Definition.} \(R_{\mathrm{src}}=\mathbb C[\lambda_D,\omega_B,\omega_{\mathrm{diag}}]\).
  \item \texttt{phi\_\allowbreak{}G} --- \texttt{source-\allowbreak{}incidence homomorphism on covariance coordinates}; \(\phi_G:R_{\sigma}\to R_{\mathrm{src}}\)
  \par\smallskip\noindent\emph{Definition.} The \(\mathbb C\)-algebra homomorphism defined by \(\phi_G(\sigma_{ij})=[(I_n-\Lambda)^{-\mathsf T}\Omega(I_n-\Lambda)^{-1}]_{ij}\), with every missing off-diagonal disturbance entry fixed at zero; acyclicity makes every image entry polynomial in \(R_{\mathrm{src}}\).
  \item \texttt{P\_\allowbreak{}G} --- \texttt{covariance-\allowbreak{}model prime ideal}; prime ideal of \(R_{\sigma}\)
  \par\smallskip\noindent\emph{Definition.} \(P_G=\ker(\phi_G)\).
  \item \texttt{R\_\allowbreak{}G} --- \texttt{covariance-\allowbreak{}model coordinate ring}; \texttt{integral domain}
  \par\smallskip\noindent\emph{Definition.} \(R_G=R_{\sigma}/P_G\cong\operatorname{im}(\phi_G)\).
  \item \texttt{S\_\allowbreak{}G} --- \texttt{generic-\allowbreak{}covariance multiplicative set}; multiplicative subset of \(R_{\sigma}\)
  \par\smallskip\noindent\emph{Definition.} \(S_G=R_{\sigma}\setminus P_G=\{f\in R_{\sigma}:\phi_G(f)\ne0\}\).
  \item \texttt{R\_\allowbreak{}inc} --- \texttt{covariance-\allowbreak{}coefficient incidence polynomial ring}; polynomial ring over \(\mathbb C\)
  \par\smallskip\noindent\emph{Definition.} \(R_{\mathrm{inc}}=R_{\sigma}[\lambda_D]\).
  \item \texttt{psi\_\allowbreak{}G} --- \texttt{acyclic SEM incidence homomorphism}; \(\psi_G:R_{\mathrm{inc}}\to R_{\mathrm{src}}\)
  \par\smallskip\noindent\emph{Definition.} The \(\mathbb C\)-algebra homomorphism extending \(\phi_G\), fixing each coordinate in \(\lambda_D\), and sending each \(\sigma_{ij}\) to its acyclic covariance parametrization.
  \item \texttt{L} --- \texttt{complexified source function field}; field extension of \(\mathbb C\)
  \par\smallskip\noindent\emph{Definition.} \(L=\mathbb C(\lambda_D,\omega_B,\omega_{\mathrm{diag}})\).
  \item \texttt{K\_\allowbreak{}G} --- \texttt{observed covariance subfield}; subfield of \(L\)
  \par\smallskip\noindent\emph{Definition.} \(K_G=\operatorname{Frac}(R_G)=\operatorname{Frac}(\operatorname{im}\phi_G)=\mathbb C(\phi_G(\sigma_{ij}):i,j\in V)\subseteq L\).
  \item \texttt{lambda\_\allowbreak{}e} --- \texttt{queried structural coefficient}; \(L\)
  \par\smallskip\noindent\emph{Definition.} \(\lambda_e=\Lambda_{pq}\) for \(e=(p,q)\).
  \item \texttt{d\_\allowbreak{}e} --- \texttt{single-\allowbreak{}coordinate complex generic degree}; \(\mathbb N\cup\{\infty\}\)
  \par\smallskip\noindent\emph{Definition.} \(d_e(G)=[K_G(\lambda_e):K_G]\) when \(\lambda_e\) is algebraic over \(K_G\), and \(d_e(G)=\infty\) otherwise.
  \item \texttt{G\_\allowbreak{}nk} --- \texttt{tree-\allowbreak{}plus-\allowbreak{}excess-\allowbreak{}arrow graph-\allowbreak{}query class}
  \par\smallskip\noindent\emph{Definition.} The class \(\mathcal G_{n,k}\) specified in \(\mathrm{def:graph\mbox{-}class}\).
  \item \texttt{D\_\allowbreak{}kn} --- \texttt{sharp coordinate-\allowbreak{}degree frontier}; \(\mathbb N\)
  \par\smallskip\noindent\emph{Definition.} \(D_k(n)=\max\{d_e(G):(G,T,e)\in\mathcal G_{n,k},\ d_e(G)<\infty\}\).
  \item \texttt{F\_\allowbreak{}ij} --- \texttt{missing-\allowbreak{}bidirected residual polynomial}; \(F_{ij}(\Lambda,\Sigma)\in K_G[\lambda_D]\)
  \par\smallskip\noindent\emph{Definition.} For \(i\leftrightarrow j\notin B\), \(F_{ij}=\sigma_{ij}-\sum_{p\in\operatorname{pa}_D(i)}\lambda_{pi}\sigma_{pj}-\sum_{q\in\operatorname{pa}_D(j)}\lambda_{qj}\sigma_{iq}+\sum_{p,q}\lambda_{pi}\lambda_{qj}\sigma_{pq}\).
  \item \texttt{I\_\allowbreak{}G} --- \texttt{source-\allowbreak{}incidence prime ideal}; prime ideal of \(R_{\mathrm{inc}}\)
  \par\smallskip\noindent\emph{Definition.} \(I_G=\ker(\psi_G)=(F_{ij}:i\ne j,\ i\leftrightarrow j\notin B)\).
  \item \texttt{A\_\allowbreak{}K} --- \texttt{localized generic-\allowbreak{}fiber coordinate algebra}
  \par\smallskip\noindent\emph{Definition.} \(A_K=S_G^{-1}(R_{\mathrm{inc}}/I_G)\cong K_G\otimes_{R_G}(R_{\mathrm{inc}}/I_G)\), the integral generic-fiber algebra obtained by inverting exactly \(S_G\) and no source-coordinate expression.
  \item \texttt{U} --- \texttt{coordinate transcendence basis}; \(2^{\lambda_D}\)
  \par\smallskip\noindent\emph{Definition.} A maximal subset \(U\subseteq\lambda_D\setminus\{\lambda_e\}\) algebraically independent over \(K_G\).
  \item \texttt{FrontierHandle} --- \texttt{orientation-\allowbreak{}realization program}
  \par\smallskip\noindent\emph{Definition.} The proof handle specified in \(\mathrm{def:frontier\mbox{-}handle}\).
  \item \texttt{CensusHandle} --- \texttt{bounded characteristic-\allowbreak{}zero elimination audit}
  \par\smallskip\noindent\emph{Definition.} The exact small-size falsification handle specified in \(\mathrm{def:census\mbox{-}handle}\).
  \item \texttt{G\_\allowbreak{}star} --- \texttt{one-\allowbreak{}sink attaining graph family}
  \par\smallskip\noindent\emph{Definition.} The family \(G^{\star}_{k,n}\) specified in \(\mathrm{def:attaining\mbox{-}family}\).
  \item \texttt{e\_\allowbreak{}star} --- \texttt{separating edge query}; directed edge of \(G^{\star}_{k,n}\)
  \par\smallskip\noindent\emph{Definition.} The first sink coefficient in the family \(G^{\star}_{k,n}\).
  \item \texttt{Core\_\allowbreak{}Ge} --- \texttt{localized quadratic coordinate core}
  \par\smallskip\noindent\emph{Definition.} The core produced by \(\mathrm{def:quadratic\mbox{-}core}\).
  \item \texttt{delta} --- \texttt{symbolic failure tolerance}; \((0,1)\)
  \par\smallskip\noindent\emph{Definition.} \texttt{The requested upper bound on randomized symbolic error.}
  \item \texttt{C} --- \texttt{universal positive constant}; \((0,\infty)\)
  \par\smallskip\noindent\emph{Definition.} A constant independent of \(G,T,e,n,k,\delta\).
  \item \texttt{ArithmeticModel} --- \texttt{exact arithmetic and identity-\allowbreak{}testing model}
  \par\smallskip\noindent\emph{Definition.} Unit-cost arithmetic over \(\mathbb Q\) and finite fields, rational arithmetic circuits in \(\sigma_{ij}\), Schwartz--Zippel testing at independently sampled field points, and exact rational reconstruction with bit complexity charged separately.
  \item \texttt{N\_\allowbreak{}Ge} --- \texttt{transcendence certificate}
  \par\smallskip\noindent\emph{Definition.} \(\mathcal N_{G,e}\) is a source-parametric dominance certificate showing that the \(\lambda_e\)-projection of a positive-dimensional irreducible component of the generic covariance fiber is Zariski dense in \(\mathbb A^1_{K_G}\), and hence that \(\lambda_e\) is transcendental over \(K_G\).
  \item \texttt{Z\_\allowbreak{}e} --- \texttt{complex queried-\allowbreak{}coordinate projection}
  \par\smallskip\noindent\emph{Definition.} \(\mathcal Z_e(G,\Sigma)=\{z\in\mathbb C:\exists\Lambda\in\mathbb C^D,\ \lambda_e=z,\ F_{ij}(\Lambda,\Sigma)=0\text{ for every }i\leftrightarrow j\notin B\}\), the finite-coordinate projection of the specialization of the prime incidence ideal \(I_G\).
  \item \texttt{R\_\allowbreak{}e} --- \texttt{retained real queried-\allowbreak{}coordinate set}
  \par\smallskip\noindent\emph{Definition.} \(\mathcal R_e(G,\Sigma)=\{z\in\mathbb R:\exists\Lambda\in\mathbb R^D,\ \lambda_e=z,\ F_{ij}(\Lambda,\Sigma)=0\text{ for every }i\leftrightarrow j\notin B\}\).
  \item \texttt{m} --- \texttt{sample size}; \(\mathbb N\)
  \par\smallskip\noindent\emph{Definition.} \texttt{The number of independent observational draws.}
  \item \texttt{X\_\allowbreak{}sample} --- \texttt{observational sample}; \((\mathbb R^n)^m\)
  \par\smallskip\noindent\emph{Definition.} \(X^{(1)},\ldots,X^{(m)}\) are the observed random vectors.
  \item \texttt{S\_\allowbreak{}m} --- \texttt{sample covariance}; \(\mathbb R^{n\times n}_{\mathrm{sym}}\)
  \par\smallskip\noindent\emph{Definition.} \(S_m=m^{-1}\sum_{r=1}^m X^{(r)}X^{(r)\mathsf T}\).
  \item \texttt{JCert\_\allowbreak{}Ge} --- \texttt{Jacobian-\allowbreak{}rank algebraicity transcript}
  \par\smallskip\noindent\emph{Definition.} \(\mathcal J_{G,e}\) records the exact source specialization, pivot positions, pivot minors, and identity-test transcript for the covariance Jacobian and the Jacobian augmented by \(d\lambda_e\); on the error-free event it certifies equality of their generic ranks.
  \item \texttt{M\_\allowbreak{}G} --- \texttt{real covariance model}
  \par\smallskip\noindent\emph{Definition.} \(\mathcal M_G=\{(I_n-\Lambda)^{-\mathsf T}\Omega(I_n-\Lambda)^{-1}:\Lambda\in\mathbb R^D,\ \Omega\succ0,\ \Omega_{ij}=0\text{ if }i\leftrightarrow j\notin B\}\).
  \item \texttt{Chart\_\allowbreak{}Ge} --- \texttt{source-\allowbreak{}prime chart package}
  \par\smallskip\noindent\emph{Definition.} \(\mathfrak C_{G,e}\) is the valid localized reduced regular-sequence and lifting package defined in \(\mathrm{def:source\mbox{-}prime\mbox{-}chart}\).
  \item \texttt{gamma} --- \texttt{branch-\allowbreak{}separation margin}; \((0,\infty)\)
  \par\smallskip\noindent\emph{Definition.} A positive uniform lower bound on distances between distinct retained roots whose existence is derived in \(\mathrm{lem:derived\mbox{-}root\mbox{-}separation}\).
  \item \texttt{p\_\allowbreak{}e} --- \texttt{queried-\allowbreak{}coordinate minimal polynomial}; \(K_G[t]\)
  \par\smallskip\noindent\emph{Definition.} When \(d_e(G)<\infty\), \(p_e(t)\) is the unique monic minimal polynomial of \(\lambda_e\) over \(K_G\). A bounded rational-circuit representation in the covariance entries is output only under the valid-chart and successful-output conditions of \(\mathrm{thm:conditional\mbox{-}geometric\mbox{-}resolution}\).
  \item \texttt{Q\_\allowbreak{}Ge} --- \texttt{conditional finite family of certificate circuits}
  \par\smallskip\noindent\emph{Definition.} Given \(d_e(G)<\infty\), a valid \(\mathfrak C_{G,e}\), and successful conditional geometric resolution, \(\mathcal Q_{G,e}\) is the returned family of covariance-localization, incidence-lifting, pole, projective-leading-form, specialized-leading-coefficient, squarefreeness, and real-lifting circuits in the entries of \(\Sigma\).
  \item \texttt{h\_\allowbreak{}Ge} --- \texttt{conditional bad-\allowbreak{}locus circuit}; \(\mathbb Q[\sigma_{ij}]\)
  \par\smallskip\noindent\emph{Definition.} On a successful conditional finite output, \(h_{G,e}=\prod_{q\in\mathcal Q_{G,e}}q\), with repeated nonzero factors removed.
  \item \texttt{Alg} --- \texttt{randomized exact zero-\allowbreak{}excess/\allowbreak{}tree and positive-\allowbreak{}excess/\allowbreak{}tag dispatcher}; \(r=|D\setminus T|\) and \(\mathsf A(G,T,e,\delta)=\begin{cases}(\operatorname{tree},d_e(G),(\rho_i)_{i=1}^{d_e(G)}),&r=0,\ d_e(G)<\infty,\\(\operatorname{tree},\infty),&r=0,\ d_e(G)=\infty,\\(\operatorname{algebraic},\mathcal J_{G,e},2^{\min\{|D|-n+2,\binom{n-2}{2}\}}),&r>0,\ d_e(G)<\infty,\\(\infty,\mathcal N_{G,e}),&r>0,\ d_e(G)=\infty.\end{cases}\)
  \par\smallskip\noindent\emph{Definition.} The dispatcher specified in \(\mathrm{def:certificate\mbox{-}algorithm}\). On the zero-excess tree class it returns the exact \(1/2/\infty\) classification and all one or two finite \(\operatorname{FASTP}\) expressions \(\rho_i\). On positive excess it returns only the randomized Jacobian-rank algebraic/transcendental tag, the size-sensitive universal degree upper bound on the algebraic branch, or the dominance certificate on the transcendental branch; it does not return the exact positive-excess degree or eliminant.
  \item \texttt{E\_\allowbreak{}Ge} --- \texttt{conditional algebraic exceptional set}
  \par\smallskip\noindent\emph{Definition.} Under a valid source-prime chart and successful conditional finite output, \(\mathcal E_{G,e}=\{\Sigma:h_{G,e}(\Sigma)=0\}\).
  \item \texttt{J} --- \texttt{number of retained real branches}; \(\mathbb N\)
  \par\smallskip\noindent\emph{Definition.} \(J=|\mathcal R_e(G,\Sigma)|\), which is positive and constant on the selected connected regular cell.
  \item \texttt{LiftCert} --- \texttt{conditional real-\allowbreak{}extension certificate}
  \par\smallskip\noindent\emph{Definition.} Under a valid source-prime chart and successful conditional finite output, a rational-univariate lifting representation, isolating interval, and exact sign data certify whether a simple real root of \(p_e\) belongs to \(\mathcal R_e(G,\Sigma)\).
  \item \texttt{C\_\allowbreak{}Ge} --- \texttt{conditional smooth connected regular real-\allowbreak{}model cell}
  \par\smallskip\noindent\emph{Definition.} Under a valid source-prime chart and successful conditional finite output, \(\mathcal C_{G,e}\) is the connected smooth coordinate cell specified in \(\mathrm{def:regular\mbox{-}cell}\), refined so that every real-root extension label is constant.
  \item \texttt{V\_\allowbreak{}Sigma} --- \texttt{Gaussian covariance limit matrix}
  \par\smallskip\noindent\emph{Definition.} \((V_\Sigma)_{ij,kl}=\sigma_{ik}\sigma_{jl}+\sigma_{il}\sigma_{jk}\).
  \item \texttt{d\_\allowbreak{}H} --- \texttt{Hausdorff distance}; \(d_H:2^{\mathbb R}\times2^{\mathbb R}\to[0,\infty]\)
  \par\smallskip\noindent\emph{Definition.} For finite \(A,B\subset\mathbb R\), \(d_H(A,B)\) is the ordinary Hausdorff distance when both sets are nonempty, \(d_H(\varnothing,\varnothing)=0\), and \(d_H(A,B)=\infty\) when exactly one of \(A,B\) is empty.
  \item \texttt{Kappa\_\allowbreak{}G} --- \texttt{conditional compact regular covariance class}
  \par\smallskip\noindent\emph{Definition.} Under a valid source-prime chart and successful conditional finite output, \(\mathcal K_G\) is a nonempty compact subset with nonempty relative interior in \(\mathcal C_{G,e}\), as specified in \(\mathrm{def:regular\mbox{-}class}\).
  \item \texttt{pi\_\allowbreak{}G} --- \texttt{local model projection}; \(\pi_G:\mathcal N(\mathcal K_G)\to\mathcal M_G\)
  \par\smallskip\noindent\emph{Definition.} A smooth nearest-point projection from a tubular neighborhood of \(\mathcal K_G\) onto its smooth covariance-model stratum.
  \item \texttt{r\_\allowbreak{}vec} --- \texttt{locally labeled retained-\allowbreak{}root map}; \(r=(r_1,\ldots,r_J):\mathcal K_G\to\mathbb R^J\)
  \par\smallskip\noindent\emph{Definition.} The increasing local labeling of \(\mathcal R_e(G,\Sigma)\) on a connected regular neighborhood.
  \item \texttt{alpha} --- \texttt{nominal error level}; \((0,1)\)
  \par\smallskip\noindent\emph{Definition.} \texttt{The target simultaneous noncoverage probability.}
  \item \texttt{Rhat\_\allowbreak{}m} --- \texttt{conditional finite-\allowbreak{}set coefficient estimator}
  \par\smallskip\noindent\emph{Definition.} Under a valid source-prime chart and successful conditional finite output, \(\widehat{\mathcal R}_{e,m}\) is the estimator specified in \(\mathrm{def:finite\mbox{-}set\mbox{-}estimator}\).
  \item \texttt{H\_\allowbreak{}m} --- \texttt{conditional high-\allowbreak{}probability estimator-\allowbreak{}validity event}
  \par\smallskip\noindent\emph{Definition.} Under a valid source-prime chart and successful conditional finite output, \(\mathcal H_m\) is the event that \(S_m\) lies in the projection neighborhood, every conditional certificate circuit is valid at \(\pi_G(S_m)\), and all retained roots are matched to their population branches.
  \item \texttt{A\_\allowbreak{}m} --- \texttt{conditional simultaneous branch confidence neighborhoods}
  \par\smallskip\noindent\emph{Definition.} Under a valid source-prime chart and successful conditional finite output, \(\mathcal A_m\) is the image under the smooth root map \(r\) of a joint asymptotic confidence ellipsoid for \(\pi_G(S_m)\), with one neighborhood for each retained branch.
  \item \texttt{tau\_\allowbreak{}b} --- \texttt{maximum labeled-\allowbreak{}tournament score-\allowbreak{}fiber cardinality}; \(\mathbb N\)
  \par\smallskip\noindent\emph{Definition.} \(\tau_b\) is the maximum prescribed-indegree coefficient defined in \(\mathrm{def:tournament\mbox{-}score\mbox{-}fiber}\) for an integer \(b\geq1\).
  \item \texttt{vartheta\_\allowbreak{}Ng} --- \texttt{joint vertex/\allowbreak{}cycle-\allowbreak{}rank orientation frontier}; \(\mathbb N\)
  \par\smallskip\noindent\emph{Definition.} \(\vartheta_{N,g}\) is the maximum prescribed-indegree orientation-fiber cardinality among connected simple labeled graphs on at most \(N\) vertices and with cyclomatic number at most \(g\), as specified in \(\mathrm{def:joint\mbox{-}orientation\mbox{-}frontier}\).
\end{itemize}

\section{Assumptions}

\begin{assumption}[ass:size]\label{ass:size}
\(n\ge2\).
\par\smallskip\noindent\emph{Standard} (Nontrivial finite SEM size, \cite{GuptaBlaser2024TreeSCM}).
\end{assumption}

\begin{assumption}[ass:k-domain]\label{ass:k-domain}
\(k\in\{0,\ldots,(n-1)(n-2)/2\}\).
\par\smallskip\noindent\emph{Novel.} The index records a transparent departure from a supplied hierarchy: applied path diagrams often add a small number of cross-level parents to a rooted causal tree.
\end{assumption}

\begin{assumption}[ass:acyclic]\label{ass:acyclic}
The directed graph \( (V,D) \) is acyclic.
\par\smallskip\noindent\emph{Standard} (Recursive linear structural equation model, \cite{FoygelDraismaDrton2012HTC}).
\end{assumption}

\begin{assumption}[ass:arborescence]\label{ass:arborescence}
\(T\subseteq D\) is a rooted spanning arborescence on \(V\).
\par\smallskip\noindent\emph{Standard} (Rooted tree-shaped directed support, \cite{GuptaBlaser2024TreeSCM}).
\end{assumption}

\begin{assumption}[ass:excess-budget]\label{ass:excess-budget}
\(|D\setminus T|\le k\).
\par\smallskip\noindent\emph{Novel.} The condition permits arbitrary placement of a bounded number of extra parents while retaining an interpretable hierarchical backbone, including skip-level and converging-parent structures used in path analysis.
\end{assumption}

\begin{assumption}[ass:noise-support]\label{ass:noise-support}
\(\Omega_{ij}=0\) whenever \(i\ne j\) and \(i\leftrightarrow j\notin B\).
\par\smallskip\noindent\emph{Standard} (Mixed-graph disturbance support, \cite{FoygelDraismaDrton2012HTC}).
\end{assumption}

\begin{assumption}[ass:iid-gaussian]\label{ass:iid-gaussian}
\(X^{(1)},\ldots,X^{(m)}\) are independent draws from \(N(0,\Sigma)\).
\par\smallskip\noindent\emph{Standard} (Independent Gaussian SEM sampling, \cite{FoygelDraismaDrton2012HTC}).
\end{assumption}

\section{Definitions}

\begin{definition}[def:graph-class]\label{def:graph-class}
\par\smallskip\noindent\emph{Defined object.} Tree-plus-\(k\) graph-query class
\par\smallskip\noindent\emph{Construction.} \(\mathcal G_{n,k}=\{(G,T,e):|V|=n,\ e\in D,\ (V,D)\text{ is acyclic},\ T\subseteq D\text{ is a rooted spanning arborescence},\ |D\setminus T|\le k\}\) for every admissible \(n\) and \(k\).
\par\smallskip\noindent\emph{Member properties.} \texttt{ass:\allowbreak{}size}; \texttt{ass:\allowbreak{}k-\allowbreak{}domain}; \texttt{ass:\allowbreak{}acyclic}; \texttt{ass:\allowbreak{}arborescence}; \texttt{ass:\allowbreak{}excess-\allowbreak{}budget}.
\end{definition}

\begin{definition}[def:generic-fiber]\label{def:generic-fiber}
\par\smallskip\noindent\emph{Defined object.} \texttt{Generic residual fiber algebra}
\par\smallskip\noindent\emph{Construction.} Let \(\phi_G:R_{\sigma}\to R_{\mathrm{src}}\) send each covariance coordinate to its acyclic SEM polynomial, put \(P_G=\ker\phi_G\), extend it to \(\psi_G:R_{\mathrm{inc}}=R_{\sigma}[\lambda_D]\to R_{\mathrm{src}}\) by fixing \(\lambda_D\), and put \(I_G=\ker\psi_G=(F_{ij}:i\ne j,\ i\leftrightarrow j\notin B)\). With \(R_G=R_{\sigma}/P_G\), \(K_G=\operatorname{Frac}(R_G)\), and the exact multiplicative set \(S_G=R_{\sigma}\setminus P_G\), define \(A_K=S_G^{-1}(R_{\mathrm{inc}}/I_G)\cong K_G\otimes_{R_G}(R_{\mathrm{inc}}/I_G)\). Equivalently, the covariance saturation is \(I_G:S_G^{\infty}=I_G\); no polynomial involving a directed source coordinate is inverted.
\par\smallskip\noindent\emph{Inputs.} \texttt{G}; \texttt{Sigma}.
\end{definition}

\begin{definition}[def:frontier-handle]\label{def:frontier-handle}
\par\smallskip\noindent\emph{Defined object.} \texttt{Size-\allowbreak{}sensitive orientation-\allowbreak{}realization handle}
\par\smallskip\noindent\emph{Construction.} Freeze a coordinate transcendence basis \(U\), select square residual-Jacobian subsystems blockwise by child vertex, encode their multihomogeneous Bezout coefficients as prescribed-indegree orientation counts, and determine exactly which counts are realized in the prime incidence algebra after localization at \(S_G=R_{\sigma}\setminus P_G\), with one original edge coordinate separating all finite branches; the output is a recurrence in \(n,k\), not a guessed exponent.
\par\smallskip\noindent\emph{Inputs.} \texttt{G\_\allowbreak{}nk}.
\end{definition}

\begin{definition}[def:census-handle]\label{def:census-handle}
\par\smallskip\noindent\emph{Defined object.} \texttt{Rooted five-\allowbreak{}node incidence handle}
\par\smallskip\noindent\emph{Construction.} Enumerate rooted graph-query isomorphism orbits with \(n\le5\) and \(k\le2\); for each orbit compute the characteristic-zero queried-edge elimination ideal from \(A_K=S_G^{-1}(R_{\mathrm{inc}}/I_G)\), homogenize before removing projective-infinity points, certify the affine squarefree degree by a rational univariate representation, and compare it with the prescribed-orientation coefficient. Lift the supplied modular quartic by rational reconstruction, exact substitution into \(I_G\), and verification that every inverted covariance factor lies in \(S_G\).
\par\smallskip\noindent\emph{Inputs.} \texttt{G\_\allowbreak{}nk}.
\end{definition}

\begin{definition}[def:attaining-family]\label{def:attaining-family}
\par\smallskip\noindent\emph{Defined object.} \texttt{One-\allowbreak{}sink family}
\par\smallskip\noindent\emph{Construction.} For \(d=k+1\) and \(n\ge2k+4\), take vertices \(0,P_1,\ldots,P_d,Q_1,\ldots,Q_d,h\), directed edges \(0\to P_i\), \(0\to Q_i\), and \(P_i\to h\), supply \(T=\{0\to P_i,0\to Q_i:1\le i\le d\}\cup\{P_1\to h\}\), omit exactly \(P_i\leftrightarrow Q_i\) and every ordinary-vertex pair with \(h\), allow all other bidirected pairs, set \(e^{\star}=P_1\to h\), and append any remaining vertices as fully bidirected children of \(0\) whose tree arrows are added to \(T\).
\par\smallskip\noindent\emph{Inputs.} \texttt{k}; \texttt{n}.
\end{definition}

\begin{definition}[def:attaining-membership-check]\label{def:attaining-membership-check}
\par\smallskip\noindent\emph{Defined object.} \texttt{One-\allowbreak{}sink family membership verification}
\par\smallskip\noindent\emph{Construction.} For \(d=k+1\), the displayed backbone has \(2d+1=2k+3\) directed edges on \(2d+2=2k+4\) vertices. Exactly \(d-1=k\) arrows \(P_i\to h\) lie outside the supplied backbone, and each appended vertex adds one tree edge and no excess arrow. Thus the supplied edge set is a rooted spanning arborescence, the directed support is acyclic, \(e^{\star}\in D\), and \((G^{\star}_{k,n},T,e^{\star})\in\mathcal G_{n,k}\). This is an inline support verification, not an identification target.
\par\smallskip\noindent\emph{Inputs.} \texttt{G\_\allowbreak{}star}; \texttt{e\_\allowbreak{}star}; \texttt{k}; \texttt{n}.
\end{definition}

\begin{definition}[def:quadratic-core]\label{def:quadratic-core}
\par\smallskip\noindent\emph{Defined object.} \texttt{Source-\allowbreak{}integral Möbius quadratic core}
\par\smallskip\noindent\emph{Construction.} Expose every incoming coefficient at a vertex of indegree at least two; on the other vertices factor rank-one residual relations in the integral source field, propagate rank-two relations by Möbius maps, record every genuine pole and every identically zero-over-zero interaction, and retain the resulting degree-at-most-two equations in at most \(2k\) exposed coefficients as \(\mathsf{Core}_{G,e}\).
\par\smallskip\noindent\emph{Inputs.} \texttt{G}; \texttt{T}; \texttt{e}.
\end{definition}

\begin{definition}[def:certificate-algorithm]\label{def:certificate-algorithm}
\par\smallskip\noindent\emph{Defined object.} \texttt{Zero-\allowbreak{}excess exact and positive-\allowbreak{}excess tag dispatcher}
\par\smallskip\noindent\emph{Construction.} Let \(r=|D\setminus T|\). In \(\mathsf{ArithmeticModel}\), dispatch as follows. If \(r=0\), invoke the complete published tree procedure on the equal class \(\mathcal G_{n,0}\). When \(d_e(G)<\infty\), return \((\operatorname{tree},d_e(G),(\rho_i)_{i=1}^{d_e(G)})\), where \(d_e(G)\in\{1,2\}\) and the list contains all returned \(\operatorname{FASTP}\) expressions; when \(d_e(G)=\infty\), return \((\operatorname{tree},\infty)\). If \(r>0\), construct the acyclic covariance parametrization in source coordinates, form its Jacobian and the Jacobian augmented by the row of \(\lambda_e\), and compute their generic ranks by independently randomized exact source specializations. When the ranks agree, return \((\operatorname{algebraic},\mathcal J_{G,e},2^{\min\{|D|-n+2,\binom{n-2}{2}\}})\); when the augmented rank is larger, return \((\infty,\mathcal N_{G,e})\). The positive-excess algebraic branch reports only a rank transcript and the size-sensitive universal degree upper bound and does not compute \(d_e(G)\), \(p_e\), \(\mathcal Q_{G,e}\), or \(h_{G,e}\). The tree branch does not claim to supply \(\mathcal Q_{G,e}\), \(h_{G,e}\), or the conditional source-prime-chart certificate.
\par\smallskip\noindent\emph{Inputs.} \texttt{G}; \texttt{T}; \texttt{e}; \texttt{delta}.
\end{definition}

\begin{definition}[def:exceptional-set]\label{def:exceptional-set}
\par\smallskip\noindent\emph{Defined object.} \texttt{Conditional certificate exceptional set}
\par\smallskip\noindent\emph{Construction.} Assume \(d_e(G)<\infty\), a valid source-prime chart package \(\mathfrak C_{G,e}\), and the successful finite-output event of \(\mathrm{thm:conditional\mbox{-}geometric\mbox{-}resolution}\). Define \(\mathcal E_{G,e}=V(h_{G,e})=\{\Sigma:h_{G,e}(\Sigma)=0\}\), where \(h_{G,e}\) is the product of the nonzero covariance-localization, incidence-lifting, pole, projective-leading-form, specialized-leading-coefficient, squarefreeness, and real-lifting circuits returned by the conditional algorithm.
\par\smallskip\noindent\emph{Inputs.} \texttt{Q\_\allowbreak{}Ge}; \texttt{Chart\_\allowbreak{}Ge}.
\end{definition}

\begin{definition}[def:real-retained-set]\label{def:real-retained-set}
\par\smallskip\noindent\emph{Defined object.} \texttt{Real extendable coordinate set}
\par\smallskip\noindent\emph{Construction.} \(\mathcal R_e(G,\Sigma)=\{z\in\mathbb R:\exists\Lambda\in\mathbb R^D,\ \lambda_e=z,\ F_{ij}(\Lambda,\Sigma)=0\text{ for every }i\leftrightarrow j\notin B\}\).
\par\smallskip\noindent\emph{Inputs.} \texttt{G}; \texttt{Sigma}; \texttt{e}.
\end{definition}

\begin{definition}[def:regular-cell]\label{def:regular-cell}
\par\smallskip\noindent\emph{Defined object.} \texttt{Conditional connected regular real-\allowbreak{}model cell}
\par\smallskip\noindent\emph{Construction.} For fixed \(G,e\) with \(d_e(G)<\infty\), a supplied valid source-prime chart package \(\mathfrak C_{G,e}\), and a successful finite output from \(\mathrm{thm:conditional\mbox{-}geometric\mbox{-}resolution}\), \(\mathcal C_{G,e}\) is a connected smooth coordinate neighborhood inside one manifold of a finite smooth-manifold stratification of the semialgebraic partition of \(\mathcal M_G\setminus\mathcal E_{G,e}\) induced by the signs used in \(\mathsf{LiftCert}\), on which the number and extendability labels of the simple real roots of \(p_e\) are constant.
\par\smallskip\noindent\emph{Inputs.} \texttt{M\_\allowbreak{}G}; \texttt{E\_\allowbreak{}Ge}; \texttt{LiftCert}; \texttt{Chart\_\allowbreak{}Ge}.
\end{definition}

\begin{definition}[def:regular-class]\label{def:regular-class}
\par\smallskip\noindent\emph{Defined object.} \texttt{Conditional compact regular finite-\allowbreak{}fiber class}
\par\smallskip\noindent\emph{Construction.} For fixed \(G,n,e\) with \(d_e(G)<\infty\), a supplied valid source-prime chart package \(\mathfrak C_{G,e}\), and a successful finite output from \(\mathrm{thm:conditional\mbox{-}geometric\mbox{-}resolution}\), select \(\mathcal C_{G,e}\) as in \(\mathrm{def:regular\mbox{-}cell}\) and let \(\mathcal K_G\subset\mathcal C_{G,e}\) be any nonempty compact set with nonempty relative interior in that cell.
\par\smallskip\noindent\emph{Inputs.} \texttt{C\_\allowbreak{}Ge}; \texttt{Chart\_\allowbreak{}Ge}.
\end{definition}

\begin{definition}[def:finite-set-estimator]\label{def:finite-set-estimator}
\par\smallskip\noindent\emph{Defined object.} \texttt{Conditional projected-\allowbreak{}covariance root-\allowbreak{}set estimator}
\par\smallskip\noindent\emph{Construction.} Assume \(d_e(G)<\infty\), a supplied valid source-prime chart package \(\mathfrak C_{G,e}\), and a successful finite output from \(\mathrm{thm:conditional\mbox{-}geometric\mbox{-}resolution}\). If \(S_m\) lies in the tubular neighborhood of \(\mathcal K_G\), \(\pi_G(S_m)\notin\mathcal E_{G,e}\), and the conditional-output circuits evaluate validly, set
\[
\widehat{\mathcal R}_{e,m}=\{z\in\mathbb R:p_e(\pi_G(S_m),z)=0\text{ and }\mathsf{LiftCert}(z,\pi_G(S_m))\text{ accepts}\},
\]
retaining every accepted branch and selecting none; otherwise set \(\widehat{\mathcal R}_{e,m}=\varnothing\).
\par\smallskip\noindent\emph{Inputs.} \texttt{S\_\allowbreak{}m}; \texttt{Chart\_\allowbreak{}Ge}; \texttt{p\_\allowbreak{}e}; \texttt{LiftCert}.
\end{definition}

\begin{definition}[def:branch-neighborhoods]\label{def:branch-neighborhoods}
\par\smallskip\noindent\emph{Defined object.} \texttt{Conditional simultaneous branch neighborhoods}
\par\smallskip\noindent\emph{Construction.} Assume \(d_e(G)<\infty\), a supplied valid source-prime chart package \(\mathfrak C_{G,e}\), and a successful finite output from \(\mathrm{thm:conditional\mbox{-}geometric\mbox{-}resolution}\). Construct \(\mathcal A_m\) by mapping a joint Gaussian covariance ellipsoid for \(\pi_G(S_m)\) through every locally labeled component of \(r\), using the implicit derivative
\[
D r_j(\Sigma)[H]=-\frac{D_\Sigma p_e(\Sigma,r_j)[H]}{\partial_t p_e(\Sigma,r_j)}.
\]
\par\smallskip\noindent\emph{Inputs.} \texttt{Rhat\_\allowbreak{}m}; \texttt{V\_\allowbreak{}Sigma}; \texttt{Chart\_\allowbreak{}Ge}; \texttt{p\_\allowbreak{}e}.
\end{definition}

\begin{definition}[def:source-prime-chart]\label{def:source-prime-chart}
\par\smallskip\noindent\emph{Defined object.} \texttt{Localized reduced source-\allowbreak{}prime chart package}
\par\smallskip\noindent\emph{Construction.} For each source-consistent recorded stratum of \(\mathsf{Core}_{G,e}\), let \(x=(x_1,\ldots,x_r)\), with \(r\leq2k\), be its exposed coordinates, let \(\xi:K_G[x]\to L\) be source evaluation, and put \(P=\ker\xi\). A source-prime chart package \(\mathfrak C_{G,e}\) supplies a product \(H\in K_G[x]\), elements \(h_1,\ldots,h_c\in P\) with \(c=\operatorname{ht}P\), a \(c\)-by-\(c\) Jacobian minor \(\Delta\) of the \(h_i\), and rational lifting circuits for all eliminated directed coordinates. It is valid when: \(\xi(H\Delta)\neq0\); \(h_1,\ldots,h_c\) form a reduced regular sequence after localization at \(H\Delta\); 
\[
P K_G[x]_{H\Delta}=(h_1,\ldots,h_c)K_G[x]_{H\Delta};
\]
and the lifting circuits satisfy every original residual equation in \(A_K\) and recover \(\lambda_e\). All chart coefficients are rational arithmetic circuits in the covariance entries. The total chart straight-line-program size is denoted by \(s(\mathfrak C_{G,e})\), and \(\mathfrak d(\mathfrak C_{G,e})\) denotes the maximum geometric degree encountered by successive chart intersections, localizations, slices, and lifting fibers.
\par\smallskip\noindent\emph{Inputs.} \texttt{G}; \texttt{T}; \texttt{e}; \texttt{Core\_\allowbreak{}Ge}.
\end{definition}

\begin{definition}[def:tournament-score-fiber]\label{def:tournament-score-fiber}
\par\smallskip\noindent\emph{Defined object.} \texttt{Maximum tournament score-\allowbreak{}vector fiber}
\par\smallskip\noindent\emph{Construction.} For every integer \(b\geq 1\), define
\[
 c_b(a_1,\ldots,a_b)
 :=[X_1^{a_1}\cdots X_b^{a_b}]
 \prod_{1\leq i<j\leq b}(X_i+X_j),
 \qquad
 \tau_b:=\max_{a\in\mathbb Z_{\geq0}^{b}}c_b(a).
\]
Choosing \(X_i\) from the factor indexed by \(\{i,j\}\) means orienting that edge toward \(i\). Thus \(c_b(a)\) is the number of labeled tournaments on vertex set \([b]\) having prescribed indegree vector \(a\), and \(\tau_b\) is the largest such score-vector fiber.
\end{definition}

\begin{definition}[def:joint-orientation-frontier]\label{def:joint-orientation-frontier}
\par\smallskip\noindent\emph{Defined object.} \texttt{Joint vertex/\allowbreak{}cycle-\allowbreak{}rank orientation frontier}
\par\smallskip\noindent\emph{Construction.} For integers \(N\geq1\) and \(g\geq0\), let \(H\) range over connected simple graphs on a labeled vertex set \([b]\), where \(1\leq b\leq N\), and whose cyclomatic number is
\[
 g(H)=|E(H)|-b+1\leq g.
\]
For \(r=(r_1,\ldots,r_b)\in\mathbb Z_{\geq0}^{b}\), put
\[
 c_H(r)=[X_1^{r_1}\cdots X_b^{r_b}]
          \prod_{\{i,j\}\in E(H)}(X_i+X_j),
 \qquad
 \vartheta_{N,g}:=\max_{H,r}c_H(r).
\]
The one-vertex graph is included: its empty product is \(1\), its cyclomatic number is zero, and \(c_H(0)=1\).
\end{definition}

\section{Main results}

\begin{theorem}[thm:incidence-localization]\label{thm:incidence-localization}
The homomorphism \(\psi_G:R_{\mathrm{inc}}\to R_{\mathrm{src}}\) is surjective, \(I_G=\ker\psi_G=(F_{ij}:i\ne j,\ i\leftrightarrow j\notin B)\) is prime, and \(R_{\mathrm{inc}}/I_G\cong R_{\mathrm{src}}\). Its covariance restriction has kernel \(P_G\), so \(R_G\cong\operatorname{im}\phi_G\), \(K_G=\operatorname{Frac}(R_G)\subseteq L\), and \(I_G:S_G^{\infty}=I_G\) for the exact multiplicative set \(S_G=R_{\sigma}\setminus P_G\). Consequently \(A_K=S_G^{-1}(R_{\mathrm{inc}}/I_G)\) is integral with fraction field \(L\), and covariance localization deletes no generic point or finite \(\lambda_e\)-value: if \(\lambda_e\) is algebraic, its affine coordinate image over \(\overline{K_G}\) is exactly the root set of its minimal polynomial, even when the full generic fiber is positive dimensional; if it is transcendental, the coordinate image is Zariski dense and infinite. No directed-coordinate expression is inverted, so poles and identically zero-over-zero components created by intermediate rational formulas remain in the incidence algebra. Homogenized elimination records projective-infinity points separately, while affine dehomogenization retains exactly the finite coordinate image; denominator zeros and leading-coefficient drops are specialization loci and do not change the generic affine image.
\end{theorem}
\begin{proof}
Put \(R=I_n-\Lambda\). By acyclicity, \(\Lambda^n=0\) after a topological ordering, and hence \(R^{-1}=I_n+\Lambda+\cdots+\Lambda^{n-1}\). Thus the covariance parametrization is polynomial.

Let \(J=(F_{ij}:i\ne j,\ i\leftrightarrow j\notin B)\). In \(R_{\mathrm{inc}}/J\), set
\[
\overline\Omega=R(\overline\Lambda)^{\mathsf T}\overline\Sigma R(\overline\Lambda).
\]
For a missing pair, direct matrix multiplication gives
\[
\overline\Omega_{ij}=\sigma_{ij}
-\sum_{p\in\operatorname{pa}_D(i)}\lambda_{pi}\sigma_{pj}
-\sum_{q\in\operatorname{pa}_D(j)}\lambda_{qj}\sigma_{iq}
+\sum_{p,q}\lambda_{pi}\lambda_{qj}\sigma_{pq}=F_{ij}=0.
\]
Consequently the supported entries of \(\overline\Omega\), together with \(\overline\Lambda\), define a homomorphism \(\theta:R_{\mathrm{src}}\to R_{\mathrm{inc}}/J\). The map induced by \(\psi_G\) is its inverse, because
\[
R(\Lambda)^{\mathsf T}\{R(\Lambda)^{-\mathsf T}\Omega R(\Lambda)^{-1}\}R(\Lambda)=\Omega
\]
and \(R(\overline\Lambda)^{-\mathsf T}\overline\Omega R(\overline\Lambda)^{-1}=\overline\Sigma\). Hence
\[
R_{\mathrm{inc}}/J\cong R_{\mathrm{src}},\qquad J=\ker\psi_G=I_G.
\]
This proves surjectivity, primality, and the quotient assertion. Restriction gives \(P_G=\ker\phi_G\) and \(R_G\cong\operatorname{im}\phi_G\).

The source ring is a domain. Every member of \(S_G\) has nonzero image in it, so \(I_G:S_G^\infty=I_G\) and \(A_K\) is a domain. Moreover \(S_G^{-1}R_G=K_G\), while localization leaves the fraction field unchanged; therefore \(\operatorname{Frac}(A_K)=L\). No directed-coordinate expression was inverted.

Suppose \(\lambda_e\) is algebraic and put \(E=K_G(\lambda_e)\). Bezout's identity modulo its irreducible minimal polynomial shows that every nonzero element of \(K_G[\lambda_e]\) is invertible there, so \(E=K_G[\lambda_e]\subset A_K\). For every \(K_G\)-embedding \(\tau:E\to\overline{K_G}\),
\[
A_K\otimes_{E,\tau}\overline{K_G}\ne0
\]
by faithful flatness. This finitely generated algebra has a geometric point with coordinate \(\tau(\lambda_e)\). Conversely, every geometric point restricts to such an embedding. Thus the coordinate image is exactly the roots of the minimal polynomial, irrespective of the dimension of the full fiber. If \(\lambda_e\) is transcendental, \(K_G[t]\to A_K\), \(t\mapsto\lambda_e\), is injective; the coordinate morphism is dominant, so its image is Zariski dense and infinite.

Finally, the nonzero homogenizing chart dehomogenizes exactly to the affine incidence scheme; the complementary hyperplane consists exactly of projective-infinity points. A denominator zero or leading-coefficient drop is the zero set of a nonzero covariance element. Recording these loci therefore preserves the generic affine coordinate image.
\end{proof}
\par\smallskip\noindent \textit{Justification.} The explicit prime incidence algebra and covariance-only localization prevent coordinate branches from being silently removed during elimination. \textit{Gap.} Foygel, Draisma, and Drton factor full generic fibers, while the stated theorem identifies the exact covariance multiplicative set and preserves one coordinate through positive-dimensional fibers, poles, zero-over-zero components, and projective specialization. \textit{Consumer.} The minimal-polynomial routine can invert certified covariance-model classes while retaining every generic value of the requested coefficient and auditing every affine or projective exclusion.

\begin{theorem}[thm:orientation-upper-bound]\label{thm:orientation-upper-bound}
For every \((G,T,e)\in\mathcal G_{n,k}\) with \(d_e(G)<\infty\),
\[
 d_e(G)
 \leq \vartheta_{n-1,\,|D|-n+2}
 \leq \vartheta_{n-1,\,k+1}.
\]
More precisely,
\[
 d_e(G)\leq\min\!\left\{2^{|D|-n+2},\tau_{n-1}\right\},
 \qquad
 \vartheta_{n-1,k+1}\leq\min\!\left\{2^{k+1},\tau_{n-1}\right\}.
\]
These inequalities hold even when the full generic parameter fiber is positive dimensional.  Consequently,
\[
 D_k(n)\leq\vartheta_{n-1,k+1}\leq\min\!\left\{2^{k+1},\tau_{n-1}\right\}.
\]
For six vertices this yields
\[
 D_2(6)\leq6,\qquad D_3(6)\leq8,\qquad D_4(6)\leq12,
 \qquad D_k(6)\leq24\quad\text{for every admissible }k\geq5.
\]
\end{theorem}
\begin{proof}
Fix \((G,T,e)\in\mathcal G_{n,k}\) and suppose \(d_e(G)<\infty\).  Multiplication of the covariance identity by \((I_n-\Lambda)^{\mathsf T}\) on the left and by \(I_n-\Lambda\) on the right gives
\[
 \Omega=(I_n-\Lambda)^{\mathsf T}\Sigma(I_n-\Lambda).
\tag{66}
\]
Thus every disturbance coordinate lies in \(K_G(\lambda_D)\).  The reverse inclusion follows from \(K_G\subseteq L\) and \(\lambda_D\subseteq L\), so
\[
 L=K_G(\lambda_D).
\tag{67}
\]

Choose \(U\subseteq\lambda_D\setminus\{\lambda_e\}\) maximal algebraically independent over \(K_G\), put \(F=K_G(U)\), and enumerate the remaining directed coordinates as \(y=(y_1,\ldots,y_q)\), with \(\lambda_e\) among them.  Maximality, finiteness of \(\lambda_D\), the assumed algebraicity of \(\lambda_e\), and (67) imply that \(L/F\) is finite and, in characteristic zero, separable.

Let \(p\in K_G[t]\) be the monic minimal polynomial of \(\lambda_e\).  It remains irreducible over the purely transcendental extension \(F/K_G\).  Indeed, a factorization over \(K_G(U)[t]\) can be cleared of denominators and made primitive in the unique-factorization domain \(K_G[U]\); comparing degrees in the indeterminates \(U\) then gives a factorization over \(K_G[t]\), contrary to irreducibility.  Therefore
\[
 [F(\lambda_e):F]=[K_G(\lambda_e):K_G]=d_e(G).
\tag{68}
\]

By \(\mathrm{thm:incidence\mbox{-}localization}\) and \(\mathrm{def:generic\mbox{-}fiber}\), first adjoining \(U\) and then passing to its fraction field gives an evaluation homomorphism
\[
 \rho:F[y_1,\ldots,y_q]\longrightarrow L
\tag{69}
\]
whose kernel is generated by the images of the missing-bidirected residuals \(F_{ij}\).  The map is surjective: (67) says its image generates \(L\) as a field, while all \(y_a\) are algebraic over \(F\), so the image is a finite-dimensional domain over \(F\), hence a field.

For every \(a\), let \(g_a\in F[t]\) be the minimal polynomial of \(y_a\).  Because \(g_a(y_a)\in\ker\rho\), there exist \(c_{a,ij}\in F[y]\) such that
\[
 g_a(y_a)=\sum_{i\leftrightarrow j\notin B}c_{a,ij}(y)F_{ij}(y).
\tag{70}
\]
Differentiate (70) in the free cotangent space \(\Omega_{F[y]/F}\otimes_{F[y],\rho}L\).  Every residual vanishes after applying \(\rho\), so
\[
 g_a'(y_a)\,dy_a
 =\sum_{i\leftrightarrow j\notin B}c_{a,ij}(y)\,dF_{ij}.
\tag{71}
\]
Separability gives \(g_a'(y_a)\ne0\).  Hence the evaluated residual differentials span all \(dy_a\), and the residual Jacobian with respect to \(y\) has rank \(q\).  Select distinct residuals \(f_1,\ldots,f_q\) for which
\[
 \Delta=\det\left(\frac{\partial f_\ell}{\partial y_a}\right)_{\ell,a=1}^{q}
\tag{72}
\]
is nonzero in \(L\).

Put \(E=F(\lambda_e)\).  By (68), \(E/F\) has exactly \(d=d_e(G)\) embeddings into an algebraic closure \(\overline F\).  Each extends to an \(F\)-embedding of \(L\): for an embedding \(\eta:E\to\overline F\), the nonzero finite algebra \(L\otimes_{E,\eta}\overline F\) has a maximal ideal, and its residue map gives such an extension.  Choose one extension for each \(\eta\).  Applying them to \(y\) gives \(d\) common zeros of \(f_1,\ldots,f_q\) over \(\overline F\).  Their Jacobians are nonsingular because every embedding sends the nonzero element \(\Delta\) to a nonzero element, and their \(\lambda_e\)-coordinates are pairwise distinct because their restrictions to \(E\) differ.

We now specialize only finitely many coefficients, while preserving this selected set.  Since \(L/F\) is finite separable, choose a primitive element \(\theta\), write \(L=F(\theta)\), let \(Q\in F[t]\) be its monic minimal polynomial, and put \(h=[L:F]\).  Express every \(y_a\), \(\Delta\), and \(\lambda_e\) as a rational polynomial in \(\theta\) of degree below \(h\).  All finitely many coefficients lie in the ambient rational function field \(L=\mathbb C(\lambda_D,\omega_B,\omega_{\mathrm{diag}})\).  Form the product of the numerators and denominators needed for those expressions and identities, together with \(\operatorname{disc}(Q)\), \(\operatorname{disc}(p)\), and the field norms of all expression denominators and of \(\Delta\).  Every factor is nonzero by construction, separability, and (72).  After clearing denominators this is a nonzero polynomial in the algebraically independent source coordinates.  Since \(\mathbb C\) is infinite, those coordinates admit a complex specialization at which that polynomial is nonzero.

At this specialization, \(Q\) has \(h\) distinct complex roots; all coordinate expressions are finite; every specialized residual identity vanishes; and the specialized Jacobian determinant is nonzero at every resulting point.  Multiplication by \(\lambda_e\) on the \(F\)-vector space \(L\) has characteristic polynomial
\[
 \det(tI-m_{\lambda_e})=p(t)^{h/d},
\tag{73}
\]
because each of the \(d\) embeddings of \(E/F\) has exactly \([L:E]=h/d\) extensions to \(L\).  Identity (73) survives the specialization, and the inverted discriminant of \(p\) makes its \(d\) specialized roots distinct.  Hence the specialized square system over \(\mathbb C\) contains a selected set of \(d=d_e(G)\) distinct nonsingular affine zeros with pairwise distinct \(\lambda_e\)-coordinates.

Partition \(y\) into its nonempty blocks indexed by child vertices, and let \(r_v\) be the number of coordinates in the block indexed by \(v\).  Then
\[
 1\leq r_v\leq\operatorname{indeg}_D(v).
\tag{74}
\]
Each residual \(F_{ij}\) is affine-linear in the incoming-coefficient block of \(i\), affine-linear in that of \(j\), and independent of every other child block.  Once the coordinates in \(U\) are coefficients in \(F\), a selected residual is therefore unary or binary, with degree at most one in each incident block.

Decompose the bipartite equation--block incidence graph into connected components.  A nonzero determinant term in (72) is a perfect row--variable matching.  Because every matched row and variable lie in the same incidence component, each component has equally many scalar equations and scalar variables, and the Jacobian becomes block diagonal with a nonsingular square block after permutation.  Consider the component containing \(\lambda_e\).  Its projections of the \(d\) selected affine zeros remain distinct, because their \(\lambda_e\)-coordinates are distinct, and each projected zero has nonsingular component Jacobian.

Let this component contain \(b\) variable blocks, \(u\) unary equations, and \(m_H\) binary equations.  If \(b>1\), its binary equations form a connected simple graph \(H\): unary equations connect no two blocks, and an unordered pair of distinct child vertices supplies at most the one residual \(F_{ij}\).  If \(b=1\), let \(H\) be the one-vertex graph.  Squareness gives
\[
 \sum_{v\in V(H)}r_v=u+m_H.
\tag{75}
\]

Apply \(\mathrm{lem:finite\mbox{-}nonsingular\mbox{-}perturbation\mbox{-}bound}\) to precisely these \(d\) selected nonsingular affine zeros, with unary block-degree vector \(X_v\) and binary block-degree vector \(X_v+X_w\).  This is the paper-owned bridge from a possibly positive-dimensional subsystem to the cited finite-system multihomogeneous B\'ezout theorem: its generic perturbation stays within the same block multidegrees, block-homogenizes the system, makes the perturbed projective intersection finite, and uses the implicit-function theorem in neighborhoods contained in the affine charts where all homogenizing coordinates equal \(1\).  Thus each selected zero continues to a distinct affine zero, and its distinct \(\lambda_e\)-coordinate remains in a disjoint neighborhood; no selected queried-coordinate branch is lost to or imported from projective infinity.  The lemma then gives exactly
\[
 d_e(G)\leq
 [\prod_vX_v^{r_v}]
 \left(\prod_{\ell=1}^{u}X_{v(\ell)}\right)
 \left(\prod_{\{v,w\}\in E(H)}(X_v+X_w)\right).
\tag{76}
\]
The underlying cited multihomogeneous B\'ezout result is used only inside that bridge, after the perturbed common-zero set has been made finite; no isolated-zero extension of the cited result is asserted.

Let \(u_v\) be the number of unary equations in block \(v\), and put \(s_v=r_v-u_v\).  The coefficient in (76) must be nonzero because its left side is positive; hence each \(s_v\geq0\).  Removing the forced unary powers from (76) leaves
\[
 [\prod_vX_v^{s_v}]\prod_{\{v,w\}\in E(H)}(X_v+X_w)=c_H(s),
\tag{77}
\]
the number of orientations of \(H\) having prescribed indegree vector \(s\).  Therefore
\[
 d_e(G)\leq c_H(s).
\tag{78}
\]

The root of the supplied arborescence has no incoming directed edge: otherwise that edge together with the directed tree path from the root to its tail would form a directed cycle, contradicting \(\mathrm{ass:acyclic}\).  It supplies no nonempty child block, so
\[
 b\leq n-1.
\tag{79}
\]
For \(b>1\), the cyclomatic number of \(H\) is \(g_H=m_H-b+1\).  By (74)--(75), nonnegativity of \(u\), and the incoming tree edge at every nonroot vertex,
\[
\begin{aligned}
 g_H
 &=\sum_{v\in V(H)}r_v-u-b+1\\
 &\leq\sum_{v\in V(H)}\bigl(\operatorname{indeg}_D(v)-1\bigr)+1\\
 &\leq\sum_{v\ne\operatorname{root}(T)}
       \bigl(\operatorname{indeg}_D(v)-1\bigr)+1\\
 &=|D|-n+2.
\end{aligned}
\tag{80}
\]
The omitted summands are nonnegative because every nonroot vertex receives its tree edge.  For \(b=1\), set \(g_H=0\); then the same inequality follows from \(|D|\geq|T|=n-1\).  Equations (78)--(80) and \(\mathrm{def:joint\mbox{-}orientation\mbox{-}frontier}\) imply
\[
 d_e(G)\leq\vartheta_{n-1,\,|D|-n+2}.
\tag{81}
\]

Because \(T\subseteq D\), \(|T|=n-1\), and \(\mathrm{ass:excess\mbox{-}budget}\) holds,
\[
 |D|-n+2=|D\setminus T|+1\leq k+1.
\tag{82}
\]
Monotonicity and the bounds in \(\mathrm{lem:joint\mbox{-}orientation\mbox{-}frontier}\), together with \(\mathrm{lem:tournament\mbox{-}score\mbox{-}fiber\mbox{-}ceiling}\), now give
\[
 d_e(G)
 \leq\vartheta_{n-1,\,|D|-n+2}
 \leq\vartheta_{n-1,k+1}
 \leq\min\{2^{k+1},\tau_{n-1}\},
\tag{83}
\]
and, before (82),
\[
 d_e(G)\leq\min\{2^{|D|-n+2},\tau_{n-1}\}.
\tag{84}
\]
Maximizing (83) over all finite-degree members of \(\mathcal G_{n,k}\) yields
\[
 D_k(n)\leq\vartheta_{n-1,k+1}.
\tag{85}
\]
The argument counted only a finite selected set of nonsingular conjugate zeros and never assumed that the full generic parameter fiber was zero dimensional.

For \(n=6\), \(\mathrm{lem:joint\mbox{-}orientation\mbox{-}frontier}\) gives \(\vartheta_{5,3}=6\), \(\vartheta_{5,4}=8\), \(\vartheta_{5,5}=12\), and \(\vartheta_{5,g}=24\) for \(g\geq6\).  Substituting \(g=k+1\) in (85) proves
\[
 D_2(6)\leq6,\qquad D_3(6)\leq8,\qquad D_4(6)\leq12,
 \qquad D_k(6)\leq24\quad(k\geq5),
\tag{86}
\]
for every admissible \(k\).
\end{proof}
\par\smallskip\noindent \textit{Justification.} This gives a universal projection-aware ceiling while remaining valid for positive-dimensional full fibers: only a finite selected nonsingular conjugate set is counted. \textit{Gap.} The available multihomogeneous citation covers finite common-zero sets only, so the paper proves the same-block-multidegree perturbation bridge before applying that result. \textit{Consumer.} The result supplies the ambiguity budget and six-vertex caps used by the positive-excess dispatcher without claiming that every combinatorial orientation fiber is SEM-attained.

\begin{theorem}[thm:large-size-frontier]\label{thm:large-size-frontier}
For every \(k\ge0\) and \(n\ge2k+4\), \(d_{e^{\star}}(G^{\star}_{k,n})=2^{k+1}\) and therefore \[D_k(n)=2^{k+1}.\] One original coefficient \(e^{\star}\) separates all branches on a nonempty Zariski-open covariance set.
\end{theorem}
\begin{proof}
Put \(d=k+1\) and \(N=2d+2=2k+4\).  In the \(N\)-vertex graph of \(\mathrm{def:attaining\mbox{-}family}\), the arrows are
\[
 0\to P_i,\qquad 0\to Q_i,\qquad P_i\to h
 \quad(1\le i\le d).
\]
The supplied set consists of the first \(2d\) arrows and \(P_1\to h\).  Therefore
\[
 |T|=2d+1=N-1,\qquad
 D\setminus T=\{P_i\to h:2\le i\le d\},\qquad
 |D\setminus T|=d-1=k.
\tag{16}
\]
Every nonroot vertex has one incoming tree arrow, every vertex is reached from \(0\), and the order \(0,(P_i,Q_i)_{i=1}^d,h\) is topological.  Thus \((G^{\star}_{k,N},T,e^{\star})\in\mathcal G_{N,k}\).

Choose positive rational \(h_i\) and rational \(b_i\) satisfying \(b_i^2>4h_i\), and put \(H_0=\operatorname{diag}(h_1,\ldots,h_d)\).  Choose a rational symmetric positive-definite matrix \(H\), sufficiently close to \(H_0\), with the same diagonal; below we impose one further generic restriction on its off-diagonal entries.  In the order \((0,P_1,\ldots,P_d,Q_1,\ldots,Q_d,h)\), set
\[
 \Sigma^{\star}=
 \begin{pmatrix}
  1&0&0&1\\
  0&H&H&0\\
  0&H&C_0I_d&b\\
  1&0&b^{\mathsf T}&M
 \end{pmatrix},
 \qquad b=(b_1,\ldots,b_d)^{\mathsf T}.
\tag{17}
\]
Choose rational \(C_0\) with \(C_0I_d-H\succ0\).  The Schur complement of \(H\) proves that the \((P,Q)\)-principal block is positive definite.  Adding the independent \(0\)-coordinate preserves this property, and choosing rational \(M\) above the final Schur-complement threshold gives
\[
 \Sigma^{\star}\succ0.
\tag{18}
\]

Write
\[
 a_i=\lambda_{0P_i},\quad c_i=\lambda_{0Q_i},\quad
 z_i=\lambda_{P_i h},\quad z=(z_1,\ldots,z_d)^{\mathsf T},
 \quad x=-Hz.
\]
The missing pairs are \(P_i\leftrightarrow Q_i\), \(P_i\leftrightarrow h\), and \(Q_i\leftrightarrow h\).  Direct substitution of (17) in the definition of \(F_{ij}\) gives
\[
 0=-a_i-(Hz)_i,\qquad
 0=b_i-c_i-(Hz)_i,\qquad
 0=h_i+a_ic_i.
\tag{19}
\]
Since \(H\) is invertible, these equations are equivalent to
\[
 a_i=x_i,\qquad c_i=x_i+b_i,\qquad
 x_i^2+b_ix_i+h_i=0\quad(1\le i\le d),\qquad
 z=-H^{-1}x.
\tag{20}
\]
Every quadratic in (20) has positive discriminant, so there are exactly \(2^d\) real solutions.  Under the invertible change \(z\mapsto x=-Hz\), the residual Jacobian in \((a,c,z)\) is nonsingular exactly when
\[
 \det(H)\prod_{i=1}^d(2x_i+b_i)\ne0.
\tag{21}
\]
Indeed, for fixed \(i\) the three equations in \((a_i,c_i,x_i)\) have determinant, up to sign, \(a_i+c_i=2x_i+b_i\), and at a solution
\[
 (2x_i+b_i)^2=b_i^2-4h_i>0.
\]
Thus all \(2^d\) points are reduced.  Defining
\[
 \Omega=(I_N-\Lambda)^{\mathsf T}\Sigma^{\star}(I_N-\Lambda)
\tag{22}
\]
makes precisely the missing-pair entries vanish by (19).  The topological order makes \(I_N-\Lambda\) unit triangular, so (18) and congruence imply \(\Omega\succ0\).  Hence every solution is a real supported positive-definite source lift.

We next choose \(H\) so that \(z_1=\lambda_{e^{\star}}\) separates the solutions.  At \(H=H_0\), branches with different first coordinate \(x_1\) already have different \(z_1=-x_1/h_1\).  If branch vectors \(x\ne x'\) have the same first coordinate, choose \(j>1\) with \(x_j\ne x'_j\) and set \(H(s)=H_0+s(E_{1j}+E_{j1})\).  Differentiating the inverse gives
\[
 \left.\frac{d}{ds}\{z_1(x)-z_1(x')\}\right|_{s=0}
 =\frac{x_j-x'_j}{h_1h_j}\ne0.
\tag{23}
\]
Thus equality of the two first coordinates after applying \(-H^{-1}\) is a proper algebraic condition on the off-diagonal entries of \(H\).  There are finitely many branch pairs.  Their proper algebraic sets cannot cover a Euclidean neighborhood of \(H_0\), and rational matrices are dense, so a rational positive-definite \(H\) can be chosen outside their union.  Its \(2^d\) values of \(z_1\) are distinct.

There are \(3d\) directed coordinates and \(3d\) missing residuals.  Since all other pairs are bidirected,
\[
 |D|+|B|+|V|
 =3d+\left(\binom N2-3d\right)+N
 =\frac{N(N+1)}2=:M_0.
\tag{24}
\]
Perturb the independent covariance coordinates as \(\Sigma^{\star}+\xi\), where the \(M_0\) entries of \(\xi\) are algebraically independent.  At any reduced solution \(v_0=(a,c,z)\), seek a formal solution \(v(\xi)=v_0+\sum_{r\ge1}v_r(\xi)\), with \(v_r\) homogeneous of degree \(r\).  At degree \(r\), the residual equations read
\[
 J(v_0)v_r=-R_r(\xi,v_1,\ldots,v_{r-1}).
\tag{25}
\]
The Jacobian is invertible by (21), so (25) uniquely determines every \(v_r\).  Each displayed point therefore has a formal section over the full covariance formal neighborhood.

If \(f\in P_G\), substitution of any section gives \(f(\Sigma^{\star}+\xi)=0\) in \(\mathbb C[[\xi]]\).  Translation by \(\Sigma^{\star}\) is injective on the covariance polynomial ring and the coordinates of \(\xi\) are algebraically independent, so \(f=0\).  Hence
\[
 P_G=(0),\qquad K_G=\mathbb C(\sigma_{ij}:i\le j).
\tag{26}
\]
Equation (24) now shows that \(L\) and \(K_G\) have equal transcendence degree; hence the finitely generated extension \(L/K_G\) is finite.

The \(2^d\) formal sections constructed above define homomorphisms from the integral incidence ring to \(\mathbb C[[\xi]]\), all restricting to the same injective covariance translation. Each image has transcendence degree at least \(M_0\), equal by (24) to the dimension of the incidence domain, so each prime kernel has height zero and is zero. Passing to fraction fields produces \(2^d\) distinct \(K_G\)-embeddings of \(L\) into \(\operatorname{Frac}\mathbb C[[\xi]]\); they send \(z_1\) to series with distinct constant terms. Consequently
\[
 2^d\leq[K_G(z_1):K_G].
\tag{27}
\]
The universal bound in \(\mathrm{thm:orientation\mbox{-}upper\mbox{-}bound}\), applied with \(n=2k+4\), gives \([K_G(z_1):K_G]=d_{e^{\star}}(G^{\star}_{k,N})\leq2^{k+1}=2^d\). Therefore equality holds.

Thus \(d_{e^{\star}}(G^{\star}_{k,N})=2^{k+1}\), and the original coefficient \(e^{\star}\) separates all conjugate branches.  Characteristic zero makes the minimal polynomial separable.  Clearing its coefficient denominators and excluding their zeros and its nonzero discriminant gives a nonempty Zariski-open covariance set on which the roots remain distinct; \(\mathrm{thm:incidence\mbox{-}localization}\) identifies them with the finite queried-coordinate projection.

For \(n>N\), append a fully bidirected child \(w\) of \(0\), and place \(0\to w\) in \(T\).  No missing residual involves \(w\).  Conditional on the old source coordinates and \(x_w=\lambda_{0w}\), covariance congruence is an invertible triangular polynomial change between the new allowed disturbance coordinates and the new covariance row \(s_w\).  Hence
\[
 L_1=L_0(x_w,s_w),\qquad K_1=K_0(s_w),
\tag{29}
\]
with \(x_w\) and the entries of \(s_w\) algebraically independent over \(L_0\).  Irreducibility is preserved by the purely transcendental extension \(K_0(s_w)/K_0\), so
\[
 [K_1(z_1):K_1]=[K_0(z_1):K_0]=2^d.
\tag{30}
\]
Iteration proves the degree assertion for every \(n\ge2k+4\), while membership is preserved because each new arrow is a tree arrow.  Finally \(\mathrm{thm:orientation\mbox{-}upper\mbox{-}bound}\) gives \(D_k(n)\le2^{k+1}\), whereas (28)--(30) give the reverse inequality.  Therefore
\[
 D_k(n)=2^{k+1}\qquad(n\ge2k+4),
\]
with the claimed nonempty Zariski-open separation set.
\end{proof}
\par\smallskip\noindent \textit{Justification.} A uniform attaining region turns the exponent from a loose Bezout count into a sharp coordinate frontier for a substantial size range. \textit{Gap.} Foygel, Draisma, and Drton already exhibit cubic and quartic behavior in small SEMs, but not an all-\(k\) family with one edge separating \(2^{k+1}\) generic branches. \textit{Consumer.} Software maintainers obtain adversarial benchmark families and users obtain a worst-case ambiguity warning indexed by the number of added parents.

\begin{lemma}[lem:tree-maxima-support]\label{lem:tree-maxima-support}
The exact tree maxima are \(D_0(2)=D_0(3)=1\) and \(D_0(n)=2\) for \(n\ge4\). For \(n=2\), the sole missing-bidirected residual is either absent, giving a transcendental query, or linear in the unique edge coefficient. For \(n=3\), consider the star \(0\to1,0\to2\) and the chain \(0\to1\to2\). In either shape the parent--child residuals, when present, are linear anchors. The only remaining residual is bilinear in the two edge coefficients: alone it leaves each query transcendental, and after either linear anchor it solves the other coefficient linearly. The unconfounded chain supplies degree one. For \(n=4\), take the star \(0\to1,0\to2,0\to3\), include exactly the three bows \(0\leftrightarrow j\), and omit all three child--child bidirected edges. Writing \(a=\lambda_{01},b=\lambda_{02},c=\lambda_{03}\), its three residual equations are \(\sigma_{ij}-\lambda_{0i}\sigma_{0j}-\lambda_{0j}\sigma_{0i}+\lambda_{0i}\lambda_{0j}\sigma_{00}=0\) for \(1\le i<j\le3\). At the positive-definite source point \(a,b,c=(1,2,3)\), \(\omega_{00}=1\), \(\omega_{11}=\omega_{22}=\omega_{33}=10\), and \((\omega_{01},\omega_{02},\omega_{03})=(1/5,-1/7,1/11)\), elimination gives \((a-1)(5a-7)/385\). The roots lift to \((a,b,c)=(1,2,3)\) and \((7/5,12/7,35/11)\); the residual Jacobian determinants are respectively \(-2/385\) and \(2/385\), and the Möbius denominators are respectively \(-1/5\) and \(1/5\). Hence both branches are simple and pole-free. Appending fully bidirected children of \(0\) preserves the eliminant for every larger \(n\).
\end{lemma}
\begin{proof}
For a rooted directed tree, write \(x_v\) for the unique coefficient entering nonroot \(v\). A missing root residual is linear in one \(x_v\); a missing nonroot pair is bilinear in \(x_v,x_w\). A rank-two equation gives an invertible Möbius relation. Along each rank-two component all variables are Möbius functions of one seed \(t\), and a nonidentity cycle gives \(t=M(t)\), a polynomial of degree at most two after clearing its denominator. A rank-one residual factors into affine factors. Direct recovery by \(\Omega=(I_n-\Lambda)^{\mathsf T}\Sigma(I_n-\Lambda)\) identifies the residual quotient with the polynomial source ring, so it is a domain; therefore one factor vanishes identically and rationally anchors an endpoint. Thus a finite tree coordinate has degree at most two.

For \(n=2\), absence of the sole residual leaves the query transcendental, while its presence solves the coefficient linearly. For \(n=3\), the rooted shape is a star or chain. There are two coefficients, root residuals are linear anchors, and only one residual can involve both. Without an anchor it leaves a one-dimensional fiber; with an anchor it solves the other coefficient linearly. Unconfounded examples attain degree one. Hence \(D_0(2)=D_0(3)=1\).

For \(n=4\), take the rooted star, allow the three root--child bows, and omit all child--child bidirected pairs. At
\[
(a,b,c)=(1,2,3),\qquad
\Omega=\begin{pmatrix}
1&1/5&-1/7&1/11\\
1/5&10&0&0\\
-1/7&0&10&0\\
1/11&0&0&10
\end{pmatrix},
\]
\(\Omega\succ0\) follows from \(2|u_iu_j|\le u_i^2+u_j^2\) and strict diagonal dominance. Direct multiplication gives
\[
\Sigma=\begin{pmatrix}
1&6/5&13/7&34/11\\
6/5&57/5&79/35&203/55\\
13/7&79/35&94/7&443/77\\
34/11&203/55&443/77&215/11
\end{pmatrix}.
\]
The three residuals have lexicographic reduction
\[
5a+7b-19=0,\qquad5a-11c+28=0,\qquad(a-1)(5a-7)=0.
\]
The lifts are \((1,2,3)\) and \((7/5,12/7,35/11)\). Their Jacobian determinants are \(-2/385\) and \(2/385\), and their Möbius denominators are \(-1/5\) and \(1/5\). Thus two simple, pole-free formal branches exist, proving generic degree two. Fully bidirected root-children preserve the degree by purely transcendental extension. Hence \(D_0(n)=2\) for \(n\ge4\).
\end{proof}
\par\smallskip\noindent \textit{Justification.} Explicit residual reductions and a pole-free quadratic witness supply the size-specific lower bounds; the matching ceiling is recorded as a bibliographic proof input. \textit{Gap.} This is an appendix boundary verification, not a new tree-classification claim; the field-tier novelty remains the positive-excess all-size frontier. \textit{Consumer.} The sharp recurrence receives verified \(k=0\) base values for every graph size.

\begin{openendedquestion}[oeq:sharp-size-frontier]\label{oeq:sharp-size-frontier}
For the remaining unsettled regime \(k>0\), \(n\geq6\), and \(n<2k+4\), use \(\mathsf{FrontierHandle}\) to determine an explicit formula or finite structural recurrence solely in \(n,k\) for \(D_k(n)\).  Prove a matching universal upper bound, construct for every remaining pair an attaining graph-query family in \(\mathcal G_{n,k}\), and prove that one original directed coefficient separates the retained generic branches.  The recurrence must account for forced covariance relations, exact \(S_G=R_{\sigma}\setminus P_G\) localization, affine dehomogenization, and roots at projective infinity.  It must refine the proved coupled envelope
\[
 2^{\min\{k+1,\lfloor(n-2)/2\rfloor\}}
 \leq D_k(n)
 \leq\vartheta_{n-1,k+1}\qquad(n\geq4),
\]
agree with
\[
 D_0(2)=1,\qquad D_k(3)=1,\qquad D_k(4)=2,
\]
and
\[
 D_0(5)=2,\qquad D_1(5)=3,\qquad D_k(5)=4\quad(k\geq2),
\]
and agree with \(D_k(n)=2^{k+1}\) for \(n\geq2k+4\).  At six vertices the established boundary values and coupled upper bounds are
\[
 D_0(6)=2,\qquad D_1(6)=4,\qquad D_2(6)\leq6,
 \qquad D_3(6)\leq8,\qquad D_4(6)\leq12,
 \qquad D_k(6)\leq24\quad(k\geq5).
\]
The identity \(\vartheta_{N,g}=2^g\) for \(N\geq2g+1\) explains the same combinatorial elbow as the established SEM equality threshold, since \(N=n-1\) and \(g=k+1\); the SEM equality itself comes from the one-sink construction.  The quantity \(\vartheta_{n-1,k+1}\) is only a combinatorial ceiling and is not asserted to be SEM-attained in the unsettled regime.
\end{openendedquestion}
\par\smallskip\noindent \textit{Justification.} The remaining theorem-level kernel is exactly the positive-excess, at-least-six-vertex strip outside the large-size equality region, now constrained by the coupled orientation frontier. \textit{Gap.} The unresolved realization step must control forced covariance relations, exact localization, affine losses, and projective-infinity roots while determining how much of the joint orientation ceiling is realized by SEM incidence fibers. \textit{Consumer.} Its resolution would complete the exact ambiguity budget for every graph size supported by a positive-excess fasttreeid extension.

\begin{openendedquestion}[oeq:five-node-kill-test]\label{oeq:five-node-kill-test}
Using \(\mathsf{CensusHandle}\), complete the secondary exhaustive characteristic-zero coordinate census for every rooted graph-query orbit \((G,T,e)\in\mathcal G_{n,k}\) with \(n\le5\) and \(k\le2\). For each orbit compute \(d_e(G)\) from \(A_K=S_G^{-1}(R_{\mathrm{inc}}/I_G)\), provide a saturated rational-univariate certificate for every finite coordinate fiber, and compare the resulting degree with its prescribed-orientation coefficient, explicitly recording every strict loss caused by forced covariance relations, affine dehomogenization, or projective-infinity roots. Treat the proved five-node cubic coordinate case and the proved five-node quartic with four real positive-definite lifts as fixed base cases. The census is a regression and structural-classification problem; it is not a prerequisite for the already determined numerical frontier \(D_k(n)\) when \(n\le5\).
\end{openendedquestion}
\par\smallskip\noindent \textit{Justification.} This is the bounded destructive test that can invalidate the uncorrected orientation-realization conjecture before the all-size proof is attempted. \textit{Gap.} The five-node computations of Foygel, Draisma, and Drton classify whole graphs and selected whole-fiber degrees, not every original-edge coordinate degree under the rooted excess-arrow budget in the \(S_G\)-localized prime incidence algebra. \textit{Consumer.} The frontier proof receives certified base cases or a concrete correction term, and fasttreeid receives a complete small-instance regression suite.

\begin{lemma}[lem:quadratic-core]\label{lem:quadratic-core}
For every \((G,T,e)\in\mathcal G_{n,k}\), the construction \(\mathsf{Core}_{G,e}\) uses at most \(2k\) exposed directed coefficients, has equations of degree at most two with rational-circuit coefficients in \(K_G\), records every pole and identically zero-over-zero interaction, and preserves exactly the finite or infinite projection degree of \(\lambda_e\).
\end{lemma}
\begin{proof}
Let \(S=\{v:\operatorname{indeg}_D(v)\ge2\}\), \(s=|S|\), and expose all \(r=\sum_{v\in S}\operatorname{indeg}_D(v)\) incoming coefficients there. Since \(T\) contributes one arrow to every nonroot,
\[
k_{\mathrm{act}}=|D\setminus T|
=\sum_{v\ne\mathrm{root}}(\operatorname{indeg}_D(v)-1)\le k.
\]
Every member of \(S\) contributes at least one, so \(s\le k_{\mathrm{act}}\) and
\[
r=\sum_{v\in S}(\operatorname{indeg}_D(v)-1)+s\le2k.
\]

Every other nonroot has a scalar coefficient \(y_v\). A rank-two ordinary--ordinary residual gives a Möbius relation. A spanning forest writes \(y_v=M_v(t_C)\); each extra component edge yields a quadratic equation in \(t_C\). A rank-one residual factors into affine factors. Since \(A_K\) is a domain, one factor vanishes identically; source pullback testing selects that factor and supplies a rational anchor. Rank-zero and identically zero-over-zero cases are recorded without division.

After propagation, an ordinary--exposed equation is \(A(z)t_C+B(z)=0\), with \(A,B\) affine in one exposed block. If \(A\ne0\) in the source field, substitute \(t_C=-B/A\) and record \(A\). Two such expressions give the quadratic compatibility equation \(A_1B_2-A_2B_1=0\); a cycle equation remains quadratic after clearing the linear denominator. If \(A=0\), test and retain the affine identity \(B=0\). Exposed--exposed residuals are bilinear. Hence all core equations have degree at most two in at most \(2k\) exposed coefficients and rational-circuit coefficients in \(K_G\).

For every recorded nonzero \(g\in A_K\) and \(E=K_G(\lambda_e)\subset A_K\),
\[
(A_K[1/g])\otimes_{E,\tau}\overline{K_G}\ne0
\]
for every embedding \(\tau\), by faithful flatness. Thus localization preserves every algebraic queried conjugate; for a transcendental query it preserves a dense coordinate image. Complementary pole divisors are checked in the original incidence ideal. An autonomous ordinary component is either free or governed by one quadratic seed, and components not containing the query cannot multiply its coordinate degree. This proves exact preservation of the finite or infinite queried projection.
\end{proof}
\par\smallskip\noindent \textit{Justification.} This converts distance from a tree into a bounded-dimensional elimination problem without multiplying unrelated branch counts. \textit{Gap.} The Möbius reduction of Gupta and Bläser has one coefficient per vertex and relies on tree-specific rank-one behavior; the source-integral anchoring step is needed after extra parents couple components. \textit{Consumer.} The core is the exact symbolic intermediate representation needed to extend fasttreeid beyond one parent per node.

\begin{theorem}[thm:fpt-minimal-polynomial]\label{thm:fpt-minimal-polynomial}
For every \((G,T,e)\in\mathcal G_{n,k}\), put \(r=|D\setminus T|\). There is a universal \(C>0\) such that \(\mathsf A(G,T,e,\delta)\) returns, with total symbolic error probability at most \(\delta\), exactly one of the following four mutually exclusive outputs.
\[
\begin{array}{lll}
r=0,\ d_e(G)<\infty:
&\displaystyle
(\operatorname{tree},d_e(G),(\rho_i)_{i=1}^{d_e(G)}),
&d_e(G)\in\{1,2\},\\[3pt]
r=0,\ d_e(G)=\infty:
&(\operatorname{tree},\infty),&\\[3pt]
r>0,\ d_e(G)<\infty:
&\displaystyle
(\operatorname{algebraic},\mathcal J_{G,e},2^{\min\{|D|-n+2,\binom{n-2}{2}\}}),
&\displaystyle d_e(G)\leq2^{\min\{|D|-n+2,\binom{n-2}{2}\}}\leq2^{\min\{k+1,\binom{n-2}{2}\}},\\[3pt]
r>0,\ d_e(G)=\infty:
&(\infty,\mathcal N_{G,e}).&
\end{array}
\]
In the first case, the \(\rho_i\) are all the fractional affine square-root expressions returned by the complete tree algorithm and exhaust the generic complex queried-coordinate values. In the third case, \(\mathcal J_{G,e}\) is a Jacobian-rank transcript certifying algebraicity, but the output does not claim to compute \(d_e(G)\), \(p_e\), \(\mathcal Q_{G,e}\), or \(h_{G,e}\). In the fourth case, \(\mathcal N_{G,e}\) certifies that the queried-coordinate projection of the integral generic fiber is dominant onto \(\mathbb A^1_{K_G}\). In \(\mathsf{ArithmeticModel}\), the construction time and certificate length are at most \(n^{C(k+1)}\operatorname{poly}(\log(1/\delta))\).
\end{theorem}
\begin{proof}
Put \(r=|D\setminus T|\).  This integer is read from the input, so the cases \(r=0\) and \(r>0\) are deterministic, mutually exclusive, and exhaustive.

Suppose first that \(r=0\).  By \(\mathrm{lem:tree\mbox{-}class\mbox{-}algorithm\mbox{-}transfer}\), \(\mathcal G_{n,0}\) is exactly the rooted-tree mixed-graph class of the published complete procedure, including arbitrary bidirected support and bows.  Running that procedure with tolerance \(\delta\) returns either all one or two generic values as \(\operatorname{FASTP}\) expressions or the infinite tag.  Thus it gives exactly the first two rows of \(\mathrm{def:certificate\mbox{-}algorithm}\), with error probability at most \(\delta\).

Now suppose that \(r>0\), and enumerate the source coordinates by
\[
 u=(\lambda_D,\omega_B,\omega_{\mathrm{diag}})=(u_1,\ldots,u_N),
 \qquad N=|D|+|B|+n.
\tag{1}
\]
Then \(L=\mathbb C(u_1,\ldots,u_N)\).  By \(\mathrm{ass:acyclic}\), a topological ordering makes \(\Lambda\) strictly triangular, so \(\Lambda^n=0\) and
\[
 (I_n-\Lambda)^{-1}=I_n+\Lambda+\cdots+\Lambda^{n-1}.
\tag{2}
\]
Consequently each covariance coordinate
\[
 s_{ij}(u)=\phi_G(\sigma_{ij})
 =\bigl[(I_n-\Lambda)^{-\mathsf T}\Omega(I_n-\Lambda)^{-1}\bigr]_{ij}
\tag{3}
\]
is a polynomial of degree at most \(2n-1\), represented by an arithmetic circuit of size polynomial in \(n\).  Let \(J_0=\partial s/\partial u\), and let \(J_1\) be obtained from \(J_0\) by adjoining the row \(\partial\lambda_e/\partial u\).  Entries of \(J_0\) have degree at most \(2n-2\), and the added row is constant.

We establish the rank criterion directly.  For elements \(a_1,\ldots,a_M\in L\), choose a maximal algebraically independent subfamily \(b_1,\ldots,b_q\), and extend it by \(c_1,\ldots,c_\ell\) to a transcendence basis of \(L/\mathbb C\).  Characteristic zero makes \(L/\mathbb C(b,c)\) finite separable.  The derivations of \(\mathbb C(b,c)\) that send one \(b_j\) to \(1\) and all other basis elements to \(0\) extend uniquely through that separable extension: if \(v\) has minimal polynomial \(g(t)=\sum_h g_h t^h\), differentiating \(g(v)=0\) gives
\[
 g'(v)\,dv=-\sum_h v^h\,dg_h,
\tag{4}
\]
and \(g'(v)\ne0\).  These extended derivations show that \(db_1,\ldots,db_q\) are linearly independent over \(L\).  Conversely, each \(a_i\) is separably algebraic over \(\mathbb C(b_1,\ldots,b_q)\), and the same differentiation argument puts \(da_i\) in the span of the \(db_j\).  Therefore
\[
 \dim_L\operatorname{span}_L\{da_1,\ldots,da_M\}
 =\operatorname{trdeg}_{\mathbb C}\mathbb C(a_1,\ldots,a_M).
\tag{5}
\]
In the basis \(du_1,\ldots,du_N\) of \(\Omega_{L/\mathbb C}\), the coefficient rows of the differentials in (5) are the corresponding Jacobian rows.  Hence
\[
 \operatorname{rank}_L J_0=\operatorname{trdeg}_{\mathbb C}K_G,
 \qquad
 \operatorname{rank}_L J_1=\operatorname{trdeg}_{\mathbb C}K_G(\lambda_e).
\tag{6}
\]
Adjoining one element preserves transcendence degree precisely when that element is algebraic and otherwise raises it by one.  Thus the two generic ranks are equal if and only if \(d_e(G)<\infty\); otherwise the augmented rank is larger by one and \(d_e(G)=\infty\).

Acyclicity bounds \(|D|\leq n(n-1)/2\), while \(|B|\leq n(n-1)/2\); hence \(N\leq n^2\).  There are at most \(n(n+1)/2+1\leq n^2+1\) rows.  A nonzero minor witnessing either generic rank has degree at most
\[
 n^2(2n-2)\leq2n^3,
\tag{7}
\]
where \(n\geq2\) is \(\mathrm{ass:size}\).  Choose a finite rational sampling set \(S\) with \(|S|\geq\lceil4n^3/\delta\rceil\), and sample the source coordinates independently and uniformly from \(S\).  Specialization cannot increase rank.  If one evaluated rank is smaller than its generic rank, a fixed nonzero maximal minor vanishes.  By \(\mathrm{lem:schwartz\mbox{-}zippel\mbox{-}bound}\), this has probability at most \(\delta/2\) for each Jacobian.  A union bound therefore makes both ranks correct with probability at least \(1-\delta\).  The exact sample, pivot positions, evaluated pivot minors, and identity-test transcript constitute \(\mathcal J_{G,e}\).

On this error-free event, a rank increase means that \(\lambda_e\) is transcendental over \(K_G\).  The homomorphism
\[
 K_G[t]\longrightarrow A_K,\qquad t\longmapsto\lambda_e,
\tag{8}
\]
is then injective.  By \(\mathrm{thm:incidence\mbox{-}localization}\), \(A_K\) is an integral finitely generated \(K_G\)-algebra with fraction field \(L\).  Thus (8) induces a dominant morphism \(\operatorname{Spec}(A_K)\to\mathbb A^1_{K_G}\); moreover \(\dim A_K=\operatorname{trdeg}_{K_G}L\geq1\).  The source parametrization together with the augmented-rank transcript is exactly the dominance certificate \(\mathcal N_{G,e}\) specified in \(\mathrm{def:certificate\mbox{-}algorithm}\).  This proves the fourth output row.

If the ranks are equal, \(d_e(G)<\infty\).  By \(\mathrm{thm:orientation\mbox{-}upper\mbox{-}bound}\) and \(\mathrm{lem:tournament\mbox{-}score\mbox{-}fiber\mbox{-}ceiling}\),
\[
\begin{aligned}
 d_e(G)
 &\leq\min\{2^{|D|-n+2},\tau_{n-1}\}\\
 &\leq\min\{2^{|D|-n+2},2^{\binom{n-2}{2}}\}
 =2^{\min\{|D|-n+2,\binom{n-2}{2}\}}.
\end{aligned}
\tag{9}
\]
Since \(|T|=n-1\),
\[
 |D|-n+2=|D\setminus T|+1=r+1\leq k+1,
\tag{10}
\]
and therefore
\[
 d_e(G)\leq2^{\min\{|D|-n+2,\binom{n-2}{2}\}}
 \leq2^{\min\{k+1,\binom{n-2}{2}\}}.
\tag{11}
\]
This is exactly the third output row; neither the rank transcript nor (9) purports to compute \(d_e(G)\), \(p_e\), \(\mathcal Q_{G,e}\), or \(h_{G,e}\).

The tree procedure has randomized polynomial cost by the transfer lemma.  On the positive-excess branch, (2), matrix multiplication, circuit differentiation, exact specialization, and Gaussian elimination on matrices of dimensions at most \((n^2+1)\)-by-\(n^2\) have polynomial arithmetic cost.  Each sampled value needs \(O(\log n+\log(1/\delta))\) bits, whose separate bit cost is excluded by \(\mathsf{ArithmeticModel}\).  Since \(n\geq2\) and \(k\geq0\), one universal \(C>0\) bounds both branches by
\[
 n^{C(k+1)}\operatorname{poly}(\log(1/\delta)).
\tag{12}
\]
Only one branch is executed.  The error bounds above establish total symbolic error at most \(\delta\), and the deterministic split by \(r\), followed by the rank dichotomy (6), proves mutual exclusivity and exhaustiveness of all four outputs.
\end{proof}
\par\smallskip\noindent \textit{Justification.} The unconditional operational result is an exact algebraic-versus-transcendental tag with a certified finite-degree ceiling; it no longer launders an absent eliminant. \textit{Gap.} Neither general polynomial-space identification nor extensional quadratic-core reduction constructs the bounded source-prime chart needed for exact finite output. \textit{Consumer.} fasttreeid can safely report finite algebraic ambiguity versus an infinite queried projection; exact degree and polynomial output require the conditional chart theorem.

\begin{theorem}[thm:exceptional-set-validity]\label{thm:exceptional-set-validity}
Assume \(d_e(G)<\infty\) and a valid source-prime chart package \(\mathfrak C_{G,e}\) is supplied. With probability at least \(1-\delta\), the conditional algorithm successfully returns \((d_e(G),p_e,\mathcal Q_{G,e},h_{G,e})\); on that event, \(h_{G,e}\) is not the zero polynomial and \(\mathcal E_{G,e}\) is proper. For every complex covariance specialization \(\Sigma\notin\mathcal E_{G,e}\) in the model image, \[\{z\in\mathbb C:p_e(\Sigma,z)=0\}=\mathcal Z_e(G,\Sigma).\] This equality uses the specialized prime incidence ideal \(I_G\) and remains valid when its full \(\Lambda\)-fiber is positive dimensional. The circuits in \(\mathcal Q_{G,e}\) certify that every coordinate-preserving slice meets every generic finite branch, intermediate poles and identically zero-over-zero components are retained or lifted back to \(I_G\), the specialized leading coefficient does not drop, and no finite root is exchanged with a root at projective infinity. The roots are distinct whenever the squarefreeness circuit is nonzero.
\end{theorem}
\begin{proof}
Condition on the successful event in \(\mathrm{thm:conditional\mbox{-}geometric\mbox{-}resolution}\), whose probability is at least \(1-\delta\).  Every circuit retained in \(\mathcal Q_{G,e}\) has nonzero pullback under the covariance parametrization.  Its class in the integral domain \(R_G\cong\operatorname{im}\phi_G\) is therefore nonzero.  Their finite product is nonzero, so \(h_{G,e}\) is not the zero polynomial on the model image and
\[
 \mathcal E_{G,e}=V(h_{G,e})
\tag{43}
\]
is proper there.

The relation \(p_e(\lambda_e)=0\) holds in \(A_K\).  Clearing only covariance denominators gives \(a\in S_G\) and incidence polynomials \(c_{ij}\) such that
\[
 a(\Sigma)p_e(\Sigma,\lambda_e)
 =\sum_{i\leftrightarrow j\notin B}
 c_{ij}(\Sigma,\Lambda)F_{ij}(\Lambda,\Sigma).
\tag{44}
\]
The factor \(a\) is among the covariance-localization circuits.  If \(h_{G,e}(\Sigma)\ne0\), every solution of all specialized residuals therefore satisfies \(p_e(\Sigma,\lambda_e)=0\).  By the definition of \(\mathcal Z_e\),
\[
 \mathcal Z_e(G,\Sigma)
 \subseteq\{z\in\mathbb C:p_e(\Sigma,z)=0\}.
\tag{45}
\]

For the converse, set \(E=K_G(\lambda_e)\).  For every \(K_G\)-embedding \(\tau:E\to\overline{K_G}\), the algebra \(L\otimes_{E,\tau}\overline{K_G}\) is nonzero.  The coordinate-preserving slices in the conditional construction were chosen from a common open set meeting each conjugate component.  Their rational univariate representations consequently contain every conjugate of \(\lambda_e\).  By \(\mathrm{lem:generic\mbox{-}rur\mbox{-}specialization}\), these representations specialize as finite point sets away from loci where dimension, a denominator, a leading coefficient, a point at infinity, or separation degenerates.  Every such locus contributes a factor to \(h_{G,e}\).  The incidence-lifting identities verify by substitution that each specialized univariate point satisfies every original residual.  Hence every root of the specialized polynomial lifts to the original incidence scheme, and
\[
 \{z\in\mathbb C:p_e(\Sigma,z)=0\}
 \subseteq\mathcal Z_e(G,\Sigma).
\tag{46}
\]
The slice is used only to construct the certificate; the lifting identities return to \(I_G\), so no zero-dimensionality of the unsliced \(\Lambda\)-fiber is required.

Equations (45)--(46) prove the asserted equality.  Pole and incidence-lifting factors retain components on which an intermediate expression has a pole or is identically zero-over-zero.  Projective-leading-form and specialized-leading-coefficient factors prevent exchange with infinity and degree drop.  Finally, characteristic zero makes \(p_e\) separable over \(K_G\); outside the squarefreeness resultant factor, its specialization remains squarefree.  Thus all displayed roots are distinct.
\end{proof}
\par\smallskip\noindent \textit{Justification.} The exceptional set makes the word generic auditable and excludes poles, lost leading terms, and infinity components explicitly. \textit{Gap.} Generic RUR specialization theory supplies general bad loci, but no existing result packages \(S_G\)-localization of the prime SEM incidence ideal and single-edge projection into a causal graph certificate. \textit{Consumer.} Users can test whether an observed or simulated covariance lies inside the certified operating region before interpreting polynomial roots.

\begin{theorem}[thm:real-root-lifting]\label{thm:real-root-lifting}
Assume \(d_e(G)<\infty\), a valid source-prime chart package \(\mathfrak C_{G,e}\) is supplied, and fix a successful finite output of \(\mathrm{thm:conditional\mbox{-}geometric\mbox{-}resolution}\). There is an effectively computable finite quantifier-free Boolean combination of polynomial sign conditions with rational coefficients in \((\Sigma,z)\) which, for every real \(\Sigma\succ0\) in the covariance-model image and outside \(\mathcal E_{G,e}\) and every simple real root \(z\) of \(p_e(\Sigma,\cdot)\), holds if and only if \(z\in\mathcal R_e(G,\Sigma)\). Thus it characterizes exactly which such roots extend to complete real supported directed coefficients. If the entries of \(\Sigma\) are real algebraic numbers supplied by exact univariate representations and isolating data, an exact post-specialization procedure enumerates the simple real roots and returns a finite \(\mathsf{LiftCert}\) for each root; the certificate accepts if and only if that root belongs to \(\mathcal R_e(G,\Sigma)\). The procedure uses the original residual equations, and therefore remains valid for positive-dimensional \(\Lambda\)-fibers and for real components contained in denominator-zero loci of intermediate reductions. For every accepted complete real supported \(\Lambda\), the recovered matrix \(\Omega=(I_n-\Lambda)^{\mathsf T}\Sigma(I_n-\Lambda)\) is positive definite by congruence. No finite exact evaluation claim is made for a specialization whose real entries have no supplied finite exact representation.
\end{theorem}
\begin{proof}
Write \(s=(s_{ij}:i\leq j)\) for real symmetric covariance variables, \(z\) for the queried value, and \(\ell=(\ell_f:f\in D)\) for real directed-coefficient variables.  Define the first-order formula
\[
 \Phi_G(s,z)\quad\Longleftrightarrow\quad
 \exists\ell\in\mathbb R^D\left[
   \ell_e=z\ \wedge\
   \bigwedge_{i\leftrightarrow j\notin B}F_{ij}(\ell,s)=0
 \right].
\tag{13}
\]
Every polynomial in (13) has rational coefficients: \(D\) and \(B\) merely select variables and equations, while the displayed residual formula has coefficients in \(\{0,1,-1\}\).  By \(\mathrm{def:real\mbox{-}retained\mbox{-}set}\),
\[
 \Phi_G(\Sigma,z)\quad\Longleftrightarrow\quad
 z\in\mathcal R_e(G,\Sigma).
\tag{14}
\]
Apply the exact quantifier-elimination algorithm asserted by \(\mathrm{lem:effective\mbox{-}real\mbox{-}quantifier\mbox{-}elimination}\) over the real closed field \(\mathbb R\), with \(\ell\) quantified and \((s,z)\) free.  By the meaning of quantifier elimination, its output is a finite quantifier-free Boolean combination of polynomial equalities and strict inequalities with rational coefficients, say
\[
 \Psi_G(s,z)=\bigvee_a\bigwedge_b
 \bigl(\operatorname{sign}g_{ab}(s,z)=\epsilon_{ab}\bigr),
 \qquad
 g_{ab}\in\mathbb Q[s,z],\quad
 \epsilon_{ab}\in\{-1,0,1\},
\tag{15}
\]
and \(\Psi_G\) is equivalent to \(\Phi_G\).  This uses only the cited existence of quantifier elimination; no complexity, ordered-domain input, or sign-determination guarantee is attributed to that citation.  Equations (14)--(15) are the claimed effectively computable if-and-only-if condition.  Because (13) uses the original residual equations and makes no division, it includes positive-dimensional fibers and real components contained in denominator-zero loci of intermediate reductions.

Under the valid-chart and successful-output hypotheses, let \(\Sigma\succ0\) be real, in the covariance-model image, and outside \(\mathcal E_{G,e}\).  Then \(\mathrm{thm:exceptional\mbox{-}set\mbox{-}validity}\) gives
\[
 \{z\in\mathbb C:p_e(\Sigma,z)=0\}=\mathcal Z_e(G,\Sigma).
\tag{16}
\]
The nonvanishing squarefreeness circuit makes all roots simple.  Evaluating (15) therefore decides, for each simple real root in (16), exactly whether it extends to a complete real supported coefficient vector.  In particular, neither complex degree nor the Jacobian rank is used as a real-uniqueness assertion.

Suppose next that the entries of \(\Sigma\) are real algebraic numbers supplied by exact univariate representations and isolating data.  Every coefficient circuit of \(p_e\) is then exactly evaluable because the certified denominators are nonzero.  After multiplication by a known nonzero common denominator, \(p_e(\Sigma,t)=0\) becomes a nonzero univariate polynomial over an explicitly ordered real algebraic coefficient domain.  Apply \(\mathrm{lem:real\mbox{-}algebraic\mbox{-}sampling}\) to that univariate zero set, including among its sign polynomials the derivative and the finitely many polynomials appearing in (15).  Its real zero set is finite, and hence each semialgebraically connected component is one root.  The returned univariate representations and their isolating data consequently enumerate all simple real roots and provide the exact signs required to evaluate (15).

For an enumerated root \(z\), also form
\[
 Q_{\Sigma,z}(\ell)
 =(\ell_e-z)^2+
 \sum_{i\leftrightarrow j\notin B}F_{ij}(\ell,\Sigma)^2.
\tag{17}
\]
Over an ordered field a sum of squares is zero exactly when every summand is zero.  Thus
\[
 Z_{\mathbb R}(Q_{\Sigma,z})\ne\varnothing
 \quad\Longleftrightarrow\quad
 \Phi_G(\Sigma,z)
 \quad\Longleftrightarrow\quad
 z\in\mathcal R_e(G,\Sigma).
\tag{18}
\]
Apply the algebraic-sampling lemma again to (17), over the ordered algebraic domain generated by the supplied representations of \(\Sigma\) and \(z\).  If (17) has a real zero, the algorithm returns a real univariate representation of a complete lift \(\ell\); if it has none, the exact sign evaluation of (15) rejects.  Arithmetic in a finite tower of represented real algebraic numbers is reduced to one primitive-element representation by testing finitely many rational linear combinations and refining isolating intervals; all tests are exact polynomial remainder and sign computations.  Therefore the root representation, isolating interval, signs in (15), and, in the accepting case, the sampled lift form a finite \(\mathsf{LiftCert}\) which accepts exactly the roots in \(\mathcal R_e(G,\Sigma)\).  The theorem makes no evaluation claim when the real entries of \(\Sigma\) lack such finite exact representations.

Finally, let \(\Lambda\in\mathbb R^D\) be an accepted lift and put \(R=I_n-\Lambda\).  By \(\mathrm{ass:acyclic}\), \(R\) is unit triangular in a topological order and is invertible.  Define
\[
 \Omega=R^{\mathsf T}\Sigma R.
\tag{19}
\]
For every nonzero real vector \(x\), \(Rx\ne0\), and therefore
\[
 x^{\mathsf T}\Omega x=(Rx)^{\mathsf T}\Sigma(Rx)>0.
\tag{20}
\]
Thus \(\Omega\succ0\).  Entrywise expansion of (19) gives \(\Omega_{ij}=F_{ij}(\Lambda,\Sigma)\).  The accepting equations make this entry zero whenever \(i\leftrightarrow j\notin B\), which is exactly \(\mathrm{ass:noise\mbox{-}support}\).  Conversely, every real supported SEM lift satisfies these residual equations.  Rearranging (19) gives \(\Sigma=R^{-\mathsf T}\Omega R^{-1}\), proving both the exact real-extension characterization and the positive-definiteness conclusion by congruence.
\end{proof}
\par\smallskip\noindent \textit{Justification.} Complex coordinate degree is not a real ambiguity count, so each real polynomial root is retained only after an if-and-only-if extension test against the original residual equations. \textit{Gap.} The cited quantifier-elimination result supplies only existence of elimination over real closed fields; exact represented-root evaluation and full supported-SEM lifting are furnished by algebraic sampling and the paper's construction. \textit{Consumer.} Applied users receive all and only real supported structural explanations on the chart-certified regular domain, with positive definiteness recovered by congruence.

\begin{lemma}[lem:regular-cell-existence]\label{lem:regular-cell-existence}
Assume \(d_e(G)<\infty\), a valid source-prime chart package \(\mathfrak C_{G,e}\) is supplied, and fix a successful finite output of \(\mathrm{thm:conditional\mbox{-}geometric\mbox{-}resolution}\). Work on the resulting certificate event, on which \(h_{G,e}\) is nonzero. The real supported source set \(\{(\Lambda,\Omega):\Lambda\in\mathbb R^D,\ \Omega\succ0,\ \Omega_{ij}=0\text{ if }i\leftrightarrow j\notin B\}\) is Euclidean open in its support space and Zariski dense in the complex source space. Since \(\mathcal E_{G,e}\) is proper on the model image, there is a real supported source point \((\Lambda_0,\Omega_0)\) whose covariance \(\Sigma_0\) lies in \(\mathcal M_G\setminus\mathcal E_{G,e}\). A finite smooth-manifold stratification of the semialgebraic real-lifting refinement contains \(\Sigma_0\) in a smooth piece; a connected coordinate neighborhood \(\mathcal C_{G,e}\) inside that piece has constant real-extension status, and a smaller closed chart ball gives a nonempty compact \(\mathcal K_G\subset\mathcal C_{G,e}\) with nonempty relative interior. For every \(\Sigma\in\mathcal K_G\), \(\mathcal R_e(G,\Sigma)\ne\varnothing\), because membership in \(\mathcal M_G\) supplies at least one complete real supported \(\Lambda\).
\end{lemma}
\begin{proof}
Work throughout under the valid-chart and successful-output hypotheses in the statement.  By \(\mathrm{thm:exceptional\mbox{-}set\mbox{-}validity}\), the restriction of \(h_{G,e}\) to the covariance-model coordinate ring is nonzero.  Identify the free real entries of a supported pair \((\Lambda,\Omega)\) with \(\mathbb R^N\), where
\[
 N=|D|+|B|+n,
\]
and define
\[
 \mathcal S_+
 =\{(\Lambda,\Omega)\in\mathbb R^N:\Omega\succ0\}.
\tag{1}
\]
Here the off-diagonal coordinates of \(\Omega\) range only over \(B\), so the missing-pair restrictions are built into the ambient support space.  The set \(\mathcal S_+\) is nonempty because it contains \((0,I_n)\), and it is Euclidean open because all eigenvalues of a real symmetric matrix vary continuously.  It is Zariski dense in the corresponding complex affine source space: it contains a real open box about \((0,I_n)\), and a complex polynomial vanishing on such a box is zero, as follows by successively regarding it as a univariate polynomial in each coordinate.

For the fixed acyclic directed graph, set
\[
 \Phi_G(\Lambda,\Omega)
 =(I_n-\Lambda)^{-\mathsf T}\Omega(I_n-\Lambda)^{-1}.
\tag{2}
\]
By \(\mathrm{ass:acyclic}\), a topological ordering makes \(\Lambda\) strictly triangular, whence \(\Lambda^n=0\) and
\[
 (I_n-\Lambda)^{-1}=I_n+\Lambda+\cdots+\Lambda^{n-1}.
\tag{3}
\]
Thus \(\Phi_G\) is a polynomial map on the source coordinates.  The nonzero restriction of \(h_{G,e}\) to the model coordinate ring means precisely that
\[
 h_{G,e}\circ\Phi_G
\]
is a nonzero source polynomial.  Zariski density of \(\mathcal S_+\) therefore supplies \((\Lambda_0,\Omega_0)\in\mathcal S_+\) such that
\[
 h_{G,e}\bigl(\Phi_G(\Lambda_0,\Omega_0)\bigr)\ne0.
\tag{4}
\]
Writing \(\Sigma_0=\Phi_G(\Lambda_0,\Omega_0)\), equations (1), (2), and (4), together with \(\mathrm{def:exceptional\mbox{-}set}\), give
\[
 \Sigma_0\in\mathcal M_G\setminus\mathcal E_{G,e}.
\tag{5}
\]

We next construct the required regular cell.  The set \(\mathcal M_G\) is semialgebraic.  Indeed, it is the projection to \(\Sigma\) of the set described by the polynomial identity
\[
 (I_n-\Lambda)^{\mathsf T}\Sigma(I_n-\Lambda)=\Omega,
\tag{6}
\]
the linear equalities \(\Omega_{ij}=0\) for missing bidirected pairs, and the strict principal-minor inequalities expressing \(\Omega\succ0\).  Quantifier elimination from \(\mathrm{lem:effective\mbox{-}real\mbox{-}quantifier\mbox{-}elimination}\) turns the existential source quantifiers into an equivalent finite Boolean combination of polynomial sign conditions.  The complement of \(\mathcal E_{G,e}\) and the real-lifting conditions supplied by \(\mathrm{thm:real\mbox{-}root\mbox{-}lifting}\) are semialgebraic for the same reason.

Away from the certified leading-coefficient and squarefreeness loci, \(p_e(\Sigma,\cdot)\) has fixed degree and only simple roots.  For each possible number of real roots and each ordered real-root index, the assertion that \(z\) is the root having that index is a first-order ordered-field formula: one asserts \(p_e(\Sigma,z)=0\), counts the distinct real roots below \(z\) using finitely many quantified variables, and excludes any additional such root.  Applying the same quantifier-elimination result, then adjoining the quantifier-free lifting formula from \(\mathrm{thm:real\mbox{-}root\mbox{-}lifting}\), shows that the real-root count and every ordered root's extension label induce a finite semialgebraic partition of \(\mathcal M_G\setminus\mathcal E_{G,e}\).

Apply \(\mathrm{lem:semialgebraic\mbox{-}stratification}\) separately to the finitely many members of this partition and collect the resulting finite disjoint family of smooth manifolds.  Let \(S\) be the smooth piece containing \(\Sigma_0\).  Because \(S\) is contained in one partition member, the number of real roots and every real-extension label are constant on \(S\).

Choose a smooth coordinate chart \(u:U\to u(U)\subseteq\mathbb R^s\) on \(S\) about \(\Sigma_0\).  If \(s>0\), choose \(0<\rho<R\) with
\[
 \overline B\bigl(u(\Sigma_0),R\bigr)\subset u(U)
\]
and put
\[
 \mathcal C_{G,e}=u^{-1}\!\left(B\bigl(u(\Sigma_0),R\bigr)\right),
 \qquad
 \mathcal K_G=u^{-1}\!\left(\overline B\bigl(u(\Sigma_0),\rho\bigr)\right).
\tag{7}
\]
If \(s=0\), put \(\mathcal C_{G,e}=\mathcal K_G=\{\Sigma_0\}\).  In either case, \(\mathcal C_{G,e}\) is a connected smooth coordinate neighborhood inside one stratum, lies in \(\mathcal M_G\setminus\mathcal E_{G,e}\), and has constant real-extension labels, as required by \(\mathrm{def:regular\mbox{-}cell}\).  The set \(\mathcal K_G\) is nonempty and compact and has nonempty relative interior in \(\mathcal C_{G,e}\), as required by \(\mathrm{def:regular\mbox{-}class}\).

Finally fix \(\Sigma\in\mathcal K_G\).  Since \(\Sigma\in\mathcal M_G\), the definition of \(\mathcal M_G\) supplies real supported \(\Lambda\) and \(\Omega\succ0\) satisfying (2).  Rearranging (2) gives (6).  For each missing bidirected pair \(i\leftrightarrow j\notin B\), its \((i,j)\) entry is
\[
 \Omega_{ij}
 =\sigma_{ij}
 -\sum_{p\in\operatorname{pa}_D(i)}\lambda_{pi}\sigma_{pj}
 -\sum_{q\in\operatorname{pa}_D(j)}\lambda_{qj}\sigma_{iq}
 +\sum_{p,q}\lambda_{pi}\lambda_{qj}\sigma_{pq}
 =F_{ij}(\Lambda,\Sigma)=0,
\tag{8}
\]
where the last equality uses \(\mathrm{ass:noise\mbox{-}support}\).  Thus the real number \(\lambda_e\) has a complete real supported extension satisfying all residual equations.  By \(\mathrm{def:real\mbox{-}retained\mbox{-}set}\),
\[
 \lambda_e\in\mathcal R_e(G,\Sigma),
\]
so \(\mathcal R_e(G,\Sigma)\ne\varnothing\) for every \(\Sigma\in\mathcal K_G\).
\end{proof}
\par\smallskip\noindent \textit{Justification.} The regular statistical domain must be shown nonvacuous and its real identified set must be nonempty rather than imposed by assumption. \textit{Gap.} Standard stratification arguments isolate smooth cells, but the stated lemma ties such a cell to the certified SEM exceptional locus and to rootwise real extension. \textit{Consumer.} The inference theorem obtains a genuine compact operating region and a nonempty target set on every data-generating covariance in that region.

\begin{lemma}[lem:derived-root-separation]\label{lem:derived-root-separation}
Assume \(d_e(G)<\infty\), a valid source-prime chart package \(\mathfrak C_{G,e}\) is supplied, and fix a successful finite output of \(\mathrm{thm:conditional\mbox{-}geometric\mbox{-}resolution}\). On \(\mathcal K_G\), every nonzero denominator, specialized-leading-coefficient, and squarefreeness factor in \(h_{G,e}\) has a positive minimum in absolute value by continuity and compactness. Hence every rational-circuit coefficient of the finite monic polynomial \(p_e\) is uniformly bounded, and the monic Cauchy bound confines all \(d_e(G)\) complex roots to one fixed disk. The root-incidence set over \(\mathcal K_G\) is compact; because the squarefreeness factor does not vanish, \(\partial_t p_e\) has a positive minimum in absolute value on that set. Consequently there is \(\gamma>0\) such that every two distinct complex roots, and therefore every two distinct retained real roots, are separated by at least \(\gamma\), uniformly over \(\mathcal K_G\).
\end{lemma}
\begin{proof}
Work under the antecedents of the lemma and put \(d=d_e(G)\).  On the fixed successful finite-output event, \(\mathrm{thm:conditional\mbox{-}geometric\mbox{-}resolution}\) gives the monic polynomial
\[
 p_e(\Sigma,t)=t^d+\sum_{a=0}^{d-1}c_a(\Sigma)t^a,
\tag{9}
\]
where every \(c_a\) is a rational arithmetic circuit in the covariance entries.  By \(\mathrm{def:regular\mbox{-}class}\) and \(\mathrm{lem:regular\mbox{-}cell\mbox{-}existence}\), \(\mathcal K_G\) is compact and is contained in \(\mathcal M_G\setminus\mathcal E_{G,e}\).  Every denominator, specialized-leading-coefficient factor, and squarefreeness factor recorded in \(h_{G,e}\) is continuous and nonzero on \(\mathcal K_G\).  Each such factor consequently has a positive minimum in absolute value.  The circuit numerators have finite maxima there.  Hence every coefficient in (9) is continuous on \(\mathcal K_G\), and there is \(M<\infty\) such that
\[
 \max_{0\leq a<d}|c_a(\Sigma)|\leq M
 \qquad(\Sigma\in\mathcal K_G).
\tag{10}
\]

If \(z\) is a root of (9) and \(|z|>1+M\), then
\[
 |z|^d
 \leq M\sum_{a=0}^{d-1}|z|^a
 =M\frac{|z|^d-1}{|z|-1}
 <|z|^d,
\tag{11}
\]
a contradiction.  Therefore all roots lie in the fixed disk \(|z|\leq1+M\).  The root-incidence set
\[
 \mathcal Z_{\mathrm{root}}
 =\{(\Sigma,z)\in\mathcal K_G\times\mathbb C:p_e(\Sigma,z)=0\}
\tag{12}
\]
is closed in a compact set and hence compact.

The nonvanishing squarefreeness certificate supplied by \(\mathrm{thm:exceptional\mbox{-}set\mbox{-}validity}\) says that, for each \(\Sigma\in\mathcal K_G\), \(p_e(\Sigma,\cdot)\) and \(\partial_t p_e(\Sigma,\cdot)\) have no common root.  Thus \(\partial_t p_e\) is nonzero throughout (12).  Continuity and compactness give
\[
 \eta
 :=\min_{(\Sigma,z)\in\mathcal Z_{\mathrm{root}}}
 |\partial_t p_e(\Sigma,z)|>0.
\tag{13}
\]

Suppose, contrary to the desired conclusion, that no uniform separation exists.  Then there are \(\Sigma_\nu\in\mathcal K_G\) and distinct roots \(z_\nu,w_\nu\) of \(p_e(\Sigma_\nu,\cdot)\) such that \(|z_\nu-w_\nu|\to0\).  Compactness of (12) permits passage to a subsequence for which
\[
 \Sigma_\nu\to\Sigma_*,\qquad z_\nu\to z_*,\qquad w_\nu\to z_*.
\]
The divided difference
\[
 Q(\Sigma,z,w)
 =\frac{p_e(\Sigma,z)-p_e(\Sigma,w)}{z-w}
\tag{14}
\]
for \(z\ne w\) extends continuously, indeed polynomially in \((z,w)\) with continuous coefficient functions, to \(z=w\), where it equals \(\partial_t p_e(\Sigma,z)\).  Since \(Q(\Sigma_\nu,z_\nu,w_\nu)=0\), passage to the limit yields
\[
 p_e(\Sigma_*,z_*)=0,
 \qquad
 \partial_t p_e(\Sigma_*,z_*)=0,
\]
contradicting (13).  Consequently some \(\gamma>0\) separates every two distinct complex roots uniformly over \(\mathcal K_G\).  The retained real values in \(\mathcal R_e(G,\Sigma)\) are distinct real roots of \(p_e(\Sigma,\cdot)\) by \(\mathrm{thm:exceptional\mbox{-}set\mbox{-}validity}\) and \(\mathrm{thm:real\mbox{-}root\mbox{-}lifting}\); hence the same \(\gamma\) separates every two distinct retained roots.
\end{proof}
\par\smallskip\noindent \textit{Justification.} Denominator, derivative, and branch-separation margins follow from the certified exceptional locus and compactness, so none is imposed as a free condition. \textit{Gap.} Standard simple-root perturbation results often posit isolated roots; this lemma records the compact polynomial argument that produces the isolation margin for the certified SEM eliminant. \textit{Consumer.} The finite-set matching and simultaneous branch neighborhoods use derived regularity constants without narrowing the statistical class by an oracle margin assumption.

\begin{theorem}[thm:regular-finite-set-inference]\label{thm:regular-finite-set-inference}
Assume \(d_e(G)<\infty\), a valid source-prime chart package \(\mathfrak C_{G,e}\) is supplied, and fix a successful finite output of \(\mathrm{thm:conditional\mbox{-}geometric\mbox{-}resolution}\). There are events \(\mathcal H_m\) with \(\inf_{\Sigma\in\mathcal K_G}\Pr_{\Sigma}(\mathcal H_m)\to1\) such that, on \(\mathcal H_m\), \(\widehat{\mathcal R}_{e,m}\) is nonempty and is bijectively matched to \(\mathcal R_e(G,\Sigma)\). With the empty-set convention in \(d_H\), uniformly over \(\Sigma\in\mathcal K_G\), \(d_H(\widehat{\mathcal R}_{e,m},\mathcal R_e(G,\Sigma))=O_{\mathbb P}(m^{-1/2})\) and \(\sup_{\Sigma\in\mathcal K_G}\Pr_{\Sigma}(\widehat{\mathcal R}_{e,m}=\varnothing)\to0\). Under local branch labels,
\[
\sqrt m\{r(\pi_G(S_m))-r(\Sigma)\}\rightsquigarrow N\!\left(0,Dr(\Sigma)D\pi_G(\Sigma)V_\Sigma D\pi_G(\Sigma)^{\mathsf T}Dr(\Sigma)^{\mathsf T}\right).
\]
The simultaneous neighborhoods \(\mathcal A_m\) cover every retained branch with asymptotic probability at least \(1-\alpha\), uniformly on \(\mathcal K_G\).
\end{theorem}
\begin{proof}
Fix the valid chart and successful finite output appearing in the statement.  Let \(q=n(n+1)/2\), fix an ordering \(\operatorname{vech}\) of the symmetric matrix entries, and define
\[
 W_r=\operatorname{vech}\!\left(X^{(r)}X^{(r)\mathsf T}-\Sigma\right).
\]
Then
\[
 \sqrt m\,\operatorname{vech}(S_m-\Sigma)
 =m^{-1/2}\sum_{r=1}^m W_r.
\tag{15}
\]
By \(\mathrm{ass:iid\mbox{-}gaussian}\), the summands are independent and centered.  Differentiation of the centered Gaussian moment-generating function gives, for all indices,
\[
 \mathbb E_\Sigma(X_iX_jX_kX_l)
 =\sigma_{ij}\sigma_{kl}+\sigma_{ik}\sigma_{jl}+\sigma_{il}\sigma_{jk}.
\tag{16}
\]
Subtracting \(\mathbb E(X_iX_j)\mathbb E(X_kX_l)=\sigma_{ij}\sigma_{kl}\) therefore yields
\[
 \operatorname{Cov}_\Sigma(W_{r,ij},W_{r,kl})
 =\sigma_{ik}\sigma_{jl}+\sigma_{il}\sigma_{jk}
 =(V_\Sigma)_{ij,kl}.
\tag{17}
\]

Compactness of \(\mathcal K_G\) uniformly bounds the eigenvalues of \(\Sigma\), and hence every fixed Gaussian moment, in particular \(\sup_{\Sigma\in\mathcal K_G}\mathbb E_\Sigma\|W_r\|^3<\infty\).  For each fixed \(u\in\mathbb R^q\), Taylor's formula for \(e^{ix}\), with the remainder bounded by a constant times \(|x|^3\), gives uniformly over \(\mathcal K_G\)
\[
 \mathbb E_\Sigma\exp\!\left(\frac{i u^{\mathsf T}W_r}{\sqrt m}\right)
 =1-\frac{u^{\mathsf T}V_\Sigma u}{2m}+O(m^{-3/2}).
\tag{18}
\]
Taking the \(m\)-th power proves the Gaussian limit
\[
 \sqrt m\,\operatorname{vech}(S_m-\Sigma)
 \rightsquigarrow N(0,V_\Sigma)
\tag{19}
\]
uniformly for \(\Sigma\in\mathcal K_G\).  More explicitly, if uniform convergence failed, a sequence of offending covariance matrices would, by compactness, have a subsequence converging to some \(\Sigma_*\); equation (18) would identify every further subsequential characteristic-function limit as that of \(N(0,V_{\Sigma_*})\), while the uniformly bounded second moments give tightness, a contradiction.  The same second-moment bound and Chebyshev's inequality give
\[
 \|S_m-\Sigma\|=O_{\mathbb P}(m^{-1/2})
\tag{20}
\]
uniformly on \(\mathcal K_G\).

On the tubular neighborhood in the definition of \(\pi_G\), put
\[
 \widehat\Sigma_m=\pi_G(S_m).
\]
The nearest-point projection fixes the smooth model stratum, so \(\pi_G(\Sigma)=\Sigma\) for \(\Sigma\in\mathcal K_G\).  Smoothness on a neighborhood of the compact set and Taylor's formula imply, uniformly on \(\mathcal K_G\),
\[
 \operatorname{vech}(\widehat\Sigma_m-\Sigma)
 =D\pi_G(\Sigma)\operatorname{vech}(S_m-\Sigma)
 +O_{\mathbb P}(\|S_m-\Sigma\|^2).
\tag{21}
\]

Every coefficient circuit of \(p_e\) is smooth on a neighborhood of \(\mathcal K_G\), since each of its denominators has a positive minimum there.  By \(\mathrm{lem:derived\mbox{-}root\mbox{-}separation}\), the roots have a common separation margin and \(|\partial_t p_e|\) has a positive lower bound on the root-incidence set.  We record directly why this gives smooth root maps.  Around a root \(z_0\) at a covariance \(\Sigma_0\), choose a circle containing \(z_0\) and no other root.  Coefficient continuity keeps the circle root-free nearby, and the locally unique root is
\[
 z(\Sigma)=\frac{1}{2\pi i}\int_{|\zeta-z_0|=\varepsilon}
 \zeta\,\frac{\partial_\zeta p_e(\Sigma,\zeta)}{p_e(\Sigma,\zeta)}\,d\zeta.
\tag{22}
\]
The identity follows by factoring the polynomial and evaluating each elementary contour integral; the same formula permits differentiation under the integral and proves smoothness.  For real \(\Sigma\) and real \(z_0\), conjugation invariance and local uniqueness keep this root real.  Differentiating \(p_e(\Sigma,r_j(\Sigma))=0\) gives
\[
 Dr_j(\Sigma)[H]
 =-\frac{D_\Sigma p_e(\Sigma,r_j(\Sigma))[H]}
          {\partial_t p_e(\Sigma,r_j(\Sigma))}.
\tag{23}
\]
The increasing order of real roots, their uniform separation, and the constant extension labels in \(\mathrm{def:regular\mbox{-}cell}\) patch these local maps into the common retained-root labeling \(r=(r_1,\ldots,r_J)\).  A finite cover of \(\mathcal K_G\) in (22), together with the denominator and derivative margins, uniformly bounds the first two derivatives of all \(r_j\).

Combining (19)--(23) and applying Taylor's formula to \(r\) yields
\[
 \sqrt m\{r(\pi_G(S_m))-r(\Sigma)\}
 \rightsquigarrow N\!\left(0,
 Dr(\Sigma)D\pi_G(\Sigma)V_\Sigma
 D\pi_G(\Sigma)^{\mathsf T}Dr(\Sigma)^{\mathsf T}
 \right)
\tag{24}
\]
uniformly on \(\mathcal K_G\).  Indeed, the Taylor remainders are uniformly \(o_{\mathbb P}(1)\) after multiplication by \(\sqrt m\), by (20), while the same subsequence argument used after (18) proves uniformity of the transformed limit.

All certificate factors are continuous and nonzero on \(\mathcal K_G\), and the extension labels are locally constant on the selected cell.  Compactness therefore supplies a relative open neighborhood \(U\) of \(\mathcal K_G\) within the regular cell such that the projection is defined, every certificate factor remains valid, and precisely the same \(J\) branches are retained.  Let \(\gamma\) be the separation margin from \(\mathrm{lem:derived\mbox{-}root\mbox{-}separation}\), and define \(\mathcal H_m\) to be the event that
\[
 \begin{split}
 &S_m\text{ lies in the projection tube},\qquad \widehat\Sigma_m\in U,\\
 &\max_{1\leq j\leq J}|r_j(\widehat\Sigma_m)-r_j(\Sigma)|<\gamma/3,
 \quad\text{and every conditional certificate is valid at }\widehat\Sigma_m.
 \end{split}
\tag{25}
\]
Equations (20), (21), uniform continuity of the finitely many root maps, and the positive certificate margins give
\[
 \inf_{\Sigma\in\mathcal K_G}\Pr_\Sigma(\mathcal H_m)\longrightarrow1.
\tag{26}
\]
On \(\mathcal H_m\), \(\mathrm{thm:real\mbox{-}root\mbox{-}lifting}\) and \(\mathrm{def:finite\mbox{-}set\mbox{-}estimator}\) give
\[
 \widehat{\mathcal R}_{e,m}
 =\{r_1(\widehat\Sigma_m),\ldots,r_J(\widehat\Sigma_m)\}.
\tag{27}
\]
The margin in (25) makes (27) a bijective matching to
\(\mathcal R_e(G,\Sigma)=\{r_1(\Sigma),\ldots,r_J(\Sigma)\}\).  By \(\mathrm{lem:regular\mbox{-}cell\mbox{-}existence}\), \(J\geq1\), so the estimated set is nonempty on \(\mathcal H_m\).  Moreover, on that event,
\[
 d_H\!\left(\widehat{\mathcal R}_{e,m},\mathcal R_e(G,\Sigma)\right)
 \leq\max_j|r_j(\widehat\Sigma_m)-r_j(\Sigma)|
 =O_{\mathbb P}(m^{-1/2})
\tag{28}
\]
uniformly.  The empty-set convention causes no difficulty off \(\mathcal H_m\): for any \(\varepsilon>0\), first use (26) to make the probability of \(\mathcal H_m^c\) less than \(\varepsilon/2\) uniformly, and then choose the constant in the uniform bound in (28) so that the remaining exceedance probability is below \(\varepsilon/2\).  Hence the asserted uniform \(O_{\mathbb P}(m^{-1/2})\) statement holds even though the distance may equal infinity on the exceptional event.  Equation (26) also gives
\[
 \sup_{\Sigma\in\mathcal K_G}
 \Pr_\Sigma(\widehat{\mathcal R}_{e,m}=\varnothing)
 \leq
 \sup_{\Sigma\in\mathcal K_G}\Pr_\Sigma(\mathcal H_m^c)
 \longrightarrow0.
\tag{29}
\]

It remains to construct the simultaneous neighborhoods.  Let \(s=\dim\mathcal C_{G,e}\) and set
\[
 W_\pi(\Sigma)=D\pi_G(\Sigma)V_\Sigma D\pi_G(\Sigma)^{\mathsf T}.
\tag{30}
\]
The covariance matrix \(V_\Sigma\) is positive definite on symmetric-coordinate directions.  Indeed, every nonzero linear functional of \(\operatorname{vech}(XX^{\mathsf T})\) can be written as \(X^{\mathsf T}HX\) for a nonzero symmetric \(H\), and (16) gives
\[
 \operatorname{Var}_\Sigma(X^{\mathsf T}HX)
 =2\operatorname{tr}\!\left((\Sigma^{1/2}H\Sigma^{1/2})^2\right)>0,
\tag{31}
\]
where \(\Sigma\succ0\) because \(\mathcal K_G\subset\mathcal M_G\).  The derivative of the nearest-point projection is the identity on the tangent space of its stratum and has that \(s\)-dimensional tangent space as its image.  Thus \(W_\pi(\Sigma)\) has constant rank \(s\).  Its smallest positive eigenvalue is bounded away from zero on \(\mathcal K_G\) by continuity and compactness, so its Moore--Penrose inverse varies continuously there.

Define the plug-in matrix
\[
 \widehat W_{\pi,m}
 =D\pi_G(\widehat\Sigma_m)V_{\widehat\Sigma_m}
  D\pi_G(\widehat\Sigma_m)^{\mathsf T}.
\]
on the event \(\widehat\Sigma_m\in U\), and assign it any fixed symmetric value otherwise.  The latter assignment occurs with probability tending to zero uniformly by (26).
Equations (19)--(21), continuity, and diagonalization on the rank-\(s\) tangent space imply
\[
 m\,\operatorname{vech}(\widehat\Sigma_m-\Sigma)^{\mathsf T}
 \widehat W_{\pi,m}^{+}
 \operatorname{vech}(\widehat\Sigma_m-\Sigma)
 \rightsquigarrow\chi_s^2
\tag{32}
\]
uniformly on \(\mathcal K_G\).  To verify the displayed law, if \(Z\sim N(0,W_\pi(\Sigma))\), diagonalization shows that \(Z^{\mathsf T}W_\pi(\Sigma)^+Z\) is the sum of the squares of \(s\) independent standard normal variables; uniform plug-in validity follows from the positive eigenvalue bound and the subsequence argument used for (19).

Let \(c_{s,1-\alpha}\) be the \((1-\alpha)\)-quantile of \(\chi_s^2\), and define the model-stratum ellipsoid
\[
 E_m=\left\{\Theta\in U:
 m\,\operatorname{vech}(\widehat\Sigma_m-\Theta)^{\mathsf T}
 \widehat W_{\pi,m}^{+}
 \operatorname{vech}(\widehat\Sigma_m-\Theta)
 \leq c_{s,1-\alpha}\right\}.
\tag{33}
\]
Equation (32) yields
\[
 \liminf_{m\to\infty}\inf_{\Sigma\in\mathcal K_G}
 \Pr_\Sigma(\Sigma\in E_m)\geq1-\alpha.
\tag{34}
\]
If \(s=0\), take the degenerate ellipsoid at the locally constant projected point; its limiting coverage is one.  Finally set
\[
 \mathcal A_{m,j}=r_j(E_m),\qquad 1\leq j\leq J,
\]
as prescribed by \(\mathrm{def:branch\mbox{-}neighborhoods}\).  Whenever \(\Sigma\in E_m\), one has \(r_j(\Sigma)\in\mathcal A_{m,j}\) simultaneously for every \(j\).  Consequently (34) proves
\[
 \liminf_{m\to\infty}\inf_{\Sigma\in\mathcal K_G}
 \Pr_\Sigma\!\left(r_j(\Sigma)\in\mathcal A_{m,j}
 \text{ for every }1\leq j\leq J\right)
 \geq1-\alpha,
\]
which covers all retained branches and selects none.
\end{proof}
\par\smallskip\noindent \textit{Justification.} The observationally equivalent branches are the estimand, so inference must cover them simultaneously rather than select one. \textit{Gap.} Chernozhukov, Hong, and Tamer treat general criterion-defined sets, while Mareis, Sturma, and Drton derive regular half-trek point inference; neither gives root-specific joint inference for a finite algebraic SEM coordinate fiber. \textit{Consumer.} Researchers using a graph with added parents can report sampling uncertainty around every compatible direct-effect branch.

\begin{proposition}[prop:singleton-inference]\label{prop:singleton-inference}
Assume \(d_e(G)<\infty\), a valid source-prime chart package \(\mathfrak C_{G,e}\) is supplied, and fix a successful finite output of \(\mathrm{thm:conditional\mbox{-}geometric\mbox{-}resolution}\). If \(\mathcal R_e(G,\Sigma)=\{r_1(\Sigma)\}\) on \(\mathcal K_G\), then with probability tending to one uniformly on \(\mathcal K_G\), \(\widehat{\mathcal R}_{e,m}\) has exactly one element; that element is consistent and asymptotically normal with the scalar variance given by the corresponding diagonal entry of the covariance matrix in \(\mathrm{thm:regular\mbox{-}finite\mbox{-}set\mbox{-}inference}\).
\end{proposition}
\begin{proof}
Assume \(\mathcal R_e(G,\Sigma)=\{r_1(\Sigma)\}\) throughout \(\mathcal K_G\).  Define \(\widehat r_m\) to be the unique element of \(\widehat{\mathcal R}_{e,m}\) when that set is a singleton and define \(\widehat r_m=0\) otherwise.  On the events \(\mathcal H_m\) supplied by \(\mathrm{thm:regular\mbox{-}finite\mbox{-}set\mbox{-}inference}\), the estimated set is bijectively matched to the population singleton.  Hence
\[
 \inf_{\Sigma\in\mathcal K_G}
 \Pr_\Sigma\!\left(\widehat{\mathcal R}_{e,m}=\{\widehat r_m\}\right)
 \longrightarrow1.
\tag{35}
\]
The complement of the event in (35) has probability tending to zero uniformly.  On the event in (35), the Hausdorff distance between the two singleton sets equals absolute error, so the rate conclusion of \(\mathrm{thm:regular\mbox{-}finite\mbox{-}set\mbox{-}inference}\), together with this uniformly negligible complement, gives
\[
 |\widehat r_m-r_1(\Sigma)|=O_{\mathbb P}(m^{-1/2})
\tag{36}
\]
uniformly on \(\mathcal K_G\), which proves consistency.  Taking the sole coordinate in the joint limit of that theorem gives
\[
 \sqrt m\{\widehat r_m-r_1(\Sigma)\}
 \rightsquigarrow N(0,W_{11}(\Sigma)),
\tag{37}
\]
where
\[
 W(\Sigma)
 =Dr(\Sigma)D\pi_G(\Sigma)V_\Sigma
 D\pi_G(\Sigma)^{\mathsf T}Dr(\Sigma)^{\mathsf T}.
\tag{38}
\]
The uniformly negligible off-event definition does not alter (37).  If \(W_{11}(\Sigma)=0\), the right side of (37) is the degenerate centered normal law; no unasserted variance lower bound is needed.
\end{proof}
\par\smallskip\noindent \textit{Justification.} The singleton case connects finite-set inference to the familiar direct-effect estimator and Wald limit. \textit{Gap.} Regular half-trek asymptotic normality covers rationally identified coefficients, whereas this corollary also covers a singleton real projection of a higher complex-degree coordinate polynomial. \textit{Consumer.} When only one real supported branch survives, applied users obtain an ordinary point estimate, standard error, and confidence interval.

\begin{lemma}[lem:multihomogeneous-bezout-isolated-bound]\label{lem:multihomogeneous-bezout-isolated-bound}
Let \(f_1,\ldots,f_q\) be a square polynomial system whose variables are partitioned into blocks of sizes \(r_1,\ldots,r_b\), where \(\sum_i r_i=q\), and suppose that the degree of \(f_\ell\) in block \(i\) is at most \(\alpha_{\ell i}\). If the affine common-zero set of the system is finite, then its number of distinct points is at most the coefficient of \(\prod_i X_i^{r_i}\) in
\[
 \prod_{\ell=1}^{q}\left(\sum_{i=1}^{b}\alpha_{\ell i}X_i\right).
\]
\end{lemma}

\begin{lemma}[lem:finite-nonsingular-perturbation-bound]\label{lem:finite-nonsingular-perturbation-bound}
Let \(f_1,\ldots,f_q\in\mathbb C[x]\) be a square polynomial system, with the variables partitioned into blocks of sizes \(r_1,\ldots,r_b\), where \(\sum_i r_i=q\), and with block-degree bounds \(\alpha_{\ell i}\in\mathbb Z_{\geq0}\) and \(\deg_i(f_\ell)\leq\alpha_{\ell i}\). If \(S\) is a finite set of distinct affine common zeros at each of which the full \(q\)-by-\(q\) Jacobian is nonsingular, then
\[
 |S|\leq
 [X_1^{r_1}\cdots X_b^{r_b}]
 \prod_{\ell=1}^{q}\left(\sum_{i=1}^{b}\alpha_{\ell i}X_i\right).
\]
This conclusion does not require the original common-zero set to be zero dimensional.
\end{lemma}
\begin{proof}
Put
\[
 C=[X_1^{r_1}\cdots X_b^{r_b}]
 \prod_{\ell=1}^{q}\left(\sum_{i=1}^{b}\alpha_{\ell i}X_i\right).
\tag{16}
\]
The assertion is immediate if \(S=\varnothing\), so suppose that \(S\ne\varnothing\), fix \(z\in S\), and write \(J_f(z)\) for the full Jacobian. Since \(\det J_f(z)\ne0\), its determinant expansion contains a nonzero term. Thus there is a permutation \(\pi\) of the \(q\) scalar variables such that
\[
 \prod_{\ell=1}^{q}
 \frac{\partial f_\ell}{\partial x_{\pi(\ell)}}(z)\ne0.
\tag{17}
\]
If \(x_{\pi(\ell)}\) belongs to block \(i(\ell)\), then (17) implies \(\alpha_{\ell,i(\ell)}\geq1\). The permutation uses exactly \(r_i\) variables from block \(i\). Hence the choices \(X_{i(\ell)}\) in the product in (16) contribute a strictly positive summand to its indicated coefficient, and therefore
\[
 C>0.
\tag{18}
\]

Homogenize each \(f_\ell\) separately in every block to multidegree \(\alpha_\ell=(\alpha_{\ell1},\ldots,\alpha_{\ell b})\) on
\[
 X=\prod_{i=1}^{b}\mathbb P^{r_i}_{\mathbb C}.
\tag{19}
\]
Let \(W_\ell\) be the vector space of multihomogeneous forms of multidegree \(\alpha_\ell\), and put \(\mathcal P=\prod_{\ell=1}^{q}\mathbb P(W_\ell)\). Every \(\alpha_\ell\) has a positive entry: otherwise the \(\ell\)-th Jacobian row in (17) would vanish. Each \(W_\ell\) is therefore base-point free on \(X\). Consider the universal incidence set
\[
 \mathcal Z=
 \{(x,[g_1],\ldots,[g_q])\in X\times\mathcal P:
       g_\ell(x)=0\text{ for every }\ell\}.
\tag{20}
\]
For fixed \(x\), each equation \(g_\ell(x)=0\) is one nontrivial linear condition on \(\mathbb P(W_\ell)\). Since \(\dim X=q\), it follows directly from (20) that
\[
 \dim\mathcal Z
 =q+\sum_{\ell=1}^{q}(\dim\mathbb P(W_\ell)-1)
 =\dim\mathcal P.
\tag{21}
\]
The projection \(\mathcal Z\to\mathcal P\) is projective. Upper semicontinuity of fiber dimension for a projective morphism makes the locus of fibers of dimension at least one Zariski closed. Equation (21) shows that the generic fiber has dimension zero when the projection is dominant and is empty otherwise. Hence there is a nonempty Zariski-open set \(\mathcal U\subseteq\mathcal P\) on which every fiber is finite or empty.

The variety \(\mathcal P\) is irreducible. Consequently the proper algebraic set \(\mathcal P\setminus\mathcal U\) has empty interior in the complex topology, so systems in \(\mathcal U\) occur arbitrarily close to the homogenized tuple \((f_1,\ldots,f_q)\) after choosing local affine coefficient charts in the projective spaces \(\mathbb P(W_\ell)\). Choose pairwise disjoint analytic coordinate balls \(B_z\), one about each \(z\in S\), all contained in the standard affine chart used before homogenization. The holomorphic implicit-function theorem and \(\det J_f(z)\ne0\) give a coefficient neighborhood \(\mathcal N\) of \((f_1,\ldots,f_q)\) such that every system in \(\mathcal N\) has one nonsingular common zero in each \(B_z\). Shrinking the balls and \(\mathcal N\) if necessary makes these continued zeros distinct. Choose
\[
 (g_1,\ldots,g_q)\in\mathcal U\cap\mathcal N.
\tag{22}
\]
The block multidegrees in (22) are the same prescribed \(\alpha_\ell\), its projective common-zero scheme is finite, and it has at least \(|S|\) distinct nonsingular zeros in the original affine chart. Projective-infinity zeros, if present, are also finite and do not remove any of these affine continuations.

The affine common-zero set of (22) is therefore finite. Applying \(\mathrm{lem:multihomogeneous\mbox{-}bezout\mbox{-}isolated\mbox{-}bound}\) to that finite system bounds its number of distinct affine common zeros by \(C\). The continued zeros are distinct, so \(|S|\leq C\), as claimed. Notice that the cited finite-system theorem is invoked only after (22) has made the whole projective system, and hence its affine chart, finite.
\end{proof}

\begin{lemma}[lem:generic-rur-specialization]\label{lem:generic-rur-specialization}
\texttt{For a generically zero-\allowbreak{}dimensional polynomial system with parameter variables,\allowbreak{} a generic rational univariate representation based on a generically separating linear form specializes to a rational univariate representation away from the finite union of the algebraic loci where dimension,\allowbreak{} a leading coefficient,\allowbreak{} a point at infinity,\allowbreak{} or separation degenerates.}
\end{lemma}

\begin{lemma}[lem:semialgebraic-stratification]\label{lem:semialgebraic-stratification}
\texttt{Every real semialgebraic set can be written as a disjoint union of finitely many smooth manifolds.}
\end{lemma}

\begin{lemma}[lem:schwartz-zippel-bound]\label{lem:schwartz-zippel-bound}
If a nonzero polynomial has total degree at most \(d\), evaluation at independent uniform points of a finite set \(S\) vanishes with probability at most \(d/|S|\).
\end{lemma}

\begin{lemma}[lem:effective-real-quantifier-elimination]\label{lem:effective-real-quantifier-elimination}
For every real closed field \(R\), there exists an algorithm for performing quantifier elimination from first-order formulas over \(R\).
\end{lemma}

\begin{lemma}[lem:real-algebraic-sampling]\label{lem:real-algebraic-sampling}
Let \(R\) be a real closed field, let \(D_0\subseteq R\) be an ordered domain, let \(Q\in D_0[X_1,\ldots,X_s]\), and let \(\mathcal P\subset D_0[X_1,\ldots,X_s]\) be finite. There is an exact algebraic-sampling algorithm whose output consists of finitely many real univariate representations whose associated points meet every semialgebraically connected component of every nonempty realizable sign condition of \(\mathcal P\) on \(\{x\in R^s:Q(x)=0\}\).
\end{lemma}

\begin{lemma}[lem:groebner-free-geometric-resolution]\label{lem:groebner-free-geometric-resolution}
\texttt{Over a field of characteristic zero,\allowbreak{} a polynomial system given by straight-\allowbreak{}line programs and satisfying the reduced regular-\allowbreak{}sequence hypothesis away from a supplied inequation admits a probabilistic Gröbner-\allowbreak{}free geometric-\allowbreak{}resolution algorithm. On successful generic choices,\allowbreak{} the algorithm represents the relevant zero-\allowbreak{}dimensional section by a primitive element,\allowbreak{} its monic polynomial,\allowbreak{} and rational coordinate parametrizations,\allowbreak{} with arithmetic complexity polynomial in the straight-\allowbreak{}line-\allowbreak{}program evaluation length and the geometric degrees of the successive varieties.}
\end{lemma}

\begin{lemma}[lem:lifting-fiber-equidimensional-decomposition]\label{lem:lifting-fiber-equidimensional-decomposition}
Over a field of characteristic zero, there is a probabilistic lifting-fiber algorithm which, for arbitrary polynomial equations \(f_1=\cdots=f_s=0\) and an inequation \(g\neq0\) given by a straight-line program, computes the equidimensional decomposition of the Zariski closure of their solution set. On successful generic choices the components are encoded by lifting fibers, and the arithmetic complexity is polynomial in the straight-line-program length and the intermediate geometric-degree parameter appearing in the theorem.
\end{lemma}

\begin{theorem}[thm:conditional-geometric-resolution]\label{thm:conditional-geometric-resolution}
Let \((G,T,e)\in\mathcal G_{n,k}\) satisfy \(d_e(G)<\infty\), and suppose a valid source-prime chart package \(\mathfrak C_{G,e}\) is supplied for every source-consistent recorded stratum. There is a universal \(a>0\) such that a randomized localized geometric-resolution algorithm returns exactly \((d_e(G),p_e,\mathcal Q_{G,e},h_{G,e})\) with symbolic error probability at most \(\delta\). Its arithmetic construction time and certificate length are at most
\[
(n+s)^{a(k+1)}\mathfrak d^a\operatorname{poly}(\log(1/\delta)),
\]
where \(s=s(\mathfrak C_{G,e})\) and \(\mathfrak d=\mathfrak d(\mathfrak C_{G,e})\). In particular, for any fixed \(b>0\), if \(s\leq n^b\) and \(\mathfrak d\leq n^{b(k+1)}\), then the bound is \(n^{C(k+1)}\operatorname{poly}(\log(1/\delta))\) for a universal \(C\) depending only on \(a\) and \(b\). The output polynomial is the monic minimal polynomial over \(K_G\), not a whole-fiber eliminant, and \(\mathcal Q_{G,e}\) contains the covariance-localization, incidence-lifting, pole, projective-infinity, specialized-leading-coefficient, squarefreeness, and real-lifting circuits whose nonzero product is \(h_{G,e}\). The theorem is conditional on the supplied charts; it does not construct or select them.
\end{theorem}
\begin{proof}
Fix one source-consistent recorded stratum.  By \(\mathrm{lem:quadratic\mbox{-}core}\) and \(\mathrm{def:source\mbox{-}prime\mbox{-}chart}\), its exposed tuple is \(x=(x_1,\ldots,x_r)\) with \(r\leq2k\), and \(P=\ker\xi\subset K_G[x]\).  Chart validity gives
\[
 P K_G[x]_{H\Delta}
 =(h_1,\ldots,h_c)K_G[x]_{H\Delta},
 \qquad \xi(H\Delta)\ne0,
\tag{51}
\]
where \(c=\operatorname{ht}P\), the localized sequence is reduced and regular, and \(\Delta\) is a nonzero \(c\)-by-\(c\) Jacobian minor.  Hence (51) presents the localization of the source prime, rather than an extraneous residual component.  It is smooth and pure dimensional of dimension
\[
 q=r-c\leq2k.
\tag{52}
\]
The principal open set is represented polynomially by adjoining a variable \(u\) and the equation \(uH\Delta-1=0\).

Let \(E=K_G(\lambda_e)\).  The hypothesis \(d_e(G)<\infty\) and characteristic zero make \(E/K_G\) finite separable.  For each \(K_G\)-embedding \(\tau:E\to\overline{K_G}\), the algebra
\[
 L\otimes_{E,\tau}\overline{K_G}
\tag{53}
\]
is nonzero: after identifying \(E\) with \(\tau(E)\), this follows because extension by the field \(\overline{K_G}\) is faithfully flat.  Thus the base-changed source-prime chart has a nonempty geometric component above every conjugate of \(\lambda_e\), and \(H\Delta\) remains nonzero on the corresponding principal open component.

The generic slicing operation in \(\mathrm{lem:groebner\mbox{-}free\mbox{-}geometric\mbox{-}resolution}\), applied to the reduced regular sequence (51), supplies \(q\) affine linear forms whose successive intersections are transverse, zero dimensional, and nonempty on every one of the finitely many conjugate components.  If \(q=0\), no slice is added.  Generic choices avoiding the finitely many transversality and leading-form polynomials work simultaneously for all components.  The cited algorithm then returns, on successful choices, a finite reduced algebra \(B_{\mathfrak C}\), multiplication matrices, and a primitive-element representation for the slice.  The chart's rational lifting circuits recover every eliminated directed coordinate, including \(\lambda_e\), and the chart identities make all original residuals vanish in the source-prime localization.

Let \(M_{e,\mathfrak C}\) be multiplication by the recovered \(\lambda_e\) on \(B_{\mathfrak C}\), and put
\[
 \chi_{e,\mathfrak C}(t)=\det(tI-M_{e,\mathfrak C}).
\tag{54}
\]
Since \(p_e(\lambda_e)=0\) in the source field, every root occurring in (54) is a root of \(p_e\).  Conversely, (53) supplies a component above every root of \(p_e\), and the simultaneous slice meets every such component; hence every root of \(p_e\) occurs in at least one polynomial (54).  Repetitions may arise from multiple slice points or charts, but characteristic zero makes \(p_e\) separable.  It follows that
\[
 p_e(t)=\operatorname{sqfree}\!\left(
   \prod_{\mathfrak C}\chi_{e,\mathfrak C}(t)
 \right),
 \qquad \deg p_e=d_e(G).
\tag{55}
\]
Indeed, the right side is monic and has exactly the distinct \(K_G\)-conjugates of \(\lambda_e\) as its roots.  Thus (55) recovers the single-coordinate minimal polynomial even when the unsliced parameter fiber is positive dimensional.

All chart coefficients and lifting maps are rational arithmetic circuits in the covariance entries.  Determinants, multiplication-matrix operations, differentiation, subresultant greatest-common-divisor computations, squarefree extraction, and monic normalization are arithmetic-circuit operations over \(K_G\); hence (55) gives rational circuits for all coefficients of \(p_e\).  For each lifting identity, clear covariance denominators and record both those denominators and the resulting original-incidence identity.  Homogenization before slicing records the projective leading forms; univariate normalization records the specialized leading coefficient and squarefreeness resultant.  Covariance-localization factors are retained separately.  The rational-univariate coordinate formulas, their denominators, their isolating-data input polynomials, and the cleared original-residual verification identities form the real-lifting circuits; quantifier elimination is performed only later, in \(\mathrm{thm:real\mbox{-}root\mbox{-}lifting}\), and is not charged to this construction.  These finitely many certified nonzero circuits form \(\mathcal Q_{G,e}\), and their repetition-free cleared product is \(h_{G,e}\).  Equation (51) and \(\mathrm{thm:incidence\mbox{-}localization}\) ensure that no directed-coordinate expression is inverted, so every recorded pole or identically zero-over-zero stratum is either represented by a supplied valid chart or checked after lifting against the original incidence equations.

For arbitrary recorded equations and inequations use \(\mathrm{lem:lifting\mbox{-}fiber\mbox{-}equidimensional\mbox{-}decomposition}\); for regular localized slices use \(\mathrm{lem:groebner\mbox{-}free\mbox{-}geometric\mbox{-}resolution}\).  Their stated complexity bounds, (52), multiplication-matrix linear algebra, and construction of the finite lifting and verification circuits give a universal \(a>0\) for which the total arithmetic work and output length are at most
\[
 (n+s)^{a(k+1)}\mathfrak d^a
 \operatorname{poly}(\log(1/\delta)),
\tag{56}
\]
where \(s=s(\mathfrak C_{G,e})\) and \(\mathfrak d=\mathfrak d(\mathfrak C_{G,e})\).  By definition, \(\mathfrak d\) charges every successive intersection, localization, slice, and lifting fiber, so no created branch is omitted from (56).

Every randomized decision is a nonvanishing test for an explicit chart, transversality, leading-form, lifting, or resultant polynomial.  If their degree bounds are \(d_1,\ldots,d_M\), sample independently from finite sets \(S_j\) satisfying \(d_j/|S_j|\leq\delta/M\).  By \(\mathrm{lem:schwartz\mbox{-}zippel\mbox{-}bound}\) and a union bound,
\[
 \Pr(\text{some symbolic genericity decision is wrong})
 \leq\sum_{j=1}^M\frac{d_j}{|S_j|}\leq\delta.
\tag{57}
\]
On the complementary event, (51)--(55) show that the returned tuple is exactly \((d_e(G),p_e,\mathcal Q_{G,e},h_{G,e})\).

Finally, if \(s\leq n^b\) and \(\mathfrak d\leq n^{b(k+1)}\), then \(n\geq2\) implies, after enlarging a universal exponent,
\[
 (n+s)^{a(k+1)}\mathfrak d^a
 \leq n^{a(\max\{1,b\}+1+b)(k+1)}
 \leq n^{C(k+1)}.
\tag{58}
\]
This proves the asserted fixed-parameter specialization.  Every step began with supplied valid localized source-prime charts; no step constructs those charts or selects the embedded source prime without them.
\end{proof}
\par\smallskip\noindent \textit{Justification.} A supplied valid localized reduced source-prime chart is sufficient for exact coordinate minimal-polynomial and specialization-certificate construction. \textit{Gap.} General geometric-resolution algorithms neither select the embedded SEM source prime nor remove their dependence on supplied chart size and intermediate geometric degree. \textit{Consumer.} This is the exact higher-degree polynomial-output contract for a chart-certified extension of fasttreeid; it is not an unconditional positive-excess eliminator.

\begin{openendedquestion}[oeq:unconditional-finite-output]\label{oeq:unconditional-finite-output}
Does there exist, in \(\mathsf{ArithmeticModel}\), a chart-free randomized algorithm which, for every \((G,T,e)\in\mathcal G_{n,k}\) with \(d_e(G)<\infty\), constructs the exact tuple \((d_e(G),p_e,\mathcal Q_{G,e},h_{G,e})\) with time and certificate length \(n^{C(k+1)}\operatorname{poly}(\log(1/\delta))\) and symbolic error probability at most \(\delta\)? Constructing valid source-prime chart packages with \(s\leq n^b\) and \(\mathfrak d\leq n^{b(k+1)}\) would suffice by \(\mathrm{thm:conditional\mbox{-}geometric\mbox{-}resolution}\), but no equivalence with bounded chart construction is claimed.
\end{openendedquestion}
\par\smallskip\noindent \textit{Justification.} The unconditional finite-output guarantee is the remaining algorithmic commitment after the valid algebraicity tag is separated from chart-dependent elimination. \textit{Gap.} The quadratic core does not provide a reduced local presentation, a source-nonzero Jacobian minor, source-component selection, or bounded descent over the implicitly presented covariance field. \textit{Consumer.} Resolving this question positively would remove the chart input from the exact higher-degree fasttreeid workflow.

\begin{lemma}[lem:gupta-blaser-tree-classification]\label{lem:gupta-blaser-tree-classification}
For a linear recursive structural equation model whose directed support is a rooted spanning tree and whose bidirected support is arbitrary, including bows, the randomized polynomial-time procedure of Gupta and Bl\"aser decides for every queried directed coefficient \(\lambda_e\) whether the generic queried-coordinate set has one value, two values, or infinitely many values. In each finite case it returns all corresponding \(\operatorname{FASTP}\) expressions for the queried coefficient. For every \(\eta\in(0,1)\), the polynomial-identity-test sampling sets can be chosen so that the total error probability is at most \(\eta\).
\end{lemma}

\begin{lemma}[lem:tree-class-algorithm-transfer]\label{lem:tree-class-algorithm-transfer}
For every \(n\geq2\), the class \(\mathcal G_{n,0}\) is exactly the rooted tree-shaped mixed-graph class with arbitrary bidirected support treated by the complete tree algorithm, up to the source paper's harmless vertex-count indexing convention. Consequently, on every input \((G,T,e)\in\mathcal G_{n,0}\), that algorithm returns the exact value \(d_e(G)\in\{1,2,\infty\}\) and, when \(d_e(G)<\infty\), all one or two \(\operatorname{FASTP}\) expressions for \(\lambda_e\).
\end{lemma}
\begin{proof}
Let \((G,T,e)\in\mathcal G_{n,0}\). By \(\mathrm{def:graph\mbox{-}class}\), one has \(T\subseteq D\) and
\[
 |D\setminus T|\leq0.
\]
Cardinality is nonnegative, so \(D\setminus T=\varnothing\), whence \(D=T\). By \(\mathrm{ass:arborescence}\), this common edge set is a rooted spanning arborescence. Thus the entire directed support is a rooted spanning tree. The definition of \(\mathcal G_{n,0}\) imposes no restriction on \(B\), so its bidirected support is arbitrary and may contain bows.

Conversely, given any rooted tree-shaped mixed graph on \(n\) vertices and any queried directed edge \(e\), put \(T=D\). Then \(T\) is a rooted spanning arborescence and \(|D\setminus T|=0\), so \((G,T,e)\in\mathcal G_{n,0}\). These two inclusions prove equality of the graph-query classes. A convention that counts nonroot vertices rather than all vertices merely shifts the displayed size index and changes no graph.

Apply \(\mathrm{lem:gupta\mbox{-}blaser\mbox{-}tree\mbox{-}classification}\) to this common class. Suppose first that the published procedure returns \(m\in\{1,2\}\) generic values. By \(\mathrm{thm:incidence\mbox{-}localization}\), if \(\lambda_e\) is algebraic over \(K_G\), its generic affine coordinate image over \(\overline{K_G}\) is exactly the root set of its minimal polynomial; if it is transcendental, that image is Zariski dense in \(\mathbb A^1_{K_G}\) and hence infinite. The returned finite set therefore implies algebraicity. If the minimal polynomial \(p\) had a repeated root, then \(\gcd(p,p')\neq1\). Irreducibility would force \(p\mid p'\), which is impossible because \(\deg p'<\deg p\) and \(p'\neq0\) in characteristic zero. Thus the minimal polynomial is separable, so its number of distinct roots equals its degree. Hence
\[
 d_e(G)=[K_G(\lambda_e):K_G]=m.
\]
If the published procedure instead returns infinitely many generic values, algebraicity would contradict the finite root-set conclusion of the same incidence-localization theorem. Thus \(\lambda_e\) is transcendental over \(K_G\) and \(d_e(G)=\infty\). The returned finite formulas are exactly the one or two \(\operatorname{FASTP}\) expressions in the cited result. The cited procedure is randomized polynomial time and permits its polynomial-identity-test error to be set to any prescribed \(\eta\in(0,1)\). This is a transfer on the equal zero-excess class and makes no assertion for \(|D\setminus T|>0\).
\end{proof}

\begin{lemma}[lem:small-block-orientation-maxima]\label{lem:small-block-orientation-maxima}
In the multihomogeneous system used by \(\mathrm{thm:orientation\mbox{-}upper\mbox{-}bound}\), the query component has at most \(n-1\) nonempty child blocks and its binary-factor graph is simple. For a simple graph on at most two blocks every prescribed-indegree coefficient is at most \(1\); on at most three blocks it is at most \(2\); and on at most four blocks it is at most \(4\), with
\[
 [X_1X_2X_3^2X_4^2]\prod_{1\leq i<j\leq4}(X_i+X_j)=4.
\]
If \(n=5\) and \(|D\setminus T|\leq1\), the coefficient is at most \(3\). Hence every finite queried-coordinate degree satisfies
\[
 d_e(G)\leq1\quad(n=3),\qquad
 d_e(G)\leq2\quad(n=4),\qquad
 d_e(G)\leq3\quad(n=5,\ |D\setminus T|\leq1),
\]
and \(d_e(G)\leq4\) when \(n=5\).
\end{lemma}
\begin{proof}
For a finite simple graph \(H=(W,E_H)\) and \(r\in\mathbb Z_{\ge0}^W\), put
\[
 c_H(r)=\left[\prod_{v\in W}X_v^{r_v}\right]
 \prod_{\{u,v\}\in E_H}(X_u+X_v).
\tag{80}
\]
Choosing \(X_v\) from the factor for \(\{u,v\}\) is equivalent to orienting that edge toward \(v\).  Thus \(c_H(r)\) counts orientations with indegree vector \(r\), and
\[
 c_H(r)=0\quad\text{unless}\quad
 \sum_{v\in W}r_v=|E_H|,qquad
 0\le r_v\le\deg_H(v).
\tag{81}
\]

On at most two vertices, simplicity leaves at most one edge and its direction is fixed by \(r\), so the count is at most one.  On at most three vertices, a forest has at most one orientation with prescribed indegrees: the orientation of an edge at a leaf is forced by the leaf's required indegree, and induction deletes that leaf.  For a triangle, the two directed cycles both have score \((1,1,1)\), while every transitive orientation is uniquely determined by its labeled permutation of \((0,1,2)\).  Hence the maximum on three vertices is two.

Now suppose \(|W|\le4\).  With at most four edges, every connected component is a tree except possibly one unicyclic component.  Repeated leaf deletion forces all edges off the unique cycle, and at most the two cyclic orientations remain.  Therefore
\[
 c_H(r)\le2\qquad(|W|\le4,\ |E_H|\le4).
\tag{82}
\]
With five edges, \(H\) is \(K_4\) minus one edge, say \(\{1,2\}\).  Conditional on the required indegrees \(a\) and \(b\) of vertices \(1\) and \(2\), direct collection in
\[
 (X_1+X_3)(X_1+X_4)(X_2+X_3)(X_2+X_4)(X_3+X_4)
\]
gives the maximum coefficient over the remaining two exponents as
\[
 \begin{array}{c|ccc}
  &b=0&b=1&b=2\\ \hline
 a=0&1&2&1\\
 a=1&2&3&2\\
 a=2&1&2&1.
 \end{array}
\tag{83}
\]
To verify the table, choose the \(a\) edges into vertex \(1\) and the \(b\) edges into vertex \(2\) among their two incident edges; the final edge \(\{3,4\}\) adds one to exactly one remaining indegree.  Grouping the \(\binom2a\binom2b\) choices by the two remaining indegrees gives (83).  Thus every five-edge coefficient is at most three.

With six edges, \(H=K_4\).  Up to permutation, the score vectors allowed by (81) are
\[
 (0,0,3,3),\ (0,1,2,3),\ (0,2,2,2),\
 (1,1,1,3),\ (1,1,2,2).
\tag{84}
\]
The first is impossible because the edge between the two zero-score vertices has no allowed head.  Type \((0,1,2,3)\) forces the unique transitive tournament.  For \((0,2,2,2)\) and \((1,1,1,3)\), the extreme vertex's incident edges are forced and the remaining triangle has its two cyclic orientations.  For \((1,1,2,2)\), condition on the edge between the two score-one vertices.  Whichever of those two vertices receives it must send its two remaining edges outward, and the residual triangle has two compatible choices.  The two directions of the conditioned edge therefore give
\[
 [X_1X_2X_3^2X_4^2]
 \prod_{1\le i<j\le4}(X_i+X_j)=2+2=4.
\tag{85}
\]
This proves that four is the maximum on four blocks.

In the square subsystem selected in \(\mathrm{thm:orientation\mbox{-}upper\mbox{-}bound}\), variables are partitioned by child.  The root has no incoming arrow by acyclicity, so at most \(n-1\) blocks are nonempty.  Each residual has degree at most one in either incident child block, and an unordered child pair supplies at most one residual; hence its binary-factor graph is simple.  The finite-selected-zero bound in \(\mathrm{thm:orientation\mbox{-}upper\mbox{-}bound}\) and the preceding counts yield
\[
 d_e(G)\le1\quad(n=3),\qquad
 d_e(G)\le2\quad(n=4),\qquad
 d_e(G)\le4\quad(n=5).
\tag{86}
\]
If \(n=5\) and \(|D\setminus T|\le1\), then \(|D|\le5\), so the selected square system has at most five equations and at most five binary factors after unary factors are removed.  Four or fewer binary factors give at most two by (82).  Five force four nonempty blocks, graph \(K_4\setminus e\), no unary factor, and positive block sizes summing to five, hence an exponent permutation of \((2,1,1,1)\).  Table (83) bounds that coefficient by three.  Therefore \(d_e(G)\le3\) in this last regime, proving all assertions.
\end{proof}

\begin{lemma}[lem:foygel-five-node-degree-three]\label{lem:foygel-five-node-degree-three}
\texttt{Example 8 and Figure 5(b) of Foygel,\allowbreak{} Draisma,\allowbreak{} and Drton exhibit an acyclic five-\allowbreak{}vertex linear structural equation model whose directed support consists of five arrows and contains a rooted spanning arborescence,\allowbreak{} and whose generic covariance parametrization has degree three. Equivalently,\allowbreak{} its source function field is a degree-\allowbreak{}three extension of its covariance function field.}
\end{lemma}

\begin{lemma}[lem:five-node-cubic-coordinate-transfer]\label{lem:five-node-cubic-coordinate-transfer}
There is a graph-query \((G,T,e)\in\mathcal G_{5,1}\) for which \(d_e(G)=3\). In particular,
\[
 D_1(5)\geq3,
\]
and one original directed coefficient separates the three characteristic-zero conjugate branches.
\end{lemma}
\begin{proof}
Let \(\widetilde G\) be the graph in \(\mathrm{lem:foygel\mbox{-}five\mbox{-}node\mbox{-}degree\mbox{-}three}\). Its directed support has five arrows and contains a rooted spanning arborescence \(T\), which has four arrows on five vertices. Hence
\[
 |D\setminus T|=5-4=1,
\]
and \((\widetilde G,T,f)\in\mathcal G_{5,1}\) for every directed edge \(f\) of \(\widetilde G\).

Write \(K_{\widetilde G}\subset L_{\widetilde G}\) for the covariance and source function fields. The cited degree statement gives
\[
 [L_{\widetilde G}:K_{\widetilde G}]=3.
 \tag{5}
\]
For fixed \(\Sigma\) and \(\Lambda\), the noise matrix is not an additional fiber choice, because the covariance equation rearranges to
\[
 \Omega=(I_5-\Lambda)^{\mathsf T}\Sigma(I_5-\Lambda).
 \tag{6}
\]
If every directed coordinate \(\lambda_f\) belonged to \(K_{\widetilde G}\), then (6) would put every allowed noise coordinate in \(K_{\widetilde G}\) as well. Since the source field is generated by the directed and allowed noise coordinates, this would give \(L_{\widetilde G}=K_{\widetilde G}\), contradicting (5). Thus some original directed edge \(e\) satisfies \(\lambda_e\notin K_{\widetilde G}\).

The tower law and (5) give
\[
 [K_{\widetilde G}(\lambda_e):K_{\widetilde G}]
 \mid [L_{\widetilde G}:K_{\widetilde G}]=3.
\]
The left side is greater than one, so primality of three forces it to equal three. Therefore \(d_e(\widetilde G)=3\), proving the lower bound. If its minimal polynomial \(p\) were inseparable, irreducibility would give \(p\mid p'\); this is impossible because \(p'\neq0\) in characteristic zero and \(\deg p'<\deg p\). By \(\mathrm{thm:incidence\mbox{-}localization}\), the three distinct roots are exactly the generic affine queried-coordinate image, so this original edge coordinate separates all three conjugate branches.
\end{proof}

\begin{lemma}[lem:five-node-quartic-witness]\label{lem:five-node-quartic-witness}
Let \(G^{(4)}=(V,D,B)\) have \(V=\{0,1,2,3,4\}\),
\[
 D=\{0\to1,0\to2,1\to3,2\to3,1\to4,2\to4\},
 \qquad
 B=\{0\leftrightarrow i:1\leq i\leq4\},
\]
let
\[
 T=\{0\to1,0\to2,1\to3,1\to4\},
 \qquad e=2\to4,
\]
and write \(f=\lambda_e\). Then \((G^{(4)},T,e)\in\mathcal G_{5,2}\) and \(d_e(G^{(4)})=4\). More precisely, at the rational positive-definite covariance \(\Sigma^{\star}\) displayed in the proof, the residual fiber has four reduced points with distinct real \(f\)-coordinates, namely the roots of
\[
 q(f)=(f-7)\bigl(8599643757f^3-181150021425f^2
 +1271943861499f-2976933819575\bigr).
\]
Every root extends to real directed coefficients and a positive-definite disturbance matrix \(\Omega\) having the required missing-bidirected zero pattern. For every \(N\geq5\), appending \(N-5\) fully bidirected children of \(0\) and placing their arrows in \(T\) preserves the queried-coordinate degree four. Consequently \(D_k(5)\geq4\) for every admissible \(k\geq2\).
\end{lemma}
\begin{proof}
Write
\[
 (a,b,c,d,u,f)=(\lambda_{01},\lambda_{02},\lambda_{13},
 \lambda_{23},\lambda_{14},\lambda_{24}).
\]
The displayed \(T\) has four arrows, gives each nonroot one parent, and reaches every vertex from \(0\); it is a rooted spanning arborescence.  Moreover \(D\setminus T=\{2\to3,2\to4\}\), the vertex order is topological, and \(e=2\to4\).  Hence \((G^{(4)},T,e)\in\mathcal G_{5,2}\).

At the rational source point
\[
 (a,b,c,d,u,f)=(2,3,4,5,6,7),
 \qquad
 \Omega^{\star}=
 \begin{pmatrix}
 30&1&2&3&4\\
 1&31&0&0&0\\
 2&0&32&0&0\\
 3&0&0&33&0\\
 4&0&0&0&34
 \end{pmatrix},
\tag{87}
\]
all nonzero off-diagonal entries are allowed by \(B\).  Using \(2i|y_0y_i|\le i(y_0^2+y_i^2)\) for \(1\le i\le4\),
\[
 y^{\mathsf T}\Omega^{\star} y
 \ge20y_0^2+30\sum_{i=1}^4y_i^2>0
\]
for \(y\ne0\), so \(\Omega^{\star}\succ0\).  Direct multiplication of the covariance parametrization gives
\[
 \Sigma^{\star}=
 \begin{pmatrix}
 30&61&92&707&1014\\
 61&155&187&1561&2247\\
 92&187&314&2327&3332\\
 707&1561&2327&17981&25747\\
 1014&2247&3332&25747&36972
 \end{pmatrix}.
\tag{88}
\]
It is positive definite by congruence because \(I_5-\Lambda^{\star}\) is unit triangular.

The missing bidirected pairs are the six pairs inside \(\{1,2,3,4\}\).  Their residuals at (88) are
\[
\begin{aligned}
 R_1={}&30ab-92a-61b+187,\\
 R_2={}&61ac+92ad-707a-155c-187d+1561,\\
 R_3={}&61au+92af-1014a-155u-187f+2247,\\
 R_4={}&61bc+92bd-707b-187c-314d+2327,\\
 R_5={}&61bu+92bf-1014b-187u-314f+3332,\\
 R_6={}&155cu+187cf-2247c+187du+314df-3332d\\
 &{}-1561u-2327f+25747.
\end{aligned}
\tag{89}
\]
Exact rational Buchberger reduction in lexicographic order \(a>b>c>d>u>f\), with nonzero rational normalizations suppressed, gives
\[
\begin{aligned}
 H_1={}&384915840a+425450175590061f^3-8965411803407607f^2\\
 &{}+62974399432228195f-147445028655847225,\\
 H_2={}&741827520b-593375419233f^3+12321625507347f^2\\
 &{}-85271405126447f+196665729339485,\\
 H_3={}&615040c+25798931271f^3-543666188277f^2\\
 &{}+3818873142057f-8941504654939,\\
 H_4={}&4d-3f+1,\\
 H_5={}&153760u+8599643757f^3-181222062759f^2\\
 &{}+1272957714019f-2980501654153,\\
 H_6={}&q(f)=(f-7)c(f),\\
 c(f)={}&8599643757f^3-181150021425f^2
 +1271943861499f-2976933819575.
\end{aligned}
\tag{90}
\]
This is a two-sided characteristic-zero ideal certificate: each \(H_j\) is a rational combination obtained by reductions from the \(R_i\), and division of every \(R_i\) by the displayed list has zero remainder.  Conversely, after monic normalization the leading monomials are \(a,b,c,d,u,f^4\), which are pairwise coprime, so Buchberger's product criterion reduces every remaining critical pair to zero.  Therefore
\[
 (R_1,\ldots,R_6)=(H_1,\ldots,H_6).
\tag{91}
\]
Exact integer arithmetic gives
\[
 c(7)=-30256,
 \qquad
 \operatorname{disc}(c)
 =1802165059445823134111101206528>0.
\tag{92}
\]
A real cubic with roots \(r,s+it,s-it\), where \(t\ne0\), has discriminant \(-4A^4t^2((r-s)^2+t^2)^2<0\).  Hence (92) implies that \(c\) has three distinct real roots, and \(c(7)\ne0\) makes them distinct from \(7\).  Thus \(q\) has four distinct real roots.  The first five equations in (90) uniquely recover real \((a,b,c,d,u)\) from each root.  Consequently
\[
 \mathbb C[a,b,c,d,u,f]/(R_1,\ldots,R_6)
 \cong\mathbb C[f]/(q)\cong\mathbb C^4.
\tag{93}
\]

The residual Jacobian has rank six at every point: the triangular Jacobian has five nonzero constant pivots and final derivative \(q'(f)\ne0\), and differentiation of the two-sided identities (91) equates its row span with that of the original residual Jacobian.  Perturb the fifteen covariance entries independently as \(\Sigma^{\star}+z\).  The invertible Jacobian recursively solves each homogeneous degree of a formal directed-coordinate series, producing four formal branches with distinct constant \(f\)-terms.

The source dimension is
\[
 |D|+|B|+|V|=6+4+5=15,
\tag{94}
\]
equal to covariance dimension.  A covariance polynomial in \(P_G\) vanishes after substitution of a formal section and hence after the injective translation \(\Sigma^{\star}+z\); it is therefore zero.  Thus \(P_G=(0)\), \(K_G=\mathbb C(\sigma_{ij}:i\le j)\), and equal transcendence dimensions make \(L/K_G\) finite.  Moreover, each formal section maps the integral incidence ring into \(\mathbb C[[z]]\) and its image contains the fifteen algebraically independent translated covariance coordinates.  The incidence ring has dimension fifteen by (94) and \(\mathrm{thm:incidence\mbox{-}localization}\), so the section's prime kernel has height zero and is therefore zero.  Each section consequently extends to an embedding of \(L\) over the common embedded copy of \(K_G\).

The four formal branches give four distinct images of \(f\) under embeddings over the common covariance field, so
\[
 4\leq[K_G(f):K_G]=d_e(G^{(4)}).
\tag{95}
\]
The matching five-vertex ceiling \(d_e(G^{(4)})\leq4\) is \(\mathrm{lem:small\mbox{-}block\mbox{-}orientation\mbox{-}maxima}\), whose upper bound now factors through the finite-selected-zero perturbation argument in \(\mathrm{thm:orientation\mbox{-}upper\mbox{-}bound}\). Hence \(d_e(G^{(4)})=4\).

For each root, let \(\widetilde\Lambda\) be its real triangular recovery and put
\[
 \widetilde\Omega=(I_5-\widetilde\Lambda)^{\mathsf T}
 \Sigma^{\star}(I_5-\widetilde\Lambda).
\]
Equations (89) are exactly its six required zero entries.  Since \(I_5-\widetilde\Lambda\) is invertible and \(\Sigma^{\star}\succ0\), congruence gives \(\widetilde\Omega\succ0\).  Thus all four roots have real supported positive-definite lifts.

Finally append a child \(w\) of \(0\), place \(0\to w\) in \(T\), and allow all bidirected pairs incident with \(w\).  No new missing residual involves \(w\).  Congruence is an invertible triangular polynomial change between the new allowed disturbance row and the new covariance row \(s_w\), so
\[
 L_1=L_0(\lambda_{0w},s_w),\qquad K_1=K_0(s_w),
\]
with these added coordinates algebraically independent in the indicated order.  The old minimal polynomial remains irreducible under the purely transcendental extension \(K_0(s_w)/K_0\), preserving degree four.  Iteration proves the appending assertion.  Since \(\mathcal G_{5,2}\subseteq\mathcal G_{5,k}\) for every admissible \(k\ge2\), the definition of \(D_k(5)\) gives \(D_k(5)\ge4\).
\end{proof}

\begin{theorem}[thm:small-size-frontier]\label{thm:small-size-frontier}
For every admissible pair with \(2\leq n\leq5\), the exact sharp coordinate-degree frontier is
\[
\begin{array}{c|c}
 n & D_k(n)\\ \hline
 2 & D_0(2)=1,\\
 3 & D_k(3)=1\quad(0\leq k\leq1),\\
 4 & D_k(4)=2\quad(0\leq k\leq3),\\
 5 & D_0(5)=2,\quad D_1(5)=3,\quad D_k(5)=4\quad(2\leq k\leq6).
\end{array}
\]
In each nontrivial lower-bound construction an original directed coefficient separates all generic conjugate branches.
\end{theorem}
\begin{proof}
By \(\mathrm{ass:k\mbox{-}domain}\), the admissible budgets at a fixed size \(n\) are the integers \(0\leq k\leq (n-1)(n-2)/2\).  We first record a monotonicity consequence of \(\mathrm{def:graph\mbox{-}class}\).  If \(0\leq j\leq k\), then every triple satisfying \(|D\setminus T|\leq j\) also satisfies \(|D\setminus T|\leq k\), while all other membership conditions are unchanged.  Hence
\[
 \mathcal G_{n,j}\subseteq\mathcal G_{n,k}.
\tag{1}
\]
Because \(D_k(n)\) is the maximum of the finite queried-coordinate degrees over \(\mathcal G_{n,k}\), (1) implies
\[
 D_j(n)\leq D_k(n)
\tag{2}
\]
whenever both budgets are admissible.

For \(n=2\), the budget condition gives \(k=0\), and \(\mathrm{lem:tree\mbox{-}maxima\mbox{-}support}\) gives
\[
 D_0(2)=1.
\tag{3}
\]
For \(n=3\), the admissible budgets are \(k=0,1\).  The two-block conclusion of \(\mathrm{lem:small\mbox{-}block\mbox{-}orientation\mbox{-}maxima}\) gives \(d_e(G)\leq1\) for every finite query in either class, and therefore \(D_k(3)\leq1\).  On the other hand, \(\mathrm{lem:tree\mbox{-}maxima\mbox{-}support}\) gives \(D_0(3)=1\); (2) then gives \(D_k(3)\geq1\).  Thus
\[
 D_k(3)=1\qquad(0\leq k\leq1).
\tag{4}
\]

For \(n=4\), the admissible budgets are \(0\leq k\leq3\).  The three-block conclusion of \(\mathrm{lem:small\mbox{-}block\mbox{-}orientation\mbox{-}maxima}\) gives \(D_k(4)\leq2\).  The tree equality \(D_0(4)=2\) from \(\mathrm{lem:tree\mbox{-}maxima\mbox{-}support}\), followed by (2), gives \(D_k(4)\geq2\).  Consequently
\[
 D_k(4)=2\qquad(0\leq k\leq3).
\tag{5}
\]

For \(n=5\), the admissible budgets are \(0\leq k\leq6\).  The tree equality in \(\mathrm{lem:tree\mbox{-}maxima\mbox{-}support}\) gives
\[
 D_0(5)=2.
\tag{6}
\]
When \(k=1\), the refined five-vertex conclusion of \(\mathrm{lem:small\mbox{-}block\mbox{-}orientation\mbox{-}maxima}\) gives \(D_1(5)\leq3\).  The graph-query furnished by \(\mathrm{lem:five\mbox{-}node\mbox{-}cubic\mbox{-}coordinate\mbox{-}transfer}\) belongs to \(\mathcal G_{5,1}\) and has a queried original coefficient of degree three, so the definition of \(D_1(5)\) gives the reverse inequality.  Therefore
\[
 D_1(5)=3.
\tag{7}
\]
For every admissible \(k\geq2\), the four-block conclusion of \(\mathrm{lem:small\mbox{-}block\mbox{-}orientation\mbox{-}maxima}\) gives \(D_k(5)\leq4\).  By \(\mathrm{lem:five\mbox{-}node\mbox{-}quartic\mbox{-}witness}\), there is a graph-query in \(\mathcal G_{5,2}\) whose queried original coefficient has degree four.  Inclusion (1) with \(j=2\) places the same query in \(\mathcal G_{5,k}\), whence \(D_k(5)\geq4\).  Thus
\[
 D_k(5)=4\qquad(2\leq k\leq6).
\tag{8}
\]
Equations (3)--(8) are precisely the displayed table.  Finally, the degree-two tree construction in \(\mathrm{lem:tree\mbox{-}maxima\mbox{-}support}\), the cubic construction in \(\mathrm{lem:five\mbox{-}node\mbox{-}cubic\mbox{-}coordinate\mbox{-}transfer}\), and the quartic construction in \(\mathrm{lem:five\mbox{-}node\mbox{-}quartic\mbox{-}witness}\) each use an original directed coefficient and assert that its distinct generic conjugate values separate the corresponding coordinate branches.  These are exactly the nontrivial lower-bound constructions used above.
\end{proof}
\par\smallskip\noindent \textit{Justification.} The exact table removes every vertex size through five from the unresolved frontier. \textit{Gap.} The remaining orientation-realization issue begins at six vertices with positive excess budget in the small-size strip. \textit{Consumer.} The frontier open question and the software ambiguity budget use these exact base cases.

\begin{proposition}[prop:frontier-lower-envelope]\label{prop:frontier-lower-envelope}
For every admissible \(k\) and every \(n\geq4\),
\[
 D_k(n)\geq
 2^{\min\{k+1,\lfloor(n-2)/2\rfloor\}}.
\]
The lower bound is attained within the one-sink family by one original directed coefficient. At \(n=5\), the separate exact values are the sharper conclusions \(D_1(5)=3\) and \(D_k(5)=4\) for \(k\geq2\).
\end{proposition}
\begin{proof}
Fix an admissible pair \((n,k)\) with \(n\geq4\), and define the auxiliary integer
\[
 j:=\min\!\left\{k,\left\lfloor\frac{n-4}{2}\right\rfloor\right\}.
\tag{9}
\]
Since \(k\geq0\) by \(\mathrm{ass:k\mbox{-}domain}\) and \(n\geq4\), we have \(j\geq0\).  Moreover, \(j\leq\lfloor(n-4)/2\rfloor\), so
\[
 n\geq2j+4.
\tag{10}
\]
Thus \(\mathrm{thm:large\mbox{-}size\mbox{-}frontier}\), applied with excess budget \(j\), gives the one-sink graph-query \((G^{\star}_{j,n},T,e^{\star})\) and the equality
\[
 d_{e^{\star}}(G^{\star}_{j,n})=2^{j+1}.
\tag{11}
\]
Its membership in \(\mathcal G_{n,j}\) is also the explicit conclusion of \(\mathrm{def:attaining\mbox{-}membership\mbox{-}check}\).  Since \(j\leq k\), \(\mathrm{def:graph\mbox{-}class}\) gives
\[
 \mathcal G_{n,j}\subseteq\mathcal G_{n,k}:
\tag{12}
\]
indeed, the sole budget inequality strengthens from \(|D\setminus T|\leq j\) to \(|D\setminus T|\leq k\), and the remaining membership conditions do not change.  The finite degree in (11) is therefore among the values maximized in the definition of \(D_k(n)\).  Equations (11)--(12) imply
\[
 D_k(n)\geq2^{j+1}.
\tag{13}
\]

It remains only to rewrite the exponent.  For every integer \(n\geq4\),
\[
 \left\lfloor\frac{n-4}{2}\right\rfloor+1
 =\left\lfloor\frac{n-2}{2}\right\rfloor.
\tag{14}
\]
For nonnegative integers \(u,v\), one also has \(\min\{u,v\}+1=\min\{u+1,v+1\}\).  Applying this identity to (9) and then using (14) yields
\[
 j+1
 =\min\!\left\{k+1,\left\lfloor\frac{n-2}{2}\right\rfloor\right\}.
\tag{15}
\]
Substitution of (15) into (13) proves
\[
 D_k(n)\geq
 2^{\min\{k+1,\lfloor(n-2)/2\rfloor\}}.
\]
The witness is the same original directed coefficient \(e^{\star}\) specified by \(\mathrm{def:attaining\mbox{-}family}\) and certified by \(\mathrm{thm:large\mbox{-}size\mbox{-}frontier}\).  Finally, \(\mathrm{thm:small\mbox{-}size\mbox{-}frontier}\) gives the separate exact five-vertex conclusions \(D_1(5)=3\) and \(D_k(5)=4\) for every admissible \(k\geq2\), as claimed.
\end{proof}
\par\smallskip\noindent \textit{Justification.} Class monotonicity turns the large-size one-sink construction into a certified lower bound at every size. \textit{Gap.} The envelope need not be sharp in the remaining strip, as the sharper five-vertex cubic and quartic values demonstrate. \textit{Consumer.} Any eventual all-size recurrence must dominate this explicit lower envelope.

\begin{lemma}[lem:tournament-score-fiber-ceiling]\label{lem:tournament-score-fiber-ceiling}
Let \(N\geq1\). If \(H\) is a simple graph on \(b\leq N\) labeled vertices and \(r\in\mathbb Z_{\geq0}^{b}\), then the number of orientations of \(H\) having indegree vector \(r\) is at most \(\tau_N\). Moreover,
\[
 \tau_N\leq 2^{\binom{N-1}{2}},
 \qquad
 \tau_1=1,
 \quad \tau_2=1,
 \quad \tau_3=2,
 \quad \tau_4=4,
 \quad \tau_5=24<64.
\]
\end{lemma}
\begin{proof}
Write \(H'\) for the graph on vertex set \([N]\) obtained from \(H\) by adjoining \(N-b\) isolated vertices and relabeling if necessary. Fix once and for all an orientation \(\eta\) of every edge of \(K_N\setminus H'\). If \(\xi\) is an orientation of \(H'\) whose restriction to the nonisolated vertices has indegree vector \(r\), complete \(\xi\) to a tournament \(\overline\xi\) by using \(\eta\) on \(K_N\setminus H'\). Let \(s_i\) be the indegree contributed at vertex \(i\) by \(\eta\). Every completed tournament then has the same score vector
\[
 a_i=
 \begin{cases}
 r_i+s_i,&1\leq i\leq b,\\
 s_i,&b<i\leq N.
 \end{cases}
\]
The completion map is injective because restricting \(\overline\xi\) to \(E(H')\) recovers \(\xi\). By \(\mathrm{def:tournament\mbox{-}score\mbox{-}fiber}\), the number of such completions is \(c_N(a)\leq\tau_N\). This proves the first assertion, including disconnected \(H\) and isolated dummy vertices.

For the general numerical ceiling, fix a spanning tree of \(K_N\) when \(N\geq2\). After the \(\binom N2-(N-1)=\binom{N-1}{2}\) non-tree edges have been oriented, at most one orientation of the tree can realize a prescribed score vector. Indeed, the required residual indegree at a leaf forces its only remaining tree edge; deleting that leaf and iterating determines every tree edge. Hence
\[
 c_N(a)\leq 2^{\binom{N-1}{2}}
\]
for every \(a\). The same bound is immediate for \(N=1\), so it also holds after maximizing over \(a\).

It remains to verify the stated small values without an unsupported enumeration. Separating the factors incident to vertex \(b\) gives the finite recurrence
\[
 c_b(a_1,\ldots,a_b)
 =
 \sum_{\substack{S\subseteq[b-1]\\ |S|=a_b}}
 c_{b-1}\!\left(
 a_1-\mathbf 1\{1\notin S\},\ldots,
 a_{b-1}-\mathbf 1\{b-1\notin S\}
 \right),
\]
where \(c_1(0)=1\) and a term is zero if any argument lies outside \(\{0,\ldots,b-2\}\). Here \(S\) is the set of edges incident to \(b\) oriented toward \(b\). The polynomial defining \(c_b\) is symmetric, and a nonzero coefficient must satisfy
\[
 0\leq a_i\leq b-1,
 \qquad
 \sum_{i=1}^{b}a_i=\binom b2.
\]
Thus it suffices to apply the displayed recurrence to the nondecreasing integer vectors satisfying these two elementary restrictions. For \(b=1,2,3\), it gives respectively
\[
\begin{array}{c|c}
b&\text{nonzero sorted score vector and coefficient}\\
\hline
1&(0):1\\
2&(0,1):1\\
3&(0,1,2):1,\quad(1,1,1):2.
\end{array}
\]
For \(b=4\), all sorted candidates and their recursively computed coefficients are
\[
\begin{array}{c|ccccc}
a&(0,0,3,3)&(0,1,2,3)&(0,2,2,2)&(1,1,1,3)&(1,1,2,2)\\
\hline
c_4(a)&0&1&2&2&4.
\end{array}
\]
For \(b=5\), the corresponding exhaustive table is
\[
\begin{array}{c|rrrrrrrrrrrr}
a
&(0,0,2,4,4)
&(0,0,3,3,4)
&(0,1,1,4,4)
&(0,1,2,3,4)
&(0,1,3,3,3)
&(0,2,2,2,4)
&(0,2,2,3,3)
&(1,1,1,3,4)
&(1,1,2,2,4)
&(1,1,2,3,3)
&(1,2,2,2,3)
&(2,2,2,2,2)\\
\hline
c_5(a)&0&0&0&1&2&2&4&2&4&8&14&24.
\end{array}
\]
Every entry is obtained from the preceding \(c_4\) row by the displayed subset sum. For example, for the maximizing vector each of the \(\binom42=6\) subsets contributes \(c_4(1,1,2,2)=4\), and therefore
\[
 c_5(2,2,2,2,2)=6\cdot4=24.
\]
Taking the maximum in each exhaustive row proves \(\tau_1=1\), \(\tau_2=1\), \(\tau_3=2\), \(\tau_4=4\), and \(\tau_5=24<64\), as claimed.
\end{proof}

\begin{lemma}[lem:joint-orientation-frontier]\label{lem:joint-orientation-frontier}
For every \(N\geq1\) and \(g\geq0\), the joint orientation frontier is finite, is nondecreasing in each argument, and satisfies
\[
 \vartheta_{N,g}\leq\min\{2^g,\tau_N\}.
\]
Moreover,
\[
 \vartheta_{N,g}=2^g\qquad\text{whenever }N\geq2g+1.
\]
For five available vertices the exact finite recurrence values are
\[
 (\vartheta_{5,g})_{g=0}^{6}=(1,2,4,6,8,12,24),
 \qquad
 \vartheta_{5,g}=24\quad(g\geq6).
\]
\end{lemma}
\begin{proof}
For a fixed connected simple graph \(H\), choosing \(X_i\) from the factor indexed by \(\{i,j\}\) is the same as orienting that edge toward \(i\).  Consequently, by \(\mathrm{def:joint\mbox{-}orientation\mbox{-}frontier}\), \(c_H(r)\) is exactly the number of orientations of \(H\) with indegree vector \(r\).  For each \(b\leq N\) there are only \(2^{\binom b2}\) labeled simple graphs, and for a fixed graph only finitely many orientations and hence only finitely many realized indegree vectors occur.  Thus the maximum defining \(\vartheta_{N,g}\) is finite.  Increasing \(N\) or \(g\) only enlarges the defining feasible set, so \(\vartheta_{N,g}\) is nondecreasing in each argument.

Fix a feasible connected graph \(H\) on \(b\) vertices, fix a spanning tree \(S\) of \(H\), and put \(h=g(H)\).  Since \(S\) has \(b-1\) edges,
\[
 |E(H)\setminus E(S)|=|E(H)|-(b-1)=h.
\tag{1}
\]
Fix an indegree vector \(r\).  After the orientations of the \(h\) non-tree edges have been fixed, there is at most one compatible orientation of \(S\).  Indeed, a leaf of the current tree is incident to one remaining tree edge; after subtracting from its required indegree the already fixed contributions of all non-tree edges and all previously removed tree edges, that remaining edge has at most one possible direction.  Remove the leaf and iterate.  Induction on the number of tree vertices proves uniqueness, including the one-vertex base case.  There are \(2^h\) orientations of the non-tree edges, whence
\[
 c_H(r)\leq 2^h\leq 2^g.
\tag{2}
\]
By \(\mathrm{lem:tournament\mbox{-}score\mbox{-}fiber\mbox{-}ceiling}\), the same fiber also satisfies \(c_H(r)\leq\tau_N\) because \(b\leq N\).  Maximizing over the graphs and score vectors in \(\mathrm{def:joint\mbox{-}orientation\mbox{-}frontier}\) gives
\[
 \vartheta_{N,g}\leq\min\{2^g,\tau_N\}.
\tag{3}
\]

It remains to attain the first bound in the asserted range.  If \(g=0\), the one-vertex graph included in the definition has coefficient one, so \(\vartheta_{N,0}=1=2^0\) for every \(N\geq1\), by (2).  Let \(g\geq1\).  Form a bouquet of \(g\) triangles with common vertex \(z\): for each \(1\leq i\leq g\), introduce new vertices \(u_i,v_i\) and edges
\[
 \{z,u_i\},\qquad \{u_i,v_i\},\qquad \{v_i,z\}.
\tag{4}
\]
This is a connected simple graph on \(2g+1\) vertices with \(3g\) edges, and therefore
\[
 g(H)=3g-(2g+1)+1=g.
\tag{5}
\]
Orient every triangle as either one of its two directed cycles.  Each choice gives indegree \(g\) at \(z\) and indegree one at every \(u_i,v_i\).  The two choices are independent for the \(g\) triangles, so for
\[
 r_z=g,\qquad r_{u_i}=r_{v_i}=1\quad(1\leq i\leq g)
\tag{6}
\]
the score fiber contains \(2^g\) orientations.  If \(N\geq2g+1\), this graph is feasible; hence (2) and (6) give \(\vartheta_{N,g}=2^g\).

Finally, we certify the five-vertex values by a finite coefficient recurrence.  For a labeled graph \(H\) whose edges in lexicographic order are \(e_j=\{u_j,v_j\}\), define
\[
 A_{H,0}(r)=\mathbf 1\{r=0\},
 \qquad
 A_{H,j}(r)=A_{H,j-1}(r-\mathbf e_{u_j})
             +A_{H,j-1}(r-\mathbf e_{v_j}),
\tag{7}
\]
where a term having a negative coordinate is zero.  Expanding the factor \(X_{u_j}+X_{v_j}\) proves inductively that
\[
 A_{H,|E(H)|}(r)=c_H(r).
\tag{8}
\]
Let \(\mu(H)=\max_r c_H(r)\), and let \(N_{b,m}(s)\) count connected labeled graphs on \([b]\) having \(m\) edges and \(\mu(H)=s\).  Applying (7) to every subset of the \(\binom b2\leq10\) possible edges gives the following complete finite audit; within a row, \(N(s)\) abbreviates \(N_{b,m}(s)\):
\[
\begin{array}{c|c|l|c}
b&m&\{s:N_{b,m}(s)\ne0\}&\sum_sN_{b,m}(s)\\ \hline
1&0&N(1)=1&1\\
2&1&N(1)=1&1\\
3&2&N(1)=3&3\\
3&3&N(2)=1&1\\
4&3&N(1)=16&16\\
4&4&N(2)=15&15\\
4&5&N(3)=6&6\\
4&6&N(4)=1&1\\
5&4&N(1)=125&125\\
5&5&N(2)=222&222\\
5&6&N(3)=190,\ N(4)=15&205\\
5&7&N(4)=50,\ N(5)=60,\ N(6)=10&120\\
5&8&N(7)=15,\ N(8)=30&45\\
5&9&N(12)=10&10\\
5&10&N(24)=1&1
\end{array}
\tag{9}
\]
This enumeration involves no graph-isomorphism choice: scan all edge subsets in lexicographic order, retain a subset precisely when a search from vertex \(1\) reaches all \(b\) vertices, and compute every coefficient by (7).  Its graph-count column has the independent checksum
\[
 a_{1,0}=1,
 \qquad
 a_{b,m}=\binom{\binom b2}{m}
 -\sum_{q=1}^{b-1}\binom{b-1}{q-1}
   \sum_{\ell=0}^{m}a_{q,\ell}
   \binom{\binom{b-q}{2}}{m-\ell},
 \qquad a_{b,m}:=\sum_sN_{b,m}(s).
\tag{10}
\]
To derive (10), subtract from all \(m\)-edge graphs the disconnected graphs, partitioned uniquely by the size \(q\), internal edge count \(\ell\), and vertex set of the connected component containing vertex \(1\).  Substitution gives the last column of (9).

Matching lower certificates are as follows.  Put \(S_0=\{12,13,14,15\}\), and for \(0\leq h\leq6\) let \(H_h\) have edge set \(S_0\cup A_h\), with
\[
\begin{array}{c|l|c|c}
h&A_h&r&c_{H_h}(r)\\ \hline
0&\varnothing&(4,0,0,0,0)&1\\
1&\{23\}&(3,1,1,0,0)&2\\
2&\{25,34\}&(2,1,1,1,1)&4\\
3&\{23,24,25\}&(2,2,1,1,1)&6\\
4&\{23,24,25,34\}&(2,2,2,1,1)&8\\
5&\{23,24,25,34,35\}&(2,2,2,2,1)&12\\
6&\{23,24,25,34,35,45\}&(2,2,2,2,2)&24.
\end{array}
\tag{11}
\]
Every coefficient in (11) is obtained directly from (7).  Since \(H_h\) has five vertices and \(4+h\) edges, its cyclomatic number is \(h\).  The upper maxima in (9) and the lower certificates in (11), also including the smaller values of \(b\), therefore yield
\[
 (\vartheta_{5,g})_{g=0}^{6}=(1,2,4,6,8,12,24).
\tag{12}
\]
A connected simple graph on at most five vertices has cyclomatic number at most
\[
 \binom52-5+1=6,
\tag{13}
\]
with equality only possible at five vertices.  Hence increasing \(g\) beyond six adds no feasible graph, and (12) gives \(\vartheta_{5,g}=24\) for every \(g\geq6\).
\end{proof}

\begin{proposition}[prop:joint-orientation-maximality-threshold]\label{prop:joint-orientation-maximality-threshold}
For all integers \(N\geq1\) and \(g\geq0\),
\[
 \vartheta_{N,g}=2^g
 \quad\Longleftrightarrow\quad
 N\geq2g+1.
\]
Consequently \(\vartheta_{N,g}<2^g\) whenever \(N<2g+1\).
\end{proposition}
\begin{proof}
The forward construction and upper bound in \(\mathrm{lem:joint\mbox{-}orientation\mbox{-}frontier}\) already prove \(\vartheta_{N,g}=2^g\) when \(N\geq2g+1\).  We prove the converse.

Suppose \(\vartheta_{N,g}=2^g\).  By finiteness and attainment of the maximum in \(\mathrm{def:joint\mbox{-}orientation\mbox{-}frontier}\), there are a connected simple graph \(H\) on \(b\leq N\) vertices and a score vector \(r\) such that
\[
 c_H(r)=2^g.
\tag{14}
\]
Let \(h=g(H)\leq g\), choose a spanning tree \(S\), and put \(F=E(H)\setminus E(S)\).  The spanning-tree bound proved in \(\mathrm{lem:joint\mbox{-}orientation\mbox{-}frontier}\) gives
\[
 2^g=c_H(r)\leq2^h\leq2^g.
\tag{15}
\]
Thus \(h=g\), and equality in the first inequality says that every one of the \(2^g\) orientations of the cotree edges \(F\) extends, necessarily uniquely, to an orientation of \(S\) with total indegree vector \(r\).

Fix a reference direction on every edge.  If the reference direction of an edge \(a\) is from \(u\) to \(v\), let
\[
 \partial a=\mathbf e_v-\mathbf e_u,
\]
and encode its actual orientation by \(q_a=1\) when it agrees with the reference direction and by \(q_a=0\) otherwise.  Its contribution to the indegree vector is then \(\mathbf e_u+q_a\partial a\).  Let \(A_S\) and \(A_F\) be the matrices whose columns are these signed incidence vectors for tree and cotree edges, respectively, and let \(c\) be the sum of all reference-tail vectors \(\mathbf e_u\).  Writing \(x\in\{0,1\}^g\) for the cotree orientations and \(y(x)\in\{0,1\}^{b-1}\) for their unique tree extension, the fixed-indegree condition is
\[
 A_Sy(x)+A_Fx=r-c.
\tag{16}
\]
The incidence matrix of a tree has column rank \(b-1\): if \(A_Sz=0\), the row at a leaf forces the coefficient of its incident edge to be zero, and induction after deleting that leaf forces \(z=0\).  Hence (16) uniquely determines \(y(x)\) even over \(\mathbb R\), and subtraction of its value at \(x=0\) gives
\[
 y(x)=y(0)+\sum_{j=1}^g x_jv_j,
 \qquad A_Sv_j=-\partial f_j,
\tag{17}
\]
where \(F=\{f_1,\ldots,f_g\}\).

Let \(P_j\) be the unique path in \(S\) between the endpoints of \(f_j\).  Summing signed incidence columns along this path, with the signs chosen consistently from one endpoint to the other, gives \(-\partial f_j\).  Injectivity of \(A_S\) therefore shows that \(v_j\) is zero off \(P_j\) and is \(1\) or \(-1\) on every edge of \(P_j\).  In particular, for a fixed tree edge \(a\), its coordinate in (17) has the form
\[
 y_a(x)=y_a(0)+\sum_{j:a\in P_j}\epsilon_{aj}x_j,
 \qquad \epsilon_{aj}\in\{-1,1\}.
\tag{18}
\]
Every value in (18) belongs to \(\{0,1\}\), because every cotree orientation has a tree-orientation extension.  If \(y_a(0)=0\), evaluating at a unit vector forces every occurring \(\epsilon_{aj}\) to equal \(1\), while evaluating at the sum of two such unit vectors rules out two occurrences.  If \(y_a(0)=1\), the same argument forces every occurring coefficient to equal \(-1\) and again rules out two occurrences.  Thus each tree edge belongs to at most one of the paths \(P_1,\ldots,P_g\); these paths are pairwise edge-disjoint.

Because \(H\) is simple, the endpoints of a cotree edge \(f_j\) cannot already be joined by a single tree edge.  Therefore every \(P_j\) contains at least two tree edges.  Pairwise edge-disjointness now gives
\[
 b-1=|E(S)|
 \geq\sum_{j=1}^g|E(P_j)|
 \geq2g.
\tag{19}
\]
Hence \(N\geq b\geq2g+1\), proving the converse.  The displayed strict inequality follows from the already-proved bound \(\vartheta_{N,g}\leq2^g\), the equivalence, and integrality of \(\vartheta_{N,g}\).
\end{proof}

\begin{theorem}[thm:maximality-threshold]\label{thm:maximality-threshold}
For every admissible pair \(n\geq2\) and \(k\geq0\),
\[
 D_k(n)=2^{k+1}
 \quad\Longleftrightarrow\quad
 n\geq2k+4.
\]
More precisely, if \(n<2k+4\), then
\[
 D_k(n)\leq\vartheta_{n-1,k+1}<2^{k+1}.
\]
Thus the exact values in the remaining below-threshold strip are left open; in particular, no equality between \(D_k(n)\) and \(\vartheta_{n-1,k+1}\) is asserted there.
\end{theorem}
\begin{proof}
Fix an admissible pair \(n\geq2\) and \(k\geq0\), as required by \(\mathrm{ass:size}\) and \(\mathrm{ass:k\mbox{-}domain}\).  If
\[
 n\geq2k+4,
\tag{59}
\]
then \(\mathrm{thm:large\mbox{-}size\mbox{-}frontier}\) gives
\[
 D_k(n)=2^{k+1}.
\tag{60}
\]

Conversely, suppose \(n<2k+4\), and set \(N=n-1\) and \(g=k+1\).  Then \(N\geq1\), \(g\geq1\), and
\[
 N=n-1<2(k+1)+1=2g+1.
\tag{61}
\]
By \(\mathrm{prop:joint\mbox{-}orientation\mbox{-}maximality\mbox{-}threshold}\), (61) gives
\[
 \vartheta_{n-1,k+1}=\vartheta_{N,g}<2^g=2^{k+1}.
\tag{62}
\]
On the other hand, \(\mathrm{thm:orientation\mbox{-}upper\mbox{-}bound}\), on the graph class fixed by \(\mathrm{def:graph\mbox{-}class}\), yields
\[
 D_k(n)\leq\vartheta_{n-1,k+1}.
\tag{63}
\]
Combining (62)--(63) gives
\[
 D_k(n)\leq\vartheta_{n-1,k+1}<2^{k+1}.
\tag{64}
\]
Equations (60) and (64) prove
\[
 D_k(n)=2^{k+1}\quad\Longleftrightarrow\quad n\geq2k+4.
\tag{65}
\]
Below the threshold, (63) is only an upper bound; the argument neither determines \(D_k(n)\) nor identifies it with \(\vartheta_{n-1,k+1}\), as claimed.
\end{proof}

\section*{Technical internal limitation (diagnostic only)}

For \(k>0\), \(n\geq6\), and \(n<2k+4\), strict submaximality \(D_k(n)<2^{k+1}\) is proved, but the exact frontier remains between the constructive lower envelope and \(\vartheta_{n-1,k+1}\). The latter is a combinatorial upper bound, not a claim of SEM attainment: branches can be lost through forced covariance relations, affine dehomogenization, or projective infinity. The sound upper-bound proof controls this by counting only selected nonsingular conjugate branches and by perturbing them to a finite projective system before applying the finite-system B\'ezout citation. Independently, the quadratic core does not construct a reduced localized presentation of the embedded source prime. Exact finite output, exceptional sets, real lifting, and inference therefore retain their valid-chart and successful-output antecedents. The real-lifting proof uses quantifier elimination only at its attested scope and does not infer real uniqueness from complex degree or local rank.

\section{Honest scope}

The paper proves the full frontier through five vertices, the all-size lower envelope, the coupled orientation upper bound, and the equivalences \(\vartheta_{N,g}=2^g\) if and only if \(N\geq2g+1\) and \(D_k(n)=2^{k+1}\) if and only if \(n\geq2k+4\). Below the latter threshold, exact positive-excess values from six vertices onward remain open; \(\vartheta_{n-1,k+1}\) is not claimed to be attained. The cited multihomogeneous bound is used only on a finite perturbed projective system, never directly on a positive-dimensional residual variety. Unconditionally, positive excess yields only the randomized algebraic-versus-transcendental tag, rank or dominance certificate, and efficiently available ceiling \(d_e(G)\leq2^{\min\{|D|-n+2,\binom{n-2}{2}\}}\). Exact \(d_e(G)\), \(p_e\), certificate circuits, and the exceptional set are returned only from supplied valid source-prime charts. On that same conditional domain, the real-lifting and regular-inference results apply; quantifier elimination contributes only an equivalent finite sign formula, and exact evaluation is claimed only for algebraic specializations with finite representations. No discriminant-uniform inference, automatic branch selection, real uniqueness from complex degree, or all-singular-covariance guarantee is asserted.

\begin{thebibliography}{99}

  \bibitem{VanDerZanderEtAl2022TreeID} B. van der Zander, M. Wienöbst, M. Bläser, and M. Liśkiewicz. Identification in Tree-shaped Linear Structural Causal Models. Proceedings of the 25th International Conference on Artificial Intelligence and Statistics, Proceedings of Machine Learning Research 151, 2022, 6770--6792. arXiv:2203.01852.

  \bibitem{GuptaBlaser2024TreeSCM} A. Gupta and M. Bläser. Identification for Tree-Shaped Structural Causal Models in Polynomial Time. Proceedings of the AAAI Conference on Artificial Intelligence 38, 2024, 20404--20411. arXiv:2311.14058.

  \bibitem{BriefsBlaser2025FasterTree} Y. Briefs and M. Bläser. Faster Generic Identification in Tree-Shaped Structural Causal Models. Advances in Neural Information Processing Systems 38, 2025. DOI: 10.52202/085713-2989.

  \bibitem{FoygelDraismaDrton2012HTC} R. Foygel, J. Draisma, and M. Drton. Half-Trek Criterion for Generic Identifiability of Linear Structural Equation Models. The Annals of Statistics 40, 2012, 1682--1713. arXiv:1107.5552.

  \bibitem{DrtonFoygelSullivant2011Global} M. Drton, R. Foygel, and S. Sullivant. Global Identifiability of Linear Structural Equation Models. The Annals of Statistics 39, 2011, 865--886. arXiv:1003.1146.

  \bibitem{DrtonWeihs2016Ancestor} M. Drton and L. Weihs. Generic Identifiability of Linear Structural Equation Models by Ancestor Decomposition. Scandinavian Journal of Statistics 43, 2016, 1035--1045. arXiv:1504.02992.

  \bibitem{WeihsEtAl2018DeterminantalIV} L. Weihs, B. Robinson, E. Dufresne, J. Kenkel, K. McGee II, and N. Nguyen. Determinantal Generalizations of Instrumental Variables. Journal of Causal Inference 6, 2018. arXiv:1702.03884.

  \bibitem{DoerflerZanderBlaserLiskiewicz2024Complexity} J. Dörfler, B. van der Zander, M. Bläser, and M. Liśkiewicz. On the Complexity of Identification in Linear Structural Causal Models. Advances in Neural Information Processing Systems 37, 2024. arXiv:2407.12528.

  \bibitem{HolleringMisraSturma2026Symbolic} B. Hollering, P. Misra, and N. Sturma. Efficient Symbolic Computations for Identifying Causal Effects. arXiv:2604.20516, 2026.

  \bibitem{Corniquel2026GenericRUR} F. Corniquel. Solving Parametric Polynomial Systems Using Generic Rational Univariate Representation. arXiv:2602.06817, 2026.

  \bibitem{JeronimoSabia2007MultihomogeneousResultant} G. Jeronimo and J. Sabia. Computing Multihomogeneous Resultants Using Straight-Line Programs. Journal of Symbolic Computation 42, 2007, 218--235.

  \bibitem{Emiris1996SparseElimination} I. Z. Emiris. On the Complexity of Sparse Elimination. Journal of Complexity 12, 1996, 134--166.

  \bibitem{BouzidiLazardPougetRouillier2015RUR} Y. Bouzidi, S. Lazard, M. Pouget, and F. Rouillier. Separating Linear Forms and Rational Univariate Representations of Bivariate Systems. Journal of Symbolic Computation 68, 2015, 84--119. arXiv:1303.5042.

  \bibitem{ChernozhukovHongTamer2007SetInference} V. Chernozhukov, H. Hong, and E. Tamer. Estimation and Confidence Regions for Parameter Sets in Econometric Models. Econometrica 75, 2007, 1243--1284.

  \bibitem{Drton2009Singularities} M. Drton. Likelihood Ratio Tests and Singularities. The Annals of Statistics 37, 2009, 979--1012. arXiv:math/0703360.

  \bibitem{MareisSturmaDrton2026Inference} L. Mareis, N. Sturma, and M. Drton. Semiparametric Inference for Half-Trek Estimators in Linear Structural Equation Models. arXiv:2606.26931, 2026.

  \bibitem{BartzosEmirisSchicho2020Orientations} E. Bartzos, I. Z. Emiris, and J. Schicho. On the Multihomogeneous Bezout Bound on the Number of Embeddings of Minimally Rigid Graphs. arXiv:2005.14485, 2020.

  \bibitem{BasuPollackRoy1996QuantifierElimination} S. Basu, R. Pollack, and M.-F. Roy. On the Combinatorial and Algebraic Complexity of Quantifier Elimination. Journal of the ACM 43(6), 1996, 1002–1045. DOI: 10.1145/235809.235813.

  \bibitem{BasuPollackRoy2000Roadmaps} S. Basu, R. Pollack, and M.-F. Roy. Computing Roadmaps of Semi-algebraic Sets on a Variety. Journal of the American Mathematical Society 13(1), 2000, 55–82. DOI: 10.1090/S0894-0347-99-00311-2.

  \bibitem{GiustiLecerfSalvy2001GroebnerFree} M. Giusti, G. Lecerf, and B. Salvy. A Gröbner Free Alternative for Polynomial System Solving. Journal of Complexity 17(1), 2001, 154–211. DOI: 10.1006/jcom.2000.0571.

  \bibitem{Lecerf2003LiftingFibers} G. Lecerf. Computing the Equidimensional Decomposition of an Algebraic Closed Set by Means of Lifting Fibers. Journal of Complexity 19(4), 2003, 564–596. DOI: 10.1016/S0885-064X(03)00031-1.

\end{thebibliography}

\end{document}
